\documentclass[uplatex, dvipdfmx]{jsreport}

% ============================================================
%  preamble.tex — 日本語数学ノート テンプレート共通設定
%
%  定理環境のカウンタ規約は tools/site が再現している：
%    - 全定理環境が definition のカウンタを共有（use counter from）
%    - カウンタは節ごとにリセット（number within=section）
%  この規約を変えるときは tools/site/lib/tex2md.mjs の
%  buildNumberMap() も合わせること（PDF とサイトで番号がずれる）。
% ============================================================

% --- 基本パッケージ ---
\usepackage{amsmath,amssymb,amsthm,mathtools}
\usepackage[dvipdfmx,hidelinks]{hyperref}
\usepackage{pxjahyper}

% --- 図・表 ---
\usepackage{tikz}
\usepackage{float}
\usetikzlibrary{decorations.pathreplacing,calc,arrows.meta,patterns,positioning}
\usepackage{pgfplots}
\pgfplotsset{compat=1.18}
\usepackage{graphicx}
\graphicspath{{fig/}}

% --- 定理環境 (tcolorbox) ---
% 第 2 引数はラベル接頭辞。tools/site はこの接頭辞で \ref を解決するので、
% def / prop / thm / clm / lem / cor / rem / ex は変えないこと。
\usepackage{tcolorbox}
\tcbuselibrary{theorems,breakable}

% Definition: 青系
\newtcbtheorem[number within=section]{problem}{問題}{
  colback=teal!4, colframe=teal!55!black, fonttitle=\bfseries, breakable
}{prob}
% Definition: 青系
\newtcbtheorem[number within=section]{definition}{Def.}{
  colback=blue!5, colframe=blue!40!black, fonttitle=\bfseries, breakable
}{def}
% Proposition: オレンジ系
\newtcbtheorem[use counter from=definition]{proposition}{Prop.}{
  colback=orange!5, colframe=orange!50!black, fonttitle=\bfseries, breakable
}{prop}
% Theorem: 赤系（章・節の目玉となる重要結果）
\newtcbtheorem[use counter from=definition]{theorem}{Thm.}{
  colback=red!5, colframe=red!40!black, fonttitle=\bfseries, breakable
}{thm}
% Claim: 緑系（Thm ほど深くはないが，証明を要する主張）
\newtcbtheorem[use counter from=definition]{claim}{Claim}{
  colback=green!5, colframe=green!45!black, fonttitle=\bfseries, breakable
}{clm}
% Lemma: シアン系（命題の証明に用いる補助結果）
\newtcbtheorem[use counter from=definition]{lemma}{Lem.}{
  colback=cyan!5, colframe=cyan!50!black, fonttitle=\bfseries, breakable
}{lem}
% Corollary: すみれ系（定理から直ちに従う結果）
\newtcbtheorem[use counter from=definition]{corollary}{Cor.}{
  colback=violet!5, colframe=violet!40!black, fonttitle=\bfseries, breakable
}{cor}
% Remark: 黄系
\newtcbtheorem[use counter from=definition]{remark}{Rem.}{
  colback=yellow!5, colframe=yellow!40!black, fonttitle=\bfseries, breakable
}{rem}
% Example: 灰色系
\newtcbtheorem[use counter from=definition]{example}{Ex.}{
  colback=gray!5, colframe=gray!50!black, fonttitle=\bfseries, breakable
}{ex}
% Memo: 薄い黒系（番号なし，発展的補足）
\newtcolorbox{memo*}{
  colback=black!7, colframe=black!40,
  fonttitle=\bfseries, title={Memo},
  breakable
}
% 証明の直後が enumerate でも、最初の項目だけ見出しの横へずれないようにする。
% 通常の文章で始まる証明は従来どおり見出しと同じ行から始まる。
\makeatletter
\renewenvironment{proof}[1][\proofname]{\par
  \pushQED{\qed}%
  \normalfont \topsep6\p@\@plus6\p@\relax
  \trivlist
  \item[\hskip\labelsep\itshape #1\@addpunct{.}]\ignorespaces\leavevmode
}{%
  \popQED\endtrivlist\@endpefalse
}
\makeatother

% 証明環境の装飾
\tcolorboxenvironment{proof}{
  colback=white, colframe=black!20,
  leftrule=2pt, rightrule=0pt, toprule=0pt, bottomrule=0pt,
  breakable
}

% --- アルゴリズム環境 ---
\newcounter{algcounter}[section]
\newtcolorbox{algorithm}[1]{
  colback=white, colframe=black!70,
  fonttitle=\bfseries,
  title={Algorithm~\refstepcounter{algcounter}\thealgcounter\label{#1}},
  boxrule=0.8pt
}

% --- サイト向けの指示マクロ（PDF では何も出力しない）---
% \demohint{名前}: site.config.mjs の demos に登録した図をサイト側で差し込む
\newcommand{\demohint}[1]{}
% \blockmeta{...}: ブロックの機械可読メタデータ。PDF にもサイトにも出さない
\newcommand{\blockmeta}[1]{}

% ============================================================
%  数式マクロ
%  ここに書いたものは tools/site が自動抽出して MathJax にも渡す。
%  サイト側の設定に書き写す必要はない。
%  MathJax が解釈できない綴りだけ site.config.mjs の macroOverrides で
%  差し替えること（例: stmaryrd の \llbracket）。
% ============================================================

\newcommand{\R}{\mathbb{R}}
\newcommand{\N}{\mathbb{N}}
\newcommand{\Z}{\mathbb{Z}}
\newcommand{\Q}{\mathbb{Q}}
\newcommand{\E}{\mathbb{E}}

\newcommand{\X}{\mathcal{X}}
\newcommand{\Y}{\mathcal{Y}}
\newcommand{\Bb}{\mathcal{B}}
\newcommand{\Prob}{\mathcal{P}}

\DeclareMathOperator*{\argmin}{arg\,min}
\DeclareMathOperator*{\argmax}{arg\,max}
\DeclareMathOperator{\diag}{diag}
\DeclareMathOperator{\tr}{tr}
\DeclareMathOperator{\supp}{supp}
\DeclareMathOperator{\Id}{Id}

\newcommand{\abs}[1]{\lvert #1 \rvert}
\newcommand{\norm}[1]{\lVert #1 \rVert}
\newcommand{\inner}[2]{\langle #1,\, #2 \rangle}
\newcommand{\defeq}{\overset{\mathrm{def}}{=}}
\renewcommand{\d}{\mathrm{d}}

% --- 参考文献 ---
\usepackage[numbers,sort&compress]{natbib}

% ============================================================
%  演習用・解答編で共通の紙面設計
%
%  見出しは本文に出さず、柱（左上）と目次にだけ回す。
%  本文に残るのは問題文と、その終わりを示す細い罫線だけにする。
%  原稿 generated/{workbook,answers}/*.tex は
%  tools/split-pdf-sources.mjs が生成する。
%
%  この後で \wbend を定義すると紙面が決まる。
%    演習用 … 罫線から下をそのまま解答欄にする（何も足さない）
%    解答編 … 罫線の下に見通しと解答を続ける
% ============================================================

\usepackage[top=18mm, bottom=14mm, left=20mm, right=20mm, headsep=5mm]{geometry}
\usepackage{fancyhdr}

% 目次は「年度・区分」までにとどめて1ページに収め、
% 問題名は PDF のしおりからだけ辿れるようにする。
\hypersetup{bookmarksdepth=2}

\makeatletter

\newcommand{\wb@year}{}
\newcommand{\wb@sec}{}
\newcommand{\wb@prob}{}
\newcommand{\wb@sep}{\,\textbar\,}

\pagestyle{fancy}
\fancyhf{}
\fancyhead[L]{\normalfont\footnotesize\textcolor{black!60}{%
  \wb@year\wb@sep\wb@sec\wb@sep\textbf{\wb@prob}}}
\fancyhead[R]{\normalfont\footnotesize\textcolor{black!60}{\thepage}}
\renewcommand{\headrulewidth}{0.4pt}
\renewcommand{\footrulewidth}{0pt}
\fancypagestyle{plain}{%
  \fancyhf{}%
  \fancyhead[R]{\normalfont\footnotesize\textcolor{black!60}{\thepage}}%
  \renewcommand{\headrulewidth}{0pt}%
}

% 章・節は柱の材料と目次の見出しにするだけで、本文には何も出さない。
\newcommand{\wbchapter}[1]{%
  \clearpage
  \renewcommand{\wb@year}{#1}%
  \phantomsection\addcontentsline{toc}{chapter}{#1}%
}
\newcommand{\wbsection}[1]{%
  \clearpage
  \renewcommand{\wb@sec}{#1}%
  \phantomsection\addcontentsline{toc}{section}{#1}%
}
% 問題は必ず新しいページから始める。題名は柱と目次へ送る。
\newcommand{\wbproblem}[1]{%
  \clearpage
  \renewcommand{\wb@prob}{#1}%
  \phantomsection\addcontentsline{toc}{subsection}{#1}%
  \par\noindent\ignorespaces
}
% 問題文の終わりを示す細い罫線。
\newcommand{\wbrule}{%
  \par\medskip
  \noindent\textcolor{black!25}{\rule{\linewidth}{0.4pt}}%
  \par
}

\makeatother


% 罫線の下に解答を続ける。
\newcommand{\wbend}{\wbrule}

\title{熊本大学 数学系大学院 入試対策\\{\Large 専門基礎科目｜演習用 解答編}}
\author{}
\date{\today}

\begin{document}
\maketitle

\begin{memo*}
演習用（\texttt{kumadai-workbook.pdf}）の解答編である。
過去問34題と予想問題6題を、演習用とまったく同じ順序・同じ柱で並べ、
各問題の罫線の下に解答を置いた。

答案に必要な式と論証を中心に、簡潔に記した。
解答で使う用語・定理を確認したいときは別冊「用語・定義・定理 補足資料」を、
前提から積み上げて学び直したいときは別冊「専門基礎科目のための教科書」を参照されたい。
\end{memo*}

\tableofcontents

\chapter{2026年度}
\label{ch:2026}

\section{第1期・専門科目（4題から2題選択）}

\begin{problem}{第1問｜群論}{2026-1-1}
$S_6$ の元 $\sigma=(1\ 2)(3\ 4)(5\ 6)$ と，その中心化群
\[
Z_{S_6}(\sigma)=\{\tau\in S_6\mid \tau\sigma\tau^{-1}=\sigma\}
\]
を考える。積は $\sigma\tau(i)=\sigma(\tau(i))$ で定める。
\begin{enumerate}
  \item[(1)] $\rho\in S_6$ と巡回置換 $(k_1\ k_2\ \cdots\ k_r)$ に対し
  $\rho(k_1\ k_2\ \cdots\ k_r)\rho^{-1}=(\rho(k_1)\ \rho(k_2)\ \cdots\ \rho(k_r))$
  を示せ。
  \item[(2)] $S_6$ における $\sigma$ の共役類の元の個数を求めよ。
  \item[(3)] $Z_{S_6}(\sigma)$ の位数を求めよ。
  \item[(4)] $Z_{S_6}(\sigma)$ の Sylow $3$ 部分群を一つ求めよ。
  \item[(5)] $H=\langle(1\ 2),(3\ 4),(5\ 6)\rangle$ が $Z_{S_6}(\sigma)$ の正規部分群であることを示せ。
\end{enumerate}
\end{problem}
\begin{proof}[解答]
\begin{enumerate}
  \item[(1)] $k_{r+1}=k_1$ とする。$x=\rho(k_i)$ なら左辺は $\rho(k_{i+1})$ に写し，右辺も同じ。どの $\rho(k_i)$ でもない $x$ は両辺が固定する。よって二つの置換は等しい。
  \item[(2)] 共役類は型 $(2,2,2)$ の置換全体である。6文字を並べ，3組の順序と各組内の順序を除けば
  \[
  \frac{6!}{3!\,2^3}=15.
  \]
  \item[(3)] 軌道・安定化群定理から
  $\abs{Z_{S_6}(\sigma)}=6!/15=48$。
  \item[(4)] $g=(1\ 3\ 5)(2\ 4\ 6)$ は三つの互換を巡回させるので $g\sigma g^{-1}=\sigma$。したがって $\langle g\rangle$ は位数 $3$ の Sylow $3$ 部分群である。
  \item[(5)] $H\cong(C_2)^3$ であり，各生成元は $\sigma$ と可換するから $H\subset Z_{S_6}(\sigma)$。$\rho\in Z_{S_6}(\sigma)$ は三つの組 $\{1,2\},\{3,4\},\{5,6\}$ を置換する。ゆえに $\rho(k\ k+1)\rho^{-1}$ は再び $H$ の生成元であり，$\rho H\rho^{-1}=H$。
\end{enumerate}
\end{proof}

\begin{problem}{第2問｜基本群と被覆空間}{2026-1-2}
$S^1\subset\R^2$ を単位円，$x_0=(1,0)$ とする。
\begin{enumerate}
  \item[(1)] $\pi_1(S^1\times S^1,(x_0,x_0))$ を求めよ。ただし $\pi_1(S^1,x_0)\cong\Z$ は用いてよい。
  \item[(2)] $S^1\times S^1$ の普遍被覆空間が $\R^2$ と同型であることを示せ。
  \item[(3)] $S^1\times S^1$ の被覆空間で，基本群が $\Z$ と同型なものを構成せよ。
\end{enumerate}
\end{problem}
\begin{proof}[解答]
\begin{enumerate}
  \item[(1)] 閉道を成分ごとに送る $[(f,g)]\mapsto([f],[g])$ は群同型だから
  \[
  \pi_1(S^1\times S^1)\cong\pi_1(S^1)\times\pi_1(S^1)\cong\Z^2.
  \]
  \item[(2)] $p(t)=(\cos2\pi t,\sin2\pi t)$ とおけば $p\times p:\R^2\to S^1\times S^1$ は被覆写像。$\R^2$ は単連結なので，これが普遍被覆である。
  \item[(3)] $q:S^1\times\R\to S^1\times S^1$, $q(z,t)=(z,p(t))$ は被覆写像で，$\pi_1(S^1\times\R)\cong\Z$。
\end{enumerate}
\end{proof}

\begin{problem}{第3問｜複素積分}{2026-1-3}
\[
f(z)=\frac{\sin z}{e^z-1-z},\qquad C=\{z\mid |z|=1\}
\]
とし，$C$ は反時計回りとする。
\begin{enumerate}
  \item[(1)] $g(z)=e^z-1-z$ の Maclaurin 展開と，零点 $z=0$ の位数を求めよ。
  \item[(2)] $|z|=1$ のとき
  \[
  \left|e^z-1-z-\frac{z^2}{2}\right|\le e-\frac52
  \]
  を示し，$g$ が $C$ の内部に $0$ 以外の零点を持たないことを示せ。$2.7<e<2.8$ は用いてよい。
  \item[(3)] $\displaystyle\int_C f(z)\,\d z$ を求めよ。
\end{enumerate}
\end{problem}
\begin{proof}[解答]
\begin{enumerate}
  \item[(1)] 展開すると
  \[
  g(z)=\sum_{n=2}^{\infty}\frac{z^n}{n!}
      =z^2\sum_{k=0}^{\infty}\frac{z^k}{(k+2)!}.
  \]
  後ろの級数は $z=0$ で $1/2$ だから，位数は $2$。
  \item[(2)] 三角不等式から左辺は $\sum_{n=3}^{\infty}1/n!=e-5/2<1/2=|z^2/2|$。Rouché の定理により $g=z^2/2+(g-z^2/2)$ と $z^2/2$ は円内の零点数が同じく2。$0$ が既に2位の零点なので，他にはない。
  \item[(3)] 円内の極は $0$ のみで1位。よって
  \[
  \operatorname{Res}(f;0)=\lim_{z\to0}\frac{z\sin z}{e^z-1-z}=2,
  \qquad \int_Cf(z)\,\d z=4\pi i.
  \]
\end{enumerate}
\end{proof}

\begin{problem}{第4問｜ルベーグ積分}{2026-1-4}
$G:\R\to\R$ は有界連続関数，$M=\sup_{x\in\R}|G(x)|$ とする。$f_n,f:\R^d\to\R$ はルベーグ可測で，$f_n\to f$ がほとんど至る所で成り立つ。
\begin{enumerate}
  \item[(1)] 任意の開集合 $I\subset\R$ に対し $(G\circ f)^{-1}(I)$ がルベーグ可測であることを示せ。
  \item[(2)] $G\circ f_n\to G\circ f$ がほとんど至る所で成り立つことを示せ。
  \item[(3)] 可測集合 $E\subset\R^d$ が $|E|<\infty$ を満たすとき
  \[
  \lim_{n\to\infty}\int_E(G\circ f_n)(x)\,\d x
  =\int_E(G\circ f)(x)\,\d x
  \]
  を示せ。
\end{enumerate}
\end{problem}
\begin{proof}[解答]
\begin{enumerate}
  \item[(1)] 連続性により $G^{-1}(I)$ は開集合。したがって $(G\circ f)^{-1}(I)=f^{-1}(G^{-1}(I))$ は可測である。
  \item[(2)] $f_n(x)\to f(x)$ となる各 $x$ で，$G$ の連続性から $G(f_n(x))\to G(f(x))$。
  \item[(3)] $|G(f_n)|\chi_E\le M\chi_E$ かつ $\int M\chi_E=M|E|<\infty$。優収束定理を適用すればよい。
\end{enumerate}
\end{proof}

\section{第2期・専門基礎科目（3題必修）}

\begin{problem}{第1問｜エルミート行列}{2026-2-1}
\[
A=\begin{pmatrix}
2&-i&1&i\\ i&2&i&-1\\ 1&-i&2&i\\ -i&-1&-i&2
\end{pmatrix}.
\]
\begin{enumerate}
  \item[(1)] $\det A$ を求めよ。
  \item[(2)] 固有値と各固有空間を求めよ。
  \item[(3)] $A$ をユニタリ行列で対角化せよ。
  \item[(4)] エルミート行列の固有値がすべて実数であることを示せ。
\end{enumerate}
\end{problem}
\begin{proof}[解答]
$v=(i,-1,i,1)^{\mathsf T}$ とおくと $A=I_4+vv^*$ である。
\begin{enumerate}
  \item[(1)] $v$ 方向の固有値は $1+\norm{v}^2=5$，その直交補空間では $1$。したがって $\det A=5$。
  \item[(2)] 固有値は $1$（重複度3）と $5$（重複度1）で，
  \[
  W_5=\langle v\rangle,\qquad W_1=v^\perp
  =\{x\in\mathbb C^4\mid -ix_1-x_2-ix_3+x_4=0\}.
  \]
  \item[(3)] 次の列は正規直交基底である：
  \[
  u_1=\frac1{\sqrt2}(1,-i,0,0)^{\mathsf T},\quad
  u_2=\frac1{\sqrt2}(0,0,1,i)^{\mathsf T},
  \]
  \[
  u_3=\frac12(-i,1,i,1)^{\mathsf T},\quad
  u_4=\frac12(i,-1,i,1)^{\mathsf T}.
  \]
  $P=(u_1\ u_2\ u_3\ u_4)$ とすれば $P^*P=I$ かつ $P^*AP=\diag(1,1,1,5)$。
  \item[(4)] $Xw=\lambda w$，$w\ne0$ とする。$X=X^*$ なので
  $\lambda\inner{w}{w}=\inner{Xw}{w}=\inner{w}{Xw}=\overline\lambda\inner{w}{w}$。
  $\inner{w}{w}>0$ より $\lambda=\overline\lambda\in\R$。
\end{enumerate}
\end{proof}

\begin{problem}{第2問｜重積分}{2026-2-2}
\[
D=\{(x,y)\in\R^2\mid1\le x^2+y^2\le4,\ 0\le x\le1,\ y\ge0\}.
\]
\begin{enumerate}
  \item[(1)] $D$ を図示せよ。
  \item[(2)] $-\pi/2<\theta<\pi/2$ で $\displaystyle\log\left(\tan\theta+\frac1{\cos\theta}\right)$ の導関数を求めよ。
  \item[(3)] $\displaystyle\iint_D\frac{x}{\sqrt{x^2+y^2}}\,\d x\d y$ を求めよ。
\end{enumerate}
\end{problem}
\begin{proof}[解答]
\begin{enumerate}
  \item[(1)] 上半平面にある半径1と2の半円の間を，縦線 $x=0,1$ で切った領域である。
  \item[(2)] 直接微分して整理すると $1/\cos\theta$。
  \item[(3)] 極座標で $x/r=\cos\theta$，Jacobian は $r$。$0\le\theta\le\pi/3$ では $1\le r\le\sec\theta$，$\pi/3\le\theta\le\pi/2$ では $1\le r\le2$ だから
  \[
  \int_0^{\pi/3}\int_1^{\sec\theta}r\cos\theta\,\d r\d\theta
  +\int_{\pi/3}^{\pi/2}\int_1^2r\cos\theta\,\d r\d\theta
  =\frac32-\sqrt3+\frac12\log(2+\sqrt3).
  \]
\end{enumerate}
\end{proof}

\begin{problem}{第3問｜集合までの距離}{2026-2-3}
距離空間 $(X,d)$ と空でない部分集合 $A\subset X$ に対し
$d(x,A)=\inf\{d(x,a)\mid a\in A\}$ とおく。
\begin{enumerate}
  \item[(1)] $|d(x,A)-d(y,A)|\le d(x,y)$ を示せ。
  \item[(2)] $f(x)=d(x,A)$ が連続であることを示せ。
  \item[(3)] $x$ が $A$ の触点であることと $d(x,A)=0$ が同値であることを示せ。
\end{enumerate}
\end{problem}
\begin{proof}[解答]
\begin{enumerate}
  \item[(1)] 任意の $a\in A$ に対し $d(x,a)\le d(x,y)+d(y,a)$。下限を取って $d(x,A)\le d(x,y)+d(y,A)$。$x,y$ を交換した式と合わせればよい。
  \item[(2)] (1) は $f$ が $1$-Lipschitz であることを表すので連続。
  \item[(3)] $d(x,A)=0$ は「任意の $\varepsilon>0$ に対し $d(x,a)<\varepsilon$ となる $a\in A$ がある」と同値。これは任意の開球 $B(x,\varepsilon)$ が $A$ と交わる，すなわち $x\in\overline A$ と同値である。
\end{enumerate}
\end{proof}

\chapter{2022年度}
\label{ch:2022}

\section{第1期・専門基礎科目（3題必修）}

\begin{problem}{第1問｜実行列}{2022-1-1}
$n$ 次の単位行列を $E_n$ で表す。このとき，以下の問いに答えよ。
\begin{enumerate}
  \item[(1)] $A^2=-E_n$ となる $n$ 次の実正方行列 $A$ が存在するならば，$n$ は偶数であることを示せ。また，$A^2=-E_2$ となる2次の実正方行列 $A$ の例を挙げよ。
  \item[(2)] ${}^{\mathsf T}AA=-E_n$ となる $n$ 次の実正方行列 $A$ は存在しないことを示せ。
  \item[(3)] $A$ は零行列ではない $n$ 次の実正方行列で ${}^{\mathsf T}A=-A$ を満たすとする。$P^{-1}AP$ が対角行列となるような $n$ 次の実正則行列 $P$ は存在しないことを示せ。
\end{enumerate}
\end{problem}
\begin{memo*}
\textbf{解答の見通し。}
行列の等式は成分が多く、そのままでは比較しにくい。そこで、行列式や
$x^{\mathsf T}(\cdot)x$ を使って、行列の等式を一つの実数の等式へ変える。
(1) では $|BC|=|B|\,|C|$ により $A^2=-E_n$ を
$|A|^2=(-1)^n$ へ変え、左右の符号を比べる。
(2) では $x^{\mathsf T}A^{\mathsf T}Ax=\norm{Ax}^2\ge0$ なのに、
右辺は $-\norm{x}^2<0$ になることを使う。
(3) の $A^{\mathsf T}=-A$ は交代行列の条件であり、まず全ての実固有値が0だと示す。
もし実数上で対角化できれば、対角成分は全て0となって $A=O$ になり、仮定と矛盾する。
\end{memo*}
\begin{proof}[解答]
\begin{enumerate}
  \item[(1)] 行列式を取ると
  \[
  (\det A)^2=\det(-E_n)=(-1)^n.
  \]
  左辺は正だから $n$ は偶数である。また
  \[
  A=\begin{pmatrix}0&-1\\1&0\end{pmatrix}
  \]
  とすれば $A^2=-E_2$ である。

  \item[(2)] $0\ne x\in\R^n$ に対して
  \[
  x^{\mathsf T}A^{\mathsf T}Ax=\|Ax\|^2\ge0,\qquad
  x^{\mathsf T}(-E_n)x=-\|x\|^2<0.
  \]
  従って $A^{\mathsf T}A=-E_n$ は不可能である。

  \item[(3)] $Av=\lambda v$ $(0\ne v\in\R^n,\ \lambda\in\R)$ なら
  \[
  \lambda\|v\|^2=v^{\mathsf T}Av
  =(v^{\mathsf T}Av)^{\mathsf T}
  =v^{\mathsf T}A^{\mathsf T}v=-v^{\mathsf T}Av,
  \]
  よって $\lambda=0$ である。実数上で対角化できれば対角行列は零行列となり、
  $A=O$ となるので仮定に反する。
\end{enumerate}
\end{proof}

\begin{problem}{第2問｜Cauchy の関数方程式}{2022-1-2}
$\R$ 上で定義された実数値関数 $f(x)$ が，任意の実数 $x,y$ に対して
\[
f(x+y)=f(x)+f(y)
\]
を満たすとする。このとき，以下の問いに答えよ。
\begin{enumerate}
  \item[(1)] $f(0)=0$, $f(-x)=-f(x)$ が成り立つことを示せ。
  \item[(2)] 任意の有理数 $r\in\Q$ に対して，$f(rx)=rf(x)$ が成り立つことを示せ。
  \item[(3)] $f(x)$ が $\R$ 上で連続とすると，$f(x)$ はある定数 $c\in\R$ を用いて，$f(x)=cx$ と書けることを示せ。
\end{enumerate}
\end{problem}
\begin{memo*}
\textbf{解答の見通し。}
条件 $f(x+y)=f(x)+f(y)$ は、入力の足し算が出力の足し算へ移るという意味である。
まず式へ $x=y=0$ を入れて $f(0)$ を決め、次に $y=-x$ を入れて負号の扱いを決める。
同じ $x$ を繰り返し足せば整数倍、$x=n(x/n)$ と書けば $1/n$ 倍、両者を組み合わせれば
有理数倍に対する式が得られる。最後に、任意の実数 $x$ へ近づく有理数列を取り、
連続性によって有理数上の公式 $f(r)=f(1)r$ を実数全体へ移す。
\end{memo*}
\begin{proof}[解答]
\begin{enumerate}
  \item[(1)]
  \[
  f(0)=f(0)+f(0)\quad\Longrightarrow\quad f(0)=0,
  \]
  \[
  0=f(x-x)=f(x)+f(-x)\quad\Longrightarrow\quad f(-x)=-f(x).
  \]

  \item[(2)] 加法性と (1) より $f(mx)=mf(x)$ $(m\in\Z)$。また
  \[
  f(x)=f\left(n\frac{x}{n}\right)=nf\left(\frac{x}{n}\right)
  \]
  だから、$r=m/n$ $(m\in\Z,\ n\in\N)$ に対して
  \[
  f(rx)=f\left(m\frac{x}{n}\right)=\frac{m}{n}f(x)=rf(x).
  \]

  \item[(3)] $c=f(1)$ とおくと (2) より $f(r)=cr$ $(r\in\Q)$。
  $r_n\in\Q$, $r_n\to x$ を取れば、連続性から
  \[
  f(x)=\lim_{n\to\infty}f(r_n)
  =c\lim_{n\to\infty}r_n=cx.
  \]
\end{enumerate}
\end{proof}

\begin{problem}{第3問｜連続写像の反例}{2022-1-3}
位相空間 $X$ から位相空間 $Y$ への連続写像 $f:X\to Y$ に関する次の各命題はどれも成り立たない。反例を挙げよ。
\begin{enumerate}
  \item[(1)] $X$ の開集合 $U$ に対して像 $f(U)$ は $Y$ の開集合である。
  \item[(2)] $Y$ のコンパクト部分集合 $K$ に対して逆像 $f^{-1}(K)$ は $X$ のコンパクト部分集合である。
  \item[(3)] $f$ が全射のとき，$Y$ の稠密な部分集合 $V$ に対して逆像 $f^{-1}(V)$ は $X$ の稠密な部分集合である。
\end{enumerate}
\end{problem}
\begin{memo*}
\textbf{解答の見通し。}
反例とは、仮定を全て満たすのに結論だけが失敗する具体例である。
連続性が直接保証するのは「開集合の\emph{逆像}が開」であり、開集合の像や
コンパクト集合の逆像ではない。(1), (2) は全ての点を0へ送る定数写像で、この違いを示す。
(3) では全射性も必要なので、実数直線はそのまま恒等的に写し、別に加えた孤立点 $p$ だけを0へ送る。
孤立点とは一元集合 $\{p\}$ 自身が開になる点であり、これが逆像の稠密性を壊す。
\end{memo*}
\begin{proof}[解答]
\begin{enumerate}
  \item[(1)] 定数写像 $f:\R\to\R$, $f(x)=0$ は連続である。
  開集合 $U=\R$ に対して $f(U)=\{0\}$ は開でない。

  \item[(2)] 同じ写像で $K=\{0\}$ はコンパクトだが
  $f^{-1}(K)=\R$ はコンパクトでない。

  \item[(3)] $X=\R\sqcup\{p\}$ に直和位相を入れ、$Y=\R$ とする。
  \[
  f(x)=x\ (x\in\R),\qquad f(p)=0
  \]
  とおけば $f$ は連続かつ全射である。$V=\R\setminus\{0\}$ は $Y$ で稠密だが
  \[
  f^{-1}(V)=\R\setminus\{0\}
  \]
  は $X$ の非空開集合 $\{p\}$ と交わらないので稠密でない。
\end{enumerate}
\end{proof}

\section{第1期・専門科目（4題から2題選択）}

\begin{problem}{第1問｜有限群}{2022-s-1}
以下の問いに答えよ。
\begin{enumerate}
  \item[(1)] $\langle x\rangle$ を位数5の群とする。$\langle x\rangle$ の自己同型群（$\langle x\rangle$ から $\langle x\rangle$ への全単射準同型写像全体の集合）の位数を求めよ。答だけでなく，どのように求めるかも説明せよ。
  \item[(2)] $\{1,2,3,4,5\}$ に作用する5次対称群を $G$ とする。$P=\langle(1\ 2\ 3\ 4\ 5)\rangle$ とおく。$P$ の $G$ における正規化群 $N_G(P)=\{g\in G\mid gPg^{-1}=P\}$ のシロー2部分群を1つ求めよ。
  \item[(3)] 5次対称群 $G$ のシロー5部分群全体の集合を $\Omega$ とし，$\Omega$ の置換全体の集合を $\operatorname{Sym}(\Omega)$ とおく。$g\in G$ に対して
  \[
  \varphi(g):\Omega\longrightarrow\Omega,\qquad
  P\longmapsto\varphi(g)(P)=gPg^{-1}
  \]
  とおくと，$\varphi(g)\in\operatorname{Sym}(\Omega)$ である。これにより写像 $\varphi:G\to\operatorname{Sym}(\Omega)$ が定義できる。このとき，$\varphi$ は準同型写像であることを示せ。
  \item[(4)] 6次対称群には指数6の部分群で互いに共役ではないものが存在することを示せ。
\end{enumerate}
\end{problem}
\begin{memo*}
\textbf{解答の見通し。}
位数5の部分群を数えると、$S_5$ が6個の Sylow $5$ 部分群へ作用することが分かる。
この共役作用から通常とは異なる $S_5\hookrightarrow S_6$ を作り、
最後に「6点へ推移的に作用するか、1点を固定するか」という共役で変わらない性質で
二つの指数6部分群を区別する。
\end{memo*}
\begin{proof}[解答]
\begin{enumerate}
  \item[(1)] 位数5の群 $\langle x\rangle$ は巡回群である。自己同型 $\varphi$ は生成元 $x$ の像
  $\varphi(x)$ によって一意に決まり、また自己同型は生成元を生成元へ写す。
  $\langle x\rangle$ の生成元は
  \[
  x,x^2,x^3,x^4
  \]
  の4個である。逆に $k=1,2,3,4$ に対して $x\mapsto x^k$ は自己同型を定めるので、
  自己同型は生成元の像で尽くされ
  \[
  |\operatorname{Aut}(\langle x\rangle)|=4.
  \]

  \item[(2)] $S_5$ の5-cycle の個数は、最初の文字を1に固定して残りを並べることで
  \[
  (5-1)!=24
  \]
  個である。一つの位数5の部分群には単位元でない5-cycle が4個あり、異なる位数5の部分群は
  単位元以外では交わらない。従って Sylow $5$ 部分群の個数は
  \[
  |\Omega|=\frac{24}{4}=6.
  \]
  共役作用における $P$ の安定化群は $N_G(P)$ であり、Sylow の定理により作用は推移的だから、
  軌道・安定化群定理から
  \[
  |N_G(P)|=\frac{|S_5|}{|\Omega|}
  =\frac{120}{6}=20=2^2\cdot5.
  \]

  ここで
  \[
  a=(1\ 2\ 3\ 4\ 5),
  \qquad b=(2\ 3\ 5\ 4)
  \]
  とおく。$b$ は4-cycle なので位数は4である。また巡回置換の共役公式から
  \[
  bab^{-1}
  =(b(1)\ b(2)\ b(3)\ b(4)\ b(5))
  =(1\ 3\ 5\ 2\ 4)=a^2.
  \]
  従って $bPb^{-1}=P$、すなわち $b\in N_G(P)$ である。
  よって $\langle b\rangle$ は $N_G(P)$ の位数4の部分群であり、
  $|N_G(P)|$ に現れる2の最高冪が $2^2=4$ だから Sylow $2$ 部分群である。

  \item[(3)] $g\in G$ と $Q\in\Omega$ に対して
  \[
  \varphi(g)(Q)=gQg^{-1}
  \]
  と定める。$Q$ が位数5なら $gQg^{-1}$ も位数5なので、$\varphi(g)$ は $\Omega$ の置換を定める。
  $g,h\in G$ に対して
  \begin{align*}
  \varphi(gh)(Q)
  &=(gh)Q(gh)^{-1}\\
  &=g(hQh^{-1})g^{-1}\\
  &=\varphi(g)(\varphi(h)(Q)).
  \end{align*}
  従って $\varphi(gh)=\varphi(g)\varphi(h)$ であり、
  $\varphi:G\to\operatorname{Sym}(\Omega)$ は群準同型である。

  \item[(4)] (2) より $|\Omega|=6$ なので、$\Omega$ と $\{1,\ldots,6\}$ を同一視すれば
  \[
  \varphi:S_5\longrightarrow S_6
  \]
  とみなせる。$K=\ker\varphi$ とおくと $K\triangleleft S_5$ である。
  $K$ の元は全ての $Q\in\Omega$ を固定し、特に $P$ を固定するから
  \[
  K\subset N_G(P).
  \]
  $S_5$ の正規部分群は $\{1\},A_5,S_5$ の三つである。
  一方 $|N_G(P)|=20$ なので、位数60の $A_5$ も位数120の $S_5$ も
  $N_G(P)$ に含まれない。従って $K=\{1\}$ であり、$\varphi$ は単射である。

  よって
  \[
  H_1:=\varphi(S_5)\subset S_6,
  \qquad |H_1|=120,
  \qquad [S_6:H_1]=\frac{720}{120}=6.
  \]
  Sylow 部分群は互いに共役なので、$H_1$ の6点集合 $\Omega$ への作用は推移的である。

  一方、通常の点固定部分群
  \[
  H_2:=\{\tau\in S_6\mid\tau(6)=6\}
  \]
  は $S_5$ と同型で、$|H_2|=120$、従って $[S_6:H_2]=6$ である。
  $H_2$ は点6を固定するが、$H_1$ は6点に推移的に作用するため固定点を持たない。
  もし $H_1=sH_2s^{-1}$ となる $s\in S_6$ が存在すれば、$H_1$ は点 $s(6)$ を固定することになり矛盾する。
  従って $H_1,H_2$ は指数6で互いに共役でない二つの部分群である。
\end{enumerate}
\end{proof}

\begin{problem}{第2問｜球面の面積形式}{2022-s-2}
$\R^3$ 内の球面 $S^2:=\{(x,y,z)\in\R^3\mid x^2+y^2+z^2=1\}$ を考える。点 $p\in S^2$ に対して，
\[
\omega_p(u,v):=\inner p{u\times v}
\]
と定める。ただし，$u,v$ は点 $p$ における接ベクトルであり，$u=(u_1,u_2,u_3)$, $v=(v_1,v_2,v_3)$ に対して $u\times v$ は
\[
u\times v=(u_2v_3-u_3v_2,\ u_3v_1-u_1v_3,\ u_1v_2-u_2v_1)
\]
で定義される外積である。また，$\inner{-}{-}$ は $\R^3$ の標準内積である。このとき，以下の問いに答えよ。
\begin{enumerate}
  \item[(1)] $\omega_p(u,v)=-\omega_p(v,u)$ が成り立つことおよび $\omega_p(u,v)=0$ が全ての接ベクトル $v$ で成り立つならば $u=0$ であることを示せ。
  \item[(2)] $S^2$ の開集合 $U$ を
  \[
  U=S^2\setminus\{(x,y,z)\in S^2\mid y=0,\ x\ge0\}
  \]
  で定める。$0<\theta<2\pi$, $-1<h<1$ に対し，
  \[
  (r(\theta,h)\cos\theta,\ r(\theta,h)\sin\theta,\ h)\in U
  \]
  となる連続関数 $r(\theta,h)$ を求めよ。
  \item[(3)] (2) で求めた $r(\theta,h)$ を用いて，写像 $T:(0,2\pi)\times(-1,1)\to U$ を
  \[
  T(\theta,h)=(r(\theta,h)\cos\theta,\ r(\theta,h)\sin\theta,\ h)
  \]
  で定める。$p=T(\theta,h)\in U$ に対して
  \[
  \omega_p\left(\frac{\partial T}{\partial\theta},\frac{\partial T}{\partial h}\right)
  \]
  を求めよ。
\end{enumerate}
\end{problem}
\begin{memo*}
\textbf{解答の見通し。}
(1) の反対称性は外積の順序交換から従う。非退化性では
$v=p\times u$ と選ぶと、三重積が $\norm{u}^2p$ になって $u$ の長さを取り出せる。
(2), (3) では高さ $h$ の緯線の半径を球面の方程式から決め、
座標ベクトルの外積が外向き法線 $T$ に一致することを確かめる。
\end{memo*}
\begin{proof}[解答]
\begin{enumerate}
  \item[(1)] 外積の反対称性 $v\times u=-(u\times v)$ より
  \[
  \omega_p(v,u)=\inner p{v\times u}
  =-\inner p{u\times v}=-\omega_p(u,v).
  \]
  次に $\omega_p(u,v)=0$ が全ての $v\in T_pS^2$ に対して成り立つと仮定する。
  $u\in T_pS^2$ だから $\inner pu=0$ であり、$v=p\times u$ も $p$ と直交するので
  $v\in T_pS^2$ である。ベクトル三重積の公式を用いると
  \[
  u\times(p\times u)
  =p\inner uu-u\inner up
  =\norm u^2p.
  \]
  従って
  \[
  0=\omega_p(u,p\times u)
  =\inner p{\norm u^2p}
  =\norm u^2\norm p^2=\norm u^2.
  \]
  よって $u=0$ である。

  \item[(2)] 点 $(r\cos\theta,r\sin\theta,h)$ が $S^2$ 上にある条件は
  \[
  r^2\cos^2\theta+r^2\sin^2\theta+h^2=1,
  \]
  すなわち $r^2=1-h^2$ である。$-1<h<1$ では右辺は正であり、
  $r$ を緯線の半径として正に取るので
  \[
  r(\theta,h)=\sqrt{1-h^2}.
  \]
  $r$ は $\theta$ に依存しない。また $0<\theta<2\pi$ としたことで、除いた半大円上の点は現れない。

  \item[(3)] 以下 $T_\theta=\partial T/\partial\theta$, $T_h=\partial T/\partial h$ と略記する。
  $r=\sqrt{1-h^2}$ だから $\partial r/\partial h=-h/r$ であり
  \[
  T_\theta=(-r\sin\theta,r\cos\theta,0),\qquad
  T_h=\left(-\frac h r\cos\theta,-\frac h r\sin\theta,1\right).
  \]
  外積を成分ごとに計算すると
  \begin{align*}
  T_\theta\times T_h
  &=\left(r\cos\theta,r\sin\theta,
    h(\sin^2\theta+\cos^2\theta)\right)\\
  &=(r\cos\theta,r\sin\theta,h)=T(\theta,h).
  \end{align*}
  従って $\omega_T(T_\theta,T_h)=\inner{T}{T}=1$。
\end{enumerate}
\end{proof}

\begin{problem}{第3問｜微分方程式}{2022-s-3}
$f(x)$ を $\R$ 上の有界かつ連続な関数とし，微分方程式
\[
\frac{\d^2y}{\d x^2}-y=f(x)
\tag{$*$}
\]
を考える。このとき，以下の問いに答えよ。
\begin{enumerate}
  \item[(1)] $f(x)=0$ $(x\in\R)$ とするとき，微分方程式 ($*$) の一般解を求めよ。
  \item[(2)] $a\in\R$ とする。このとき，$e^{-x}f(x)$ は区間 $[a,\infty)$ において広義積分可能であることを示せ。
  \item[(3)] $\R$ 上の関数 $g(x)$ を
  \[
  g(x)=-\frac{e^x}2\int_x^\infty e^{-t}f(t)\,\d t
  -\frac{e^{-x}}2\int_{-\infty}^xe^tf(t)\,\d t
  \qquad(x\in\R)
  \]
  によって定義する。このとき，$g(x)$ は微分方程式 ($*$) の解であることを示せ。
  \item[(4)] $\displaystyle\lim_{x\to\infty}f(x)=1$ を満たすとき，(3) で定義した関数 $g(x)$ は $\displaystyle\lim_{x\to\infty}g(x)=-1$ を満たすことを示せ。
  \item[(5)] 
  \[
  f(x)=\begin{cases}1-e^{-x}&(x\ge0)\\0&(x<0)\end{cases}
  \]
  のとき $\displaystyle\lim_{x\to\infty}y(x)=-1$ を満たす微分方程式 ($*$) の解を全て求めよ。
\end{enumerate}
\end{problem}
\begin{memo*}
\textbf{解答の見通し。}
最初に同次方程式の二つの基本解 $e^x,e^{-x}$ を確認する。
(3) の積分表示は二つの積分を別々に微分すると交差項が消えるよう設計されている。
(4) は $f=1+(f-1)$ と分けて優収束定理を使い、
(5) は $x<0$ と $x\ge0$ で解いた後、$x=0$ で $y,y'$ を接続する。
\end{memo*}
\begin{proof}[解答]
\begin{enumerate}
  \item[(1)] 特性方程式は
  \[
  \lambda^2-1=(\lambda-1)(\lambda+1)=0
  \]
  であり、相異なる特性根は $1,-1$ である。従って一般解は
  \[
  y(x)=C_1e^x+C_2e^{-x}\qquad(C_1,C_2\in\R)
  \]
  である。
  \item[(2)] $f$ は有界だから、ある $M\ge0$ が存在して
  $|f(x)|\le M$ が全ての $x$ で成り立つ。任意の $a\in\R$ に対して
  \[
  \int_a^\infty|e^{-x}f(x)|\,\d x
  \le M\int_a^\infty e^{-x}\,\d x
  =Me^{-a}<\infty.
  \]
  よって $e^{-x}f(x)$ は $[a,\infty)$ 上で絶対可積分であり、広義積分可能である。
  \item[(3)]
  \[
  A(x)=\int_x^\infty e^{-t}f(t)\,\d t,\qquad
  B(x)=\int_{-\infty}^x e^tf(t)\,\d t
  \]
  とおく。微積分学の基本定理より
  \[
  A'(x)=-e^{-x}f(x),\qquad B'(x)=e^xf(x).
  \]
  $g=-(e^xA+e^{-x}B)/2$ を積の微分法で微分すると
  \begin{align*}
  g'
  &=-\frac12\{e^xA+e^xA'-e^{-x}B+e^{-x}B'\}\\
  &=-\frac12e^xA+\frac12e^{-x}B,
  \end{align*}
  ここでは $e^xA'=-f$ と $e^{-x}B'=f$ が打ち消し合った。もう一度微分して
  \begin{align*}
  g''
  &=-\frac12(e^xA+e^xA')
    +\frac12(-e^{-x}B+e^{-x}B')\\
  &=-\frac12(e^xA+e^{-x}B)+f(x)\\
  &=g(x)+f(x).
  \end{align*}
  従って $g''-g=f$ となり、$g$ は微分方程式 ($*$) の解である。
  \item[(4)] $h=f-1$ とおくと $h(x)\to0$ で、$h$ は有界である。
  定数関数 $1$ を (3) の式へ代入して得られる解は
  \[
  -\frac{e^x}{2}\int_x^\infty e^{-t}\,\d t
  -\frac{e^{-x}}2\int_{-\infty}^x e^t\,\d t=-1.
  \]
  一方、$h$ が作る第1項は変数変換 $t=x+s$ により
  \[
  e^x\int_x^\infty e^{-t}h(t)\,\d t
  =\int_0^\infty e^{-s}h(x+s)\,\d s\longrightarrow0,
  \]
  第2項も $t=x-s$ により
  \[
  e^{-x}\int_{-\infty}^x e^th(t)\,\d t
  =\int_0^\infty e^{-s}h(x-s)\,\d s\longrightarrow0.
  \]
  いずれも有界性から $Ce^{-s}$ を優関数にして優収束定理を使った。
  従って $g(x)\to-1$。
  \item[(5)] $x\ge0$ では、定数項 $1$ の特殊解は $-1$。
  また $-e^{-x}$ は同次解 $e^{-x}$ と共振するので
  $y_p=Cxe^{-x}$ とおくと
  \[
  y_p''-y_p=-2Ce^{-x}.
  \]
  よって $C=1/2$。$x\to\infty$ で $-1$ に収束するには $e^x$ の項があってはならないから
  \[
  y(x)=-1+\left(D+\frac x2\right)e^{-x}\qquad(x\ge0).
  \]
  $x<0$ では $f=0$ なので $y=Be^x+Ce^{-x}$。
  $x=0$ で $y,y'$ を一致させると
  \[
  B+C=-1+D,\qquad B-C=\frac12-D,
  \]
  従って $B=-1/4$, $C=D-3/4$。以上より、任意の $D\in\R$ に対し
  \[
  y(x)=\begin{cases}
  -\frac14e^x+(D-\frac34)e^{-x},&x<0,\\
  -1+(D+\frac x2)e^{-x},&x\ge0
  \end{cases}
  \]
  であり、これが全てである。
\end{enumerate}
\end{proof}

\begin{problem}{第4問｜Dirichlet 積分}{2022-s-4}
以下の問いに答えよ。
\begin{enumerate}
  \item[(1)] $t$ を正の実数とする。$\dfrac{\sin x}x$ は開区間 $(0,t)$ 上でルベーグ積分可能であることを示せ。
  \item[(2)] $t,y$ を正の実数とする。等式
  \[
  \int_0^t e^{-xy}\sin x\,\d x
  =\frac1{1+y^2}-\frac{e^{-ty}}{1+y^2}(y\sin t+\cos t)
  \]
  を示せ。
  \item[(3)] $t$ を正の実数とする。等式
  \[
  \int_0^t\frac{\sin x}{x}\,\d x
  =\int_0^\infty\frac{\d y}{1+y^2}
  -\int_0^\infty\frac{e^{-ty}}{1+y^2}(y\sin t+\cos t)\,\d y
  \]
  \item[(4)] $\displaystyle\lim_{t\to\infty}\int_0^t\frac{\sin x}{x}\,\d x$ を求めよ。
\end{enumerate}
\end{problem}
\begin{memo*}
\textbf{解答の見通し。}
$\sin x/x$ は原点の除去可能特異点を埋めれば連続になる。
次に $1/x=\int_0^\infty e^{-xy}\,\d y$ を使って二重積分へ移し、
Fubini の定理で積分順序を交換する。
最後は明示式の誤差項が $1/t+1/t^2$ 以下になることを示せばよい。
\end{memo*}
\begin{proof}[解答]
\begin{enumerate}
  \item[(1)] $\lim_{x\to0}\sin x/x=1$ なので、$x=0$ で値を1と定めれば $[0,t]$ 上連続になる。
  従ってコンパクト区間上でルベーグ積分可能である。
  \item[(2)]
  \[
  F(x)=\frac{e^{-xy}(-y\sin x-\cos x)}{1+y^2}
  \]
  とおく。積の微分法により
  \begin{align*}
  F'(x)
  &=\frac{e^{-xy}}{1+y^2}
  \{(-y)(-y\sin x-\cos x)-y\cos x+\sin x\}\\
  &=e^{-xy}\sin x.
  \end{align*}
  よって
  \begin{align*}
  \int_0^t e^{-xy}\sin x\,\d x
  &=F(t)-F(0)\\
  &=\frac1{1+y^2}
   -\frac{e^{-ty}}{1+y^2}(y\sin t+\cos t).
  \end{align*}
  \item[(3)] $x>0$ では
  \[
  \frac1x=\int_0^\infty e^{-xy}\,\d y.
  \]
  また
  \[
  \int_0^t\int_0^\infty|e^{-xy}\sin x|\,\d y\d x
  =\int_0^t\frac{|\sin x|}{x}\,\d x<\infty
  \]
  なので Fubini の定理で積分順序を交換できる。(2) を用いて
  \begin{align*}
  \int_0^t\frac{\sin x}{x}\,\d x
  &=\int_0^\infty\int_0^t e^{-xy}\sin x\,\d x\d y\\
  &=\int_0^\infty\frac{\d y}{1+y^2}
  -\int_0^\infty\frac{e^{-ty}}{1+y^2}(y\sin t+\cos t)\,\d y.
  \end{align*}
  \item[(4)] 第2積分の絶対値は
  \[
  \int_0^\infty e^{-ty}(y+1)\,\d y
  =\frac1{t^2}+\frac1t\longrightarrow0
  \]
  以下である。一方
  \[
  \int_0^\infty\frac{\d y}{1+y^2}
  =[\arctan y]_0^\infty=\frac\pi2.
  \]
  従って求める極限は $\pi/2$。
\end{enumerate}
\end{proof}

\section{第2期・専門基礎科目（3題必修）}

\begin{problem}{第1問｜線形空間}{2022-2-1}
以下，$U,V$ は $\R$ 上の線形空間とする。
\begin{enumerate}
  \item[(1)] $\R$ 上の線形空間 $V$ に対して，次の問いに答えよ。\\
  (i) 線形空間 $V$ の部分空間 $W_1,W_2$ に対して，$W_1\cap W_2$ も $V$ の部分空間となることを示せ。\\
  (ii) 線形空間 $V$ の部分空間 $W_1,W_2$ に対して，$W_1\cup W_2$ が $V$ の部分空間となるとき，$W_1\subset W_2$ または $W_2\subset W_1$ であることを示せ。
  \item[(2)] $U,V$ は $\R$ 上の線形空間で，$f:U\to V$ は線形写像であるとする。$U,V$ の零ベクトルをそれぞれ $0_U,0_V$ で表す。$\ker f=\{x\in U\mid f(x)=0_V\}$ は $U$ の部分空間であることを示せ。また，$f$ が単射であるための必要十分条件は，$\ker f=\{0_U\}$ であることを示せ。
  \item[(3)] 実数係数の3次正方行列全体を $V=M_3(\R)$ とし，$A\in V$ が交代行列であるとは ${}^{\mathsf T}A=-A$ を満たすことをいう。交代行列全体の部分集合を $W=\{A\in V\mid {}^{\mathsf T}A=-A\}$ とする。$W$ が $V$ の部分空間であることを示せ。また $W$ の1組の基底を与え，$W$ の次元を求めよ。
\end{enumerate}
\end{problem}
\begin{memo*}
\textbf{解答の見通し。}
部分空間であることは、零ベクトルを含み、集合の任意の元 $x,y$ と任意の $a,b\in\R$ に対して
$ax+by$ も集合内に残ることを確認すればよい。
合併の問題では、どちらにも含まれないベクトルを一つずつ取ると、その和が矛盾を生む。
核 $\ker f$ は0へ写るベクトル全体である。単射の条件 $f(x)=f(y)\Rightarrow x=y$ を、
線形性により $f(x-y)=0$ へ直すと核との関係が見える。
交代行列では $A^{\mathsf T}=-A$ を各成分の式 $a_{ji}=-a_{ij}$ に直し、
対角成分と下三角成分が上三角成分から決まることを読む。
\end{memo*}
\begin{proof}[解答]
\begin{enumerate}
  \item[(1)] (i) $0\in W_1\cap W_2$ であり、$x,y\in W_1\cap W_2$、
  $a,b\in\R$ なら $ax+by\in W_1\cap W_2$。従って部分空間である。

  (ii) どちらの包含も成り立たないとすると
  $u\in W_1\setminus W_2$, $v\in W_2\setminus W_1$ を取れる。
  $u+v\in W_1$ なら $v\in W_1$、$u+v\in W_2$ なら $u\in W_2$ となる。
  従って $u+v\notin W_1\cup W_2$ であり、合併が部分空間であることに反する。

  \item[(2)] $x,y\in\ker f$, $a,b\in\R$ なら
  \[
  f(ax+by)=af(x)+bf(y)=0
  \]
  であり、$0_U\in\ker f$ だから $\ker f$ は部分空間である。また
  \[
  f(x)=f(y)\iff f(x-y)=0\iff x-y\in\ker f.
  \]
  よって $f$ が単射であることと $\ker f=\{0_U\}$ は同値である。

  \item[(3)] $O\in W$ であり、$A,B\in W$, $a,b\in\R$ なら
  \[
  (aA+bB)^{\mathsf T}=-aA-bB,
  \]
  従って $W$ は部分空間である。任意の $A\in W$ は一意に
  \[
  A=\begin{pmatrix}0&a&b\\-a&0&c\\-b&-c&0\end{pmatrix}
  =a(E_{12}-E_{21})+b(E_{13}-E_{31})+c(E_{23}-E_{32})
  \]
  と書ける。従って
  \[
  E_{12}-E_{21},\quad E_{13}-E_{31},\quad E_{23}-E_{32}
  \]
  が基底であり、$\dim W=3$。
\end{enumerate}
\end{proof}

\begin{problem}{第2問｜微分積分}{2022-2-2}
\begin{enumerate}
  \item[(1)] $a<b$ とする。閉区間の直積集合 $[a,b]\times[a,b]$ 上の連続関数 $u(s,t)$ と $x\in(a,b)$ に対し
  \[
  f(x)=\int_a^x\d s\int_a^su(s,t)\,\d t
  \]
  とおく。$f(x)$ は開区間 $(a,b)$ 上微分可能で，その導関数は
  \[
  \int_a^xu(x,t)\,\d t
  \]
  で与えられることを示せ。
  \item[(2)] $x>0$ に対し
  \[
  D_x=\{(s,t)\mid0\le s\le x,\ s\le t\le x\}
  \]
  とおく。このとき関数
  \[
  g(x)=\iint_{D_x}s(1-2st)\,\d s\d t
  \]
  の $x>0$ の範囲における最大値を求めよ。
\end{enumerate}
\end{problem}
\begin{memo*}
\textbf{解答の見通し。}
(1) では内側の積分全体を $F(s)=\int_a^s u(s,t)\,\d t$ と名付ける。
すると $f(x)=\int_a^xF(s)\,\d s$ なので、$F$ が連続だと確認できれば
微積分学の基本定理から $f'(x)=F(x)$ となる。
$F$ では積分区間と被積分関数の両方が $s$ とともに変わるため、差を二つに分けて連続性を示す。
(2) の不等式 $0\le s\le t\le x$ は、外側で $t$ を $0$ から $x$ まで動かし、
各 $t$ を固定した内側で $s$ を $0$ から $t$ まで動かす、という積分順序を表す。
\end{memo*}
\begin{proof}[解答]
\begin{enumerate}
  \item[(1)] $F(s)=\int_a^s u(s,t)\,\d t$ とおく。
  $u$ の一様連続性と有界性から、$h\to0$ のとき
  \[
  F(s+h)-F(s)
  =\int_a^s\{u(s+h,t)-u(s,t)\}\,\d t
   +\int_s^{s+h}u(s+h,t)\,\d t\longrightarrow0.
  \]
  従って $F$ は連続であり、微積分学の基本定理から
  \[
  f'(x)=F(x)=\int_a^x u(x,t)\,\d t.
  \]

  \item[(2)] $D_x=\{(s,t)\mid0\le s\le t\le x\}$ より
  \begin{align*}
  g(x)
  &=\int_0^x\int_0^t(s-2s^2t)\,\d s\d t\\
  &=\int_0^x\left(\frac{t^2}{2}-\frac{2t^4}{3}\right)\d t
  =\frac{x^3}{6}-\frac{2x^5}{15}.
  \end{align*}
  \[
  g'(x)=x^2\left(\frac12-\frac{2x^2}{3}\right)
  \]
  だから、$g$ は $x=\sqrt3/2$ で最大となる。その値は
  \[
  g\left(\frac{\sqrt3}{2}\right)=\frac{\sqrt3}{40}.
  \]
\end{enumerate}
\end{proof}

\begin{problem}{第3問｜位相空間}{2022-2-3}
\begin{enumerate}
  \item[(1)] $X$ を位相空間とし，$A,B$ を $X$ の部分集合とする。このとき $\overline{A\cup B}=\overline A\cup\overline B$ を示せ。但し $\overline A,\overline B,\overline{A\cup B}$ はそれぞれ $A,B,A\cup B$ の閉包である。
  \item[(2)] ハウスドルフ空間 $Y$ のコンパクトな部分集合 $K$ は閉集合であることを示せ。
  \item[(3)] $S,T$ を位相空間とし，$F:S\to T$ を全射な連続写像とする。$S$ が連結ならば，$T$ も連結であることを示せ。
\end{enumerate}
\end{problem}
\begin{memo*}
\textbf{解答の見通し。}
(1) の閉包 $\overline A$ は「$A$ を含む最小の閉集合」である。
二つの集合が等しいことは、左が右に含まれることと、右が左に含まれることを別々に示せばよい。
(2) の Hausdorff 性は異なる二点を交わらない開近傍で分けられる性質、
コンパクト性は任意の開被覆から有限個を選べる性質である。
点 $x\notin K$ と各 $y\in K$ をまず個別に分け、有限個の近傍だけで $K$ 全体を覆う。
(3) の連結とは、空でない二つの開集合へ分割できないことである。
像が分割できると仮定し、連続写像の逆像を取ると定義域まで分割されることを示す。
\end{memo*}
\begin{proof}[解答]
\begin{enumerate}
  \item[(1)] $\overline A\cup\overline B$ は $A\cup B$ を含む閉集合だから
  \[
  \overline{A\cup B}\subset\overline A\cup\overline B.
  \]
  一方、$A,B\subset A\cup B$ より
  $\overline A,\overline B\subset\overline{A\cup B}$。従って等号が成り立つ。

  \item[(2)] $x\notin K$ とする。各 $y\in K$ に対し Hausdorff 性より
  \[
  x\in U_y,\quad y\in V_y,\quad U_y\cap V_y=\varnothing
  \]
  となる開集合を取る。$K$ のコンパクト性より
  $K\subset V_{y_1}\cup\cdots\cup V_{y_m}$ とできる。このとき
  \[
  U=U_{y_1}\cap\cdots\cap U_{y_m}
  \]
  は $x$ の開近傍で $U\cap K=\varnothing$。従って $X\setminus K$ は開であり、
  $K$ は閉である。

  \item[(3)] $T=U\cup V$ が非自明な開分離だとする。連続性と全射性より
  $F^{-1}(U),F^{-1}(V)$ は空でない $S$ の開集合であり、
  \[
  S=F^{-1}(U)\cup F^{-1}(V),\qquad
  F^{-1}(U)\cap F^{-1}(V)=\varnothing.
  \]
  これは $S$ の連結性に反する。従って $T$ は連結である。
\end{enumerate}
\end{proof}

\chapter{2021年度}
\label{ch:2021}

\section{第1期・専門基礎科目（3題必修）}

\begin{problem}{第1問｜Jordan 標準形}{2021-1-1}
$a\in\mathbb C$ とし
\[
A=\begin{pmatrix}a-1&a&0\\-a&-a-1&0\\-1&1&-1\end{pmatrix}.
\]
固有多項式，固有値と固有空間の次元，Jordan 標準形を求めよ。
\end{problem}
\begin{memo*}
\textbf{解答の見通し。}
行列 $A$ の対角成分には全て $-1$ が含まれているので、$N=A+I_3$ とおいて
$-I_3$ の部分を取り除く。実際に掛け算すると $N^3=O$ になる。
このように何乗かすると零行列になる行列を冪零行列という。
$Nx=\mu x$ なら $N^3x=\mu^3x=0$ なので、$N$ の固有値は0だけである。
従って $A=N-I_3$ の固有値は $-1$ だけになる。

Jordan ブロック $J_k(-1)$ に $I_k$ を足すと、上の対角線に1が並ぶ冪零ブロックになる。
その $k$ 乗で初めて0になるため、$N^2$ と $N^3$ を調べれば最大ブロックの大きさが分かる。
$a=0$ のときだけ $N^2$ が0になるので場合分けする。
\end{memo*}
\begin{proof}[解答]
$N=A+I_3$ とおくと、$A$ の対角成分に $1$ を足して
\[
N=\begin{pmatrix}a&a&0\\-a&-a&0\\-1&1&0\end{pmatrix}
\]
である。この行列の冪を計算する。第1行 $(a,a,0)$ と $N$ の第1列 $(a,-a,-1)^{\mathsf T}$ の積は
$a^2-a^2+0=0$ であり、第2列 $(a,-a,1)^{\mathsf T}$ との積も $a^2-a^2+0=0$、
第3列は零ベクトルなので0である。第2行 $(-a,-a,0)$ は第1行の $-1$ 倍だから、
その行も全て0になる。第3行 $(-1,1,0)$ については
\[
(-1)a+1\cdot(-a)+0=-2a,\qquad
(-1)a+1\cdot(-a)+0=-2a
\]
となる。従って
\[
N^2=\begin{pmatrix}0&0&0\\0&0&0\\-2a&-2a&0\end{pmatrix}.
\]
さらに $N^3=N^2N$ の第3行は $(-2a,-2a,0)$ と $N$ の各列の積であり、
第1列では $-2a\cdot a+(-2a)(-a)=-2a^2+2a^2=0$、第2列でも同様に0、第3列は0である。
他の行はすでに0なので
\[
N^3=O
\]
である。

$Nx=\mu x$ を満たす $x\ne0$ があれば
\[
0=N^3x=\mu^3x
\]
だから $\mu^3=0$、すなわち $\mu=0$ である。従って $N$ の固有値は0だけであり、
3次の特性多項式は根を重複度込みで全て並べたものだから $|\mu I_3-N|=\mu^3$ となる。
$A=N-I_3$ なので $\lambda I_3-A=(\lambda+1)I_3-N$ であり、上の式へ $\mu=\lambda+1$ を代入して
\[
|\lambda I_3-A|=(\lambda+1)^3.
\]
従って固有値は $-1$ のみ（重複度3）である。

$a\ne0$ の場合を見る。このとき $N^2\ne O$ かつ $N^3=O$ である。
大きさ $k$ の冪零 Jordan ブロックは $k$ 乗で初めて0になるので、
最大の Jordan ブロックの大きさは3である。
全体の大きさも3だから、標準形はブロック一つの $J_3(-1)$ に限る。
固有空間の次元も直接分かる。$Nx=0$ は
\[
ax_1+ax_2=0,\qquad -ax_1-ax_2=0,\qquad -x_1+x_2=0
\]
であり、$a\ne0$ より第1式は $x_1+x_2=0$、第3式は $x_1=x_2$ だから $x_1=x_2=0$、
$x_3$ は自由である。従って
\[
\ker(A+I_3)=\ker N=\left\langle(0,0,1)^{\mathsf T}\right\rangle,
\qquad \dim\ker(A+I_3)=1.
\]

$a=0$ なら
\[
N=\begin{pmatrix}0&0&0\\0&0&0\\-1&1&0\end{pmatrix},
\]
従って $N\ne O$, $N^2=O$ である。
$x=(x_1,x_2,x_3)^{\mathsf T}$ とすると
\[
Nx=\begin{pmatrix}0\\0\\-x_1+x_2\end{pmatrix}=0
\quad\Longleftrightarrow\quad x_2=x_1.
\]
従って全ての解は
\[
x=x_1\begin{pmatrix}1\\1\\0\end{pmatrix}
 +x_3\begin{pmatrix}0\\0\\1\end{pmatrix}
\]
と表され、自由に選べる数が $x_1,x_3$ の二つあるので
$\dim\ker N=2$ である。
$N\ne O$ なので大きさ2のブロックが少なくとも一つ必要である。
大きさ $k$ の冪零 Jordan ブロック
\[
J_k(0)=\begin{pmatrix}
0&1&&0\\
&0&\ddots&\\
&&\ddots&1\\
0&&&0
\end{pmatrix}
\]
では $J_k(0)y=0$ の解は $y=(y_1,0,\ldots,0)^{\mathsf T}$ であり、
各ブロックが核の自由な数を一つずつ作る。従って
$\dim\ker N=2$ は Jordan ブロックが2個あることを表す。
全体の大きさは3だから、ブロックの大きさは2と1になり、標準形は
$J_2(-1)\oplus(-1)$ となる。
\end{proof}

\begin{problem}{第2問｜二変数関数}{2021-1-2}
\[
f(x,y)=\begin{cases}\dfrac{\sin(x^2y)}{x^2+y^2},&(x,y)\ne(0,0),\\0,&(x,y)=(0,0).
\end{cases}
\]
\begin{enumerate}
  \item[(1)] 極座標で $|f(x,y)|\le r$ および $|f(x,y)|\le r^{-2}$ を示せ。
  \item[(2)] 原点で連続であることを示せ。
  \item[(3)] $p>1$，$B(R)=\{x^2+y^2\le R^2\}$ とし $I_p(R)=\iint_{B(R)}|f|^p\,\d x\d y$ とおく。$\lim_{R\to\infty}I_p(R)$ の存在を示せ。
  \item[(4)] $\partial f/\partial y(x,0)$ を求め，$f$ が $C^1$ 級か判定せよ。
\end{enumerate}
\end{problem}
\begin{memo*}
\textbf{解答の見通し。}
極座標の半径 $r=\sqrt{x^2+y^2}$ は、点 $(x,y)$ と原点の距離であるため
$|x|\le r$, $|y|\le r$ が成り立つ。
$|\sin u|\le|u|$ を使うと $|f|\le r$ となり、原点へ近づくと0へ近づくことが分かる。
一方、$|\sin u|\le1$ を使うと $|f|\le r^{-2}$ となり、遠方での減衰が分かる。
半径1の内側と外側で二つの評価を使い分け、面積要素
$\d x\d y=r\,\d r\d\theta$ を掛けた一変数積分が収束するか確認する。
最後の偏導関数は公式を先に使わず、定義の差商を直接計算する。
\end{memo*}
\begin{proof}[解答]
\begin{enumerate}
  \item[(1)] $r=\sqrt{x^2+y^2}>0$ とおくと $|x|\le r$, $|y|\le r$ である。
  $|\sin u|\le|u|$ を $u=x^2y$ に使えば
  \[
  |f(x,y)|
  \le\frac{|x|^2|y|}{x^2+y^2}
  \le\frac{r^3}{r^2}=r.
  \]
  また $|\sin u|\le1$ より
  \[
  |f(x,y)|\le\frac1{x^2+y^2}=\frac1{r^2}.
  \]
  \item[(2)] 任意の $\varepsilon>0$ に対して $\delta=\varepsilon$ と取る。
  $0<\sqrt{x^2+y^2}=r<\delta$ なら (1) より
  \[
  |f(x,y)-f(0,0)|=|f(x,y)|\le r<\varepsilon.
  \]
  従って $f$ は原点で連続である。
  \item[(3)] $r\le1$ では $|f|^p\le r^p$、$r\ge1$ では
  $|f|^p\le r^{-2p}$ である。極座標の面積要素は
  $\d x\d y=r\,\d r\d\theta$ だから
  \begin{align*}
  \iint_{r\le1}|f|^p\,\d x\d y
  &\le\int_0^{2\pi}\int_0^1r^p\cdot r\,\d r\d\theta\\
  &=2\pi\int_0^1r^{p+1}\,\d r
  =2\pi\left[\frac{r^{p+2}}{p+2}\right]_0^1
  =\frac{2\pi}{p+2}<\infty,
  \end{align*}
  \begin{align*}
  \iint_{r\ge1}|f|^p\,\d x\d y
  &\le\int_0^{2\pi}\int_1^\infty r^{-2p}\cdot r\,\d r\d\theta\\
  &=2\pi\int_1^\infty r^{1-2p}\,\d r
  =2\pi\left[\frac{r^{2-2p}}{2-2p}\right]_1^\infty
  =\frac{2\pi}{2p-2}<\infty.
  \end{align*}
  最後の広義積分では、上端で $r^{2-2p}\to0$ となるために指数の条件
  $2-2p<0$、すなわち $p>1$ が必要であり、仮定はこれを満たしている。
  従って $|f|^p\in L^1(\R^2)$。また
  $R$ を増やすと円板 $B(R)$ は広がるので、非負関数
  $|f|^p\mathbf1_{B(R)}$ は各点で $|f|^p$ へ単調に増加する。従って単調収束定理より
  \[
  \lim_{R\to\infty}I_p(R)=\iint_{\R^2}|f|^p\,\d x\d y<\infty.
  \]
  \item[(4)] 偏微分を計算すると
  \[
  \frac{\partial f}{\partial y}(x,0)=\begin{cases}1,&x\ne0,\\0,&x=0.
  \end{cases}
  \]
  実際、$x\ne0$ では偏微分の定義から
  \begin{align*}
  \frac{\partial f}{\partial y}(x,0)
  &=\lim_{h\to0}\frac{f(x,h)-f(x,0)}h\\
  &=\lim_{h\to0}\frac{\sin(x^2h)}{h(x^2+h^2)}\\
  &=\lim_{h\to0}
    \frac{\sin(x^2h)}{x^2h}\frac{x^2}{x^2+h^2}=1.
  \end{align*}
  最後の等号では $\sin u/u\to1$（$u=x^2h\to0$）と
  $x^2/(x^2+h^2)\to1$（$x\ne0$ を固定）を使った。
  $x=0$ では、全ての $h$ について $f(0,h)=\sin0/h^2=0$ かつ $f(0,0)=0$ なので、
  同じ差商は0であり $\partial f/\partial y(0,0)=0$ である。

  $C^1$ 級とは、全ての1階偏導関数が存在して連続になることである。
  いま $x\to0$（$x\ne0$）とすると
  \[
  \frac{\partial f}{\partial y}(x,0)=1
  \longrightarrow1
  \ne0=\frac{\partial f}{\partial y}(0,0)
  \]
  だから、$\partial f/\partial y$ は原点で不連続である。従って $f$ は $C^1$ 級ではない。
\end{enumerate}
\end{proof}

\begin{problem}{第3問｜連続像}{2021-1-3}
\begin{enumerate}
  \item[(1)] 集合と写像 $f:X\to Y$, $g:Y\to Z$ で，$f$ は全射でなく，$g$ は単射でなく，$g\circ f$ は全単射となる例を挙げよ。
  \item[(2)] 連続写像による連結集合の像は連結か。
  \item[(3)] 連続写像によるコンパクト集合の像はコンパクトか。
\end{enumerate}
\end{problem}
\begin{memo*}
\textbf{解答の見通し。}
(1) で、全射は「終域の全ての元が像に現れる」、単射は
「異なる二点が同じ値へ写らない」、全単射は両方を満たすことを意味する。
一元集合と二元集合を使えば、各条件を元を列挙して確認できる。
(2) の連結とは、空でない二つの相対開集合へ分割できないことである。
像が分割できると仮定し、その二集合の逆像を取ると元の集合まで分割される。
(3) のコンパクト性は、任意の開被覆から有限部分被覆を選べることとして使う。
像の開被覆を逆像へ戻して定義域を覆い、選んだ有限個を再び像へ戻す。
\end{memo*}
\begin{proof}[解答]
\begin{enumerate}
  \item[(1)] $X=Z=\{0\}$、$Y=\{0,1\}$ とし
  \[
  f(0)=0,\qquad g(0)=g(1)=0
  \]
  と定める。$1\notin f(X)$ なので $f$ は全射でなく、$g(0)=g(1)$ なので $g$ は単射でない。
  一方、$g\circ f:X\to Z$ は一元集合から一元集合への写像なので全単射である。

  \item[(2)] 連結である。連結集合 $C$ と連続写像 $h:C\to Y$ を取り、$h(C)$ が連結でないと仮定する。
  このとき $h(C)$ の相対位相で開な非空集合 $U,V$ が存在して
  \[
  h(C)=U\sqcup V
  \]
  と書ける。$U=h(C)\cap U'$, $V=h(C)\cap V'$ となる $Y$ の開集合 $U',V'$ を取れば
  \[
  C=h^{-1}(U')\sqcup h^{-1}(V')
  \]
  となる。右辺の二集合は連続性により $C$ で開であり、$U,V$ が非空なので、ともに非空である。
  これは $C$ の連結性に矛盾する。従って $h(C)$ は連結である。

  \item[(3)] コンパクトである。コンパクト集合 $K$ と連続写像 $h:K\to Y$ を取り、
  $h(K)$ の開被覆を $\{U_\lambda\}_{\lambda\in\Lambda}$ とする。このとき
  \[
  K\subset\bigcup_{\lambda\in\Lambda}h^{-1}(U_\lambda)
  \]
  であり、各 $h^{-1}(U_\lambda)$ は $K$ で開である。$K$ のコンパクト性より有限個
  $\lambda_1,\ldots,\lambda_m$ を選んで
  \[
  K\subset\bigcup_{j=1}^m h^{-1}(U_{\lambda_j})
  \]
  とできる。両辺を $h$ で写せば
  \[
  h(K)\subset\bigcup_{j=1}^m U_{\lambda_j}
  \]
  となるので、$h(K)$ はコンパクトである。
\end{enumerate}
\end{proof}

\section{第1期・専門科目（4題から2題選択）}

\begin{problem}{第1問｜環とイデアル}{2021-s-1}
$\Phi:\mathbb C[X,Y,Z]\to\mathbb C[t]$ を
$\Phi(f)=f(t^2,t^3,t^4)$ で定め，$I=\ker\Phi$，$R=\mathbb C[X,Y,Z]/I$ とする。
\begin{enumerate}
  \item[(1)] $I$ は素イデアルか。
  \item[(2)] $I$ の生成元を求めよ。
  \item[(3)] $\bar X,\bar Y,\bar Z$ で生成される $R$ のイデアル $\mathfrak m$ が極大であることを示せ。
  \item[(4)] $\mathfrak m$ が2元で生成されることを示せ。
  \item[(5)] $\mathfrak m$ が単項イデアルでないことを示せ。
\end{enumerate}
\end{problem}
\begin{memo*}
\textbf{解答の見通し。}
準同型の像を $\mathbb C[t^2,t^3]$ と同定して商環を具体化する。
核ではまず $Z=X^2$ を消去し、次に $Y^2=X^3$ で割った余りを偶数次数・奇数次数に分ける。
極大性は原点での値写像、生成元の最小個数は
$\mathfrak m/\mathfrak m^2$ のベクトル空間次元で判定する。
\end{memo*}
\begin{proof}[解答]
\begin{enumerate}
  \item[(1)] $t^4=(t^2)^2$ だから
  \[
  \operatorname{Im}\Phi=\mathbb C[t^2,t^3,t^4]
  =\mathbb C[t^2,t^3].
  \]
  第1同型定理より
  \[
  R=\mathbb C[X,Y,Z]/\ker\Phi
  \cong\mathbb C[t^2,t^3]\subset\mathbb C[t].
  \]
  $\mathbb C[t]$ は整域なので、その部分環 $\mathbb C[t^2,t^3]$ も整域である。
  従って $I=\ker\Phi$ は素イデアルである。
  \item[(2)] 生成元は
  \[
  I=(Z-X^2,\ Y^2-X^3).
  \]
  実際、$\Phi(Z-X^2)=t^4-(t^2)^2=0$、
  $\Phi(Y^2-X^3)=t^6-(t^2)^3=0$ だから、右辺を $J$ とおけば $J\subset I$。

  逆に $f\in I$ とする。$Z-X^2$ を用いて $Z$ を消去すれば、ある
  $g(X,Y)\in\mathbb C[X,Y]$ により
  \[
  f\equiv g(X,Y)\pmod{(Z-X^2)}
  \]
  と書ける。さらに $g$ を $Y$ の多項式と見て、モニック多項式
  $Y^2-X^3$ で割ると
  \[
  g(X,Y)=q(X,Y)(Y^2-X^3)+a(X)Y+b(X)
  \]
  となる。$f\in I$ より
  \[
  0=\Phi(f)=a(t^2)t^3+b(t^2).
  \]
  $a(t^2)t^3$ は奇数次数の項だけ、$b(t^2)$ は偶数次数の項だけからなるので、
  この多項式恒等式から $a=b=0$。従って $f\in J$ であり、$I=J$ を得る。

  \item[(3)] $R\cong\mathbb C[t^2,t^3]$ のもとで
  $\mathfrak m=(t^2,t^3)$ に対応する。定数項を取る準同型
  \[
  \mathbb C[t^2,t^3]\longrightarrow\mathbb C,\qquad h(t)\longmapsto h(0)
  \]
  は全射で、核は $(t^2,t^3)$ である。従って
  $R/\mathfrak m\cong\mathbb C$ は体だから $\mathfrak m$ は極大である。

  \item[(4)] $\bar Z=\bar X^2$ だから
  \[
  \mathfrak m=(\bar X,\bar Y,\bar Z)
  =(\bar X,\bar Y,\bar X^2)=(\bar X,\bar Y).
  \]

  \item[(5)] $R\cong\mathbb C[t^2,t^3]$ で考えると
  \[
  \mathfrak m=(t^2,t^3),\qquad
  \mathfrak m^2=(t^4,t^5,t^6).
  \]
  従って $\mathfrak m/\mathfrak m^2$ では $t^2,t^3$ の剰余類が一次独立であり、
  これらが全体を生成するから
  \[
  \dim_{R/\mathfrak m}(\mathfrak m/\mathfrak m^2)
  =\dim_{\mathbb C}(\mathfrak m/\mathfrak m^2)=2.
  \]
  もし $\mathfrak m=(h)$ が単項なら、$\mathfrak m/\mathfrak m^2$ は
  $h+\mathfrak m^2$ 一つで生成される $R/\mathfrak m$ ベクトル空間なので次元は高々1となり、矛盾する。
\end{enumerate}
\end{proof}

\begin{problem}{第2問｜曲率・捩率}{2021-s-2}
$p(t)=(\cos t,\sin t,\cosh t)$ と円柱面 $S=\{x^2+y^2=1\}$ を考える。
\begin{enumerate}
  \item[(1)] $p$ の曲率と捩率を求めよ。
  \item[(2)] $S$ の主方向と $p'(t)$ の角度が $\pi/4$ となる $t$ を全て求めよ。
  \item[(3)] $p$ の $S$ における測地的曲率ベクトルを求めよ。
\end{enumerate}
\end{problem}
\begin{memo*}
\textbf{解答の見通し。}
(1) は $p',p'',p'''$ を曲率・捩率の公式へ代入する。
(2) では円柱の二つの主方向へ $p'$ を分解し、内積から角度条件を作る。
(3) は単位接ベクトルを弧長で微分して曲率ベクトルを求め、
そこから円柱の法線方向成分を引けば測地的曲率ベクトルになる。
\end{memo*}
\begin{proof}[解答]
まず
\[
p'=(-\sin t,\cos t,\sinh t),\quad
p''=(-\cos t,-\sin t,\cosh t),
\]
\[
p'''=(\sin t,-\cos t,\sinh t)
\]
であり
\[
p'\times p''=(\cos t\cosh t+\sin t\sinh t,
\sin t\cosh t-\cos t\sinh t,1).
\]
従って
\[
\norm{p'}=\cosh t,\qquad
\norm{p'\times p''}=\sqrt2\cosh t,\qquad
(p'\times p'')\cdot p'''=2\sinh t.
\]
曲率・捩率の公式に代入すると
\[
\kappa(t)=\frac{\sqrt2}{\cosh^2t},\qquad
\tau(t)=\frac{\sinh t}{\cosh^2t}.
\]
円柱の単位主方向を
\[
e_1=(-\sin t,\cos t,0),\qquad e_2=(0,0,1)
\]
とする。$p'=e_1+\sinh t\,e_2$ かつ $\norm{p'}=\cosh t$ なので、
$p'$ と $e_1$ のなす角が $\pi/4$ である条件は
\[
\frac{|p'\cdot e_1|}{\norm{p'}}=\frac1{\cosh t}=\frac1{\sqrt2}.
\]
$e_2$ との角についても $|\sinh t|/\cosh t=1/\sqrt2$ となり、同じ条件
$\cosh t=\sqrt2$ を与える。従って
$t=\pm\operatorname{arcosh}\sqrt2=\pm\log(1+\sqrt2)$。

単位接ベクトルは
\[
T=\frac{p'}{\norm{p'}}
=(-\operatorname{sech}t\sin t,
  \operatorname{sech}t\cos t,\tanh t).
\]
$\d s/\d t=\norm{p'}=\cosh t$ より $\d/\d s=\operatorname{sech}t\,\d/\d t$。
従って成分ごとに微分すると
\[
\frac{\d T}{\d s}
=\operatorname{sech}t\frac{\d T}{\d t}
\]
\[
=\bigl(
\operatorname{sech}^2t\tanh t\sin t-\operatorname{sech}^2t\cos t,
-\operatorname{sech}^2t\tanh t\cos t-\operatorname{sech}^2t\sin t,
\operatorname{sech}^3t
\bigr).
\]
円柱の単位法線は $n=(\cos t,\sin t,0)$ であり
\begin{align*}
\frac{\d T}{\d s}\cdot n
&=\operatorname{sech}^2t\tanh t\sin t\cos t
-\operatorname{sech}^2t\cos^2t\\
&\quad-\operatorname{sech}^2t\tanh t\cos t\sin t
-\operatorname{sech}^2t\sin^2t\\
&=-\operatorname{sech}^2t.
\end{align*}
従って法線成分を引けば
\[
\boldsymbol{k}_g
=\frac{\d T}{\d s}
-\left(\frac{\d T}{\d s}\cdot n\right)n
=\operatorname{sech}^2t\,
(\tanh t\sin t,-\tanh t\cos t,\operatorname{sech}t).
\]
\end{proof}

\begin{problem}{第3問｜収束べき級数}{2021-s-3}
$\varphi(z)=\sum_{n=0}^\infty a_nz^n$ は収束べき級数で $a_0\ne0$ とする。
\begin{enumerate}
  \item[(1)] $1/\varphi(z)$ が $z=0$ で収束べき級数に展開できることを示せ。
  \item[(2)] その収束半径が $\varphi$ の収束半径より真に小さくなることはあるか。
  \item[(3)] $1/\varphi(z)=\sum b_nz^n$ とするとき $b_0,b_1,b_2$ を求めよ。
  \item[(4)] $n=0,1,2,3$ に対し $\displaystyle\int_{|z|=1}\frac{\d z}{z^n\sin^2z}$ を求めよ。
\end{enumerate}
\end{problem}
\begin{memo*}
\textbf{解答の見通し。}
$\varphi(0)\ne0$ なら逆数は原点近くで正則だが、その収束半径は
$\varphi$ 自身の収束半径ではなく、逆数が最初に出会う特異点で決まる。
係数は積が1になる条件から再帰的に求める。
積分では $1/\sin^2z$ を必要な次数まで Laurent 展開し、$z^{-1}$ の係数だけを読む。
\end{memo*}
\begin{proof}[解答]
\begin{enumerate}
\item[(1)] $\varphi(0)\ne0$ なので、連続性により0の十分小さい近傍で $\varphi(z)\ne0$。
従って $1/\varphi$ はその近傍で正則であり、Taylor 展開できる。
\item[(2)] ある。例えば $\varphi(z)=1-2z$ は整関数だが
$1/\varphi(z)=\sum_{n=0}^\infty(2z)^n$ の収束半径は $1/2$。
\item[(3)] 恒等式
\[
(a_0+a_1z+a_2z^2+\cdots)(b_0+b_1z+b_2z^2+\cdots)=1
\]
で $z^0,z^1,z^2$ の係数を比較すると
\[
a_0b_0=1,\quad a_0b_1+a_1b_0=0,\quad
a_0b_2+a_1b_1+a_2b_0=0.
\]
これを順に解いて
\[
b_0=\frac1{a_0},\quad b_1=-\frac{a_1}{a_0^2},\quad
b_2=\frac{a_1^2-a_0a_2}{a_0^3}.
\]
\item[(4)]
\[
\sin z=z-\frac{z^3}{6}+\frac{z^5}{120}+O(z^7)
\]
を二乗すると
\begin{align*}
\sin^2z
&=z^2-\frac13z^4
+\left(\frac1{36}+\frac1{60}\right)z^6+O(z^8)\\
&=z^2\left(1-\frac13z^2+\frac{2}{45}z^4+O(z^6)\right).
\end{align*}
$u=-z^2/3+2z^4/45+O(z^6)$ とおき
$(1+u)^{-1}=1-u+u^2+O(u^3)$ を使うと
\begin{align*}
\left(1-\frac13z^2+\frac{2}{45}z^4+O(z^6)\right)^{-1}
&=1+\frac13z^2
+\left(\frac19-\frac{2}{45}\right)z^4+O(z^6)\\
&=1+\frac13z^2+\frac1{15}z^4+O(z^6).
\end{align*}
従って
\[
\frac1{\sin^2z}=\frac1{z^2}+\frac13+\frac{z^2}{15}+O(z^4).
\]
従って
\[
\frac1{z^n\sin^2z}
=z^{-n-2}+\frac13z^{-n}+\frac1{15}z^{2-n}+O(z^{4-n}).
\]
$z^{-1}$ の係数、すなわち留数を $n=0,1,2,3$ の順に読むと
$0,1/3,0,1/15$。留数定理から積分値は順に
\[
0,\qquad \frac{2\pi i}{3},\qquad0,\qquad\frac{2\pi i}{15}.
\]
\end{enumerate}
\end{proof}

\begin{problem}{第4問｜測度収束}{2021-s-4}
$E\subset\R^d$ は可測で $|E|<\infty$，$f_n$ は $E$ 上可測とする。$A_{n,\varepsilon}=\{x\in E\mid|f_n(x)|\ge\varepsilon\}$ とおく。
\begin{enumerate}
  \item[(1)] $g\in L^1(E)$，$F=\{|g|\ge1\}$ なら $|F|\le\int_E|g|$ を示せ。
  \item[(2)] $\int_E|f_n|\to0$ なら全ての $\varepsilon>0$ で $|A_{n,\varepsilon}|\to0$ を示せ。
  \item[(3)] $f_n\to0$ がほとんど至る所なら $\displaystyle\int_E\frac{|f_n|}{1+|f_n|}\to0$ を示せ。
  \item[(4)] 全ての $\varepsilon>0$ で $|A_{n,\varepsilon}|\to0$ なら，ほとんど至る所で0に収束する部分列があることを示せ。
\end{enumerate}
\end{problem}
\begin{memo*}
\textbf{解答の見通し。}
(1), (2) は集合の定義関数を積分して Chebyshev 型の評価を得る。
(3) は被積分関数が1以下なので、有限測度という仮定が優関数を与える。
(4) は閾値を $1/k$、例外集合の測度を $2^{-k}$ とする部分列を選び、
例外集合へ無限回入る点が零集合であることを使う。
\end{memo*}
\begin{proof}[解答]
\begin{enumerate}
\item[(1)] $F=\{|g|\ge1\}$ 上では $1\le|g|$ だから、$E$ 上で
$\mathbf1_F\le|g|$ である。両辺を積分して
\[
|F|=\int_E\mathbf1_F\le\int_E|g|.
\]

\item[(2)] $A_{n,\varepsilon}$ 上では $|f_n|/\varepsilon\ge1$ なので
\[
\mathbf1_{A_{n,\varepsilon}}
\le\frac{|f_n|}{\varepsilon}.
\]
従って
\[
|A_{n,\varepsilon}|
\le\frac1\varepsilon\int_E|f_n|\longrightarrow0.
\]

\item[(3)]
\[
0\le\frac{|f_n(x)|}{1+|f_n(x)|}\le1
\]
であり、$f_n(x)\to0$ となる点では左辺も0へ収束する。
$|E|<\infty$ なので優関数 $\mathbf1_E$ は可積分であり、優収束定理から
\[
\int_E\frac{|f_n|}{1+|f_n|}\longrightarrow0.
\]

\item[(4)] 各 $k$ について $|A_{n,1/k}|\to0$ だから、$n_1<n_2<\cdots$ を
\[
|A_{n_k,1/k}|<2^{-k}
\]
となるよう帰納的に選べる。すると
\[
\sum_{k=1}^\infty|A_{n_k,1/k}|
\le\sum_{k=1}^\infty2^{-k}<\infty.
\]
Borel--Cantelli の補題より、ほとんど全ての $x$ は
$A_{n_k,1/k}$ に有限回しか属さない。従ってそのような $x$ では、十分大きい $k$ に対し
\[
|f_{n_k}(x)|<\frac1k,
\]
ゆえに $f_{n_k}(x)\to0$ である。
\end{enumerate}
\end{proof}

\section{第2期・専門基礎科目（3題必修）}

\begin{problem}{第1問｜等角ベクトルと対称行列}{2021-2-1}
$x_1,\ldots,x_n\in\R^3$ が $|(x_i,x_j)|=r$ $(i\ne j)$，$|(x_i,x_i)|=1$ を満たす。ただし $r\ne1$。
\begin{enumerate}
  \item[(1)] 3次実対称行列全体の次元を求めよ。
  \item[(2)] $\sum_i a_i(x_i,x_j)^2=0$ が全ての $j$ で成り立つなら，全ての $a_i$ が0であることを示せ。
  \item[(3)] 行列 $x_ix_i^{\mathsf T}$ が一次独立であることを示せ。
  \item[(4)] $n\le6$ を示せ。
\end{enumerate}
\end{problem}
\begin{memo*}
\textbf{解答の見通し。}
内積 $(x_i,x_j)=x_i^{\mathsf T}x_j$ は二つのベクトルの角度を測る数である。
目標の $n\le6$ を得るには、$n$ 個の対象を6次元の空間へ入れて一次独立だと示せばよい。
そこで各ベクトルから $3\times3$ 対称行列 $X_i=x_ix_i^{\mathsf T}$ を作る。
行列の一次関係 $\sum_i a_iX_i=O$ に $X_j$ との Frobenius 内積を取ると、
\[
\operatorname{tr}(X_iX_j)=(x_i^{\mathsf T}x_j)^2
\]
となり、問題文 (2) のスカラー方程式へ変わる。
従って (2) から $X_i$ が一次独立だと分かり、その個数は対称行列空間の次元6以下になる。
\end{memo*}
\begin{proof}[解答]
(1) 3次実対称行列は $S^{\mathsf T}=S$、すなわち $s_{ji}=s_{ij}$ を満たす行列である。
従って対角成分3個と、対角より上の成分3個を決めれば残りも決まり、
\[
S=\begin{pmatrix}a&d&e\\d&b&f\\e&f&c\end{pmatrix}
\]
と一意に書ける。この表示は
\[
S=aE_{11}+bE_{22}+cE_{33}
 +d(E_{12}+E_{21})+e(E_{13}+E_{31})+f(E_{23}+E_{32})
\]
と同じであり、右辺の6個の行列は対称行列全体を生成し、
係数が成分としてそのまま現れるので一次独立でもある。従って次元は6である。

(2) 仮定より $(x_j,x_j)^2=|(x_j,x_j)|^2=1$、$i\ne j$ では $(x_i,x_j)^2=|(x_i,x_j)|^2=r^2$
である。従って $j$ を固定した式の $i=j$ の項と $i\ne j$ の項を分けると
\begin{align*}
0
&=\sum_i a_i(x_i,x_j)^2\\
&=a_j+r^2\sum_{i\ne j}a_i\\
&=(1-r^2)a_j+r^2\sum_i a_i. \tag{$*$}
\end{align*}
最後の等号では $\sum_{i\ne j}a_i=\sum_i a_i-a_j$ を使った。
$r=|(x_i,x_j)|\ge0$ と $r\ne1$ より $r^2\ne1$、すなわち $1-r^2\ne0$ である。
式 ($*$) の右辺第2項は $j$ によらないので、$j,k$ に対する二式の差を取れば
\[
(1-r^2)(a_j-a_k)=0,
\]
従って全ての $a_j$ は同じ値 $a$ である。これを ($*$) の前の行へ代入すると
\[
0=a+(n-1)r^2a=\{1+(n-1)r^2\}a.
\]
$n\ge1$, $r\ge0$ なので係数 $1+(n-1)r^2$ は正であり、$a=0$、
従って全ての $a_i$ は0である。

(3) $\sum_i a_ix_ix_i^{\mathsf T}=O$ と仮定する。行列の Frobenius 内積
$\langle C,D\rangle_F=\operatorname{tr}(C^{\mathsf T}D)$ を取る。
$x_ix_i^{\mathsf T}$ は対称なので $C^{\mathsf T}=C$ として扱えて、各 $j$ について
\begin{align*}
0
&=\left\langle\sum_i a_ix_ix_i^{\mathsf T},x_jx_j^{\mathsf T}\right\rangle_F\\
&=\sum_i a_i\operatorname{tr}(x_ix_i^{\mathsf T}x_jx_j^{\mathsf T})\\
&=\sum_i a_i(x_i^{\mathsf T}x_j)^2.
\end{align*}
最後の等式は
\[
x_ix_i^{\mathsf T}x_jx_j^{\mathsf T}
=x_i(x_i^{\mathsf T}x_j)x_j^{\mathsf T},\qquad
\operatorname{tr}(x_ix_j^{\mathsf T})=x_j^{\mathsf T}x_i
\]
を順に用いたものである。
これは (2) の式なので全ての $a_i$ は0、従って $x_ix_i^{\mathsf T}$ は一次独立である。

(4) これらは6次元の実対称行列空間に属する $n$ 個の一次独立な元なので $n\le6$。
\end{proof}

\begin{problem}{第2問｜Gamma 関数と Beta 関数}{2021-2-2}
$p,q>0$ とし
\[
\Gamma(p)=\int_0^\infty x^{p-1}e^{-x}\,\d x,
\quad B(p,q)=\int_0^1x^{p-1}(1-x)^{q-1}\,\d x.
\]
$\Gamma(p)\Gamma(q)=B(p,q)\Gamma(p+q)$ を，二重積分に (1) $x=st,y=s(1-t)$，(2) $x=(r\cos\theta)^2,y=(r\sin\theta)^2$ を用いる二通りで示せ。
\end{problem}
\begin{memo*}
\textbf{解答の見通し。}
まず二つの一変数積分の積を、第1象限 $x>0,y>0$ 上の二重積分として書く。
被積分関数が非負なので、Tonelli の定理により積分順序を変えても値は変わらない。
変数変換では、面積が何倍になるかを表す Jacobian の絶対値も必ず掛ける。

一つ目では $s=x+y$ が二点の合計、$t=x/(x+y)$ が合計に占める $x$ の割合であり、
$s$ だけの積分と $t$ だけの積分へ分離できる。
二つ目では $x,y$ を平方付き極座標で表し、動径 $r$ と角度 $\theta$ の積分へ分ける。
最後に $u=r^2$ と $t=\cos^2\theta$ を代入して、定義された Gamma・Beta 積分へ戻す。
\end{memo*}
\begin{proof}[解答]
積の中の積分変数を区別して $x,y$ と書く。被積分関数は非負なので、Tonelli の定理から
\[
\Gamma(p)\Gamma(q)
=\int_0^\infty\int_0^\infty x^{p-1}y^{q-1}e^{-(x+y)}\,\d x\d y.
\]

(1) $x=st$, $y=s(1-t)$ とおくと、$s=x+y>0$, $t=x/(x+y)\in(0,1)$ であり、
\[
\frac{\partial(x,y)}{\partial(s,t)}
=\begin{vmatrix}t&s\\1-t&-s\end{vmatrix}=-s,
\qquad
\left|\frac{\partial(x,y)}{\partial(s,t)}\right|=s.
\]
従って
\begin{align*}
\Gamma(p)\Gamma(q)
&=\int_0^\infty\int_0^1
(st)^{p-1}\{s(1-t)\}^{q-1}e^{-s}s\,\d t\d s\\
&=\left(\int_0^\infty s^{p+q-1}e^{-s}\,\d s\right)
\left(\int_0^1t^{p-1}(1-t)^{q-1}\,\d t\right)\\
&=\Gamma(p+q)B(p,q).
\end{align*}

(2) $x=(r\cos\theta)^2$, $y=(r\sin\theta)^2$ とおく。
$x,y>0$ から $r>0$ かつ $0<\theta<\pi/2$ であり、逆に
\[
r=\sqrt{x+y},\qquad \theta=\arctan\sqrt{\frac yx}
\]
と一意に定まるので、この対応は $(0,\infty)\times(0,\pi/2)$ と第1象限の間の全単射である。
また、後で使うため
\[
x+y=r^2(\cos^2\theta+\sin^2\theta)=r^2
\]
に注意しておく。各偏微分は
\[
x_r=2r\cos^2\theta,\quad x_\theta=-2r^2\sin\theta\cos\theta,
\]
\[
y_r=2r\sin^2\theta,\quad y_\theta=2r^2\sin\theta\cos\theta
\]
だから、Jacobian は
\begin{align*}
\left|\frac{\partial(x,y)}{\partial(r,\theta)}\right|
&=\left|x_ry_\theta-x_\theta y_r\right|\\
&=\left|4r^3\sin\theta\cos^3\theta
 +4r^3\sin^3\theta\cos\theta\right|\\
&=4r^3\sin\theta\cos\theta
(\cos^2\theta+\sin^2\theta)\\
&=4r^3\sin\theta\cos\theta.
\end{align*}
最後で絶対値が外れるのは $0<\theta<\pi/2$ で $\sin\theta,\cos\theta>0$ だからである。
被積分関数を新しい変数で書き直すと
\begin{align*}
x^{p-1}y^{q-1}e^{-(x+y)}
&=(r^2\cos^2\theta)^{p-1}(r^2\sin^2\theta)^{q-1}e^{-r^2}\\
&=r^{2p-2}\cos^{2p-2}\theta\cdot
  r^{2q-2}\sin^{2q-2}\theta\cdot e^{-r^2}
\end{align*}
なので、Jacobian を掛けると
\begin{align*}
&x^{p-1}y^{q-1}e^{-(x+y)}
 \left|\frac{\partial(x,y)}{\partial(r,\theta)}\right|\\
&\qquad=4r^{(2p-2)+(2q-2)+3}
 \cos^{(2p-2)+1}\theta\,\sin^{(2q-2)+1}\theta\,e^{-r^2}\\
&\qquad=4r^{2p+2q-1}
 \cos^{2p-1}\theta\,\sin^{2q-1}\theta\,e^{-r^2}
\end{align*}
となる。この式は $r$ だけの因子と $\theta$ だけの因子の積なので、積分は分離して
\begin{align*}
\Gamma(p)\Gamma(q)
&=4\int_0^\infty r^{2p+2q-1}e^{-r^2}\,\d r
\int_0^{\pi/2}\cos^{2p-1}\theta\sin^{2q-1}\theta\,\d\theta.
\end{align*}
第1積分では $u=r^2$, $\d u=2r\,\d r$ とおく。
$r^{2p+2q-1}\,\d r=r^{2p+2q-2}\cdot r\,\d r=u^{p+q-1}\cdot\frac{\d u}2$ だから
\[
\int_0^\infty r^{2p+2q-1}e^{-r^2}\,\d r
=\frac12\int_0^\infty u^{p+q-1}e^{-u}\,\d u
=\frac12\Gamma(p+q).
\]
第2積分では $t=\cos^2\theta$, $\d t=-2\sin\theta\cos\theta\,\d\theta$ とおく。
このとき $\sin^2\theta=1-t$ であり、
$\theta=0,\pi/2$ がそれぞれ $t=1,0$ に対応するので
\begin{align*}
\int_0^{\pi/2}\cos^{2p-1}\theta\sin^{2q-1}\theta\,\d\theta
&=\int_0^{\pi/2}\cos^{2p-2}\theta\sin^{2q-2}\theta
  \cdot\sin\theta\cos\theta\,\d\theta\\
&=\int_1^0t^{p-1}(1-t)^{q-1}\left(-\frac{\d t}2\right)\\
&=\frac12\int_0^1t^{p-1}(1-t)^{q-1}\,\d t
=\frac12B(p,q).
\end{align*}
係数も含めて $4\cdot(1/2)\cdot(1/2)=1$ なので、再び
$\Gamma(p)\Gamma(q)=\Gamma(p+q)B(p,q)$ を得る。
\end{proof}

\begin{problem}{第3問｜距離空間の連続写像}{2021-2-3}
距離空間 $(X,d_X)$ から $(Y,d_Y)$ への写像 $f$ について答えよ。
\begin{enumerate}
  \item[(1)] $f$ が点 $x\in X$ で連続であることを開球を用いて定義せよ。
  \item[(2)] $f$ が連続なら，任意の閉集合 $F\subset Y$ に対し $f^{-1}(F)$ が閉であることを示せ。
  \item[(3)] 任意の $A\subset X$ に対し $f(\overline A)\subset\overline{f(A)}$ を示せ。
\end{enumerate}
\end{problem}
\begin{memo*}
\textbf{解答の見通し。}
(1) の開球 $B_X(x,\delta)$ は、$x$ からの距離が $\delta$ 未満の点全体である。
連続とは、出力側で許す誤差 $\varepsilon$ を先に指定したとき、
入力側の十分小さい誤差 $\delta$ を選べることをいう。量化記号はこの順序で書く。
(2) は $x\notin f^{-1}(F)$ を一つ取り、その周りに逆像の補集合へ含まれる開球を作る。
(3) の $x\in\overline A$ は、$x$ のどんな開球にも $A$ の点があるという意味である。
連続性が与える $\delta$-球からその点を取り、像が任意の $\varepsilon$-球へ入ることを示す。
\end{memo*}
\begin{proof}[解答]
\begin{enumerate}
\item[(1)] $B_X(x,\delta)=\{y\in X\mid d_X(x,y)<\delta\}$ とする。
$f$ が $x$ で連続であるとは
\[
\forall\varepsilon>0\ \exists\delta>0\quad
f(B_X(x,\delta))\subset B_Y(f(x),\varepsilon)
\]
が成り立つこと、すなわち
\[
d_X(x,y)<\delta\quad\Longrightarrow\quad
d_Y(f(x),f(y))<\varepsilon
\]
が成り立つことである。

\item[(2)] $F\subset Y$ を閉集合とすると $Y\setminus F$ は開集合である。
$x\in f^{-1}(Y\setminus F)$ を取ると $f(x)\in Y\setminus F$ である。
補集合は開なので、ある $\varepsilon>0$ に対して
\[
B_Y(f(x),\varepsilon)\subset Y\setminus F
\]
となる。$x$ における連続性から、ある $\delta>0$ が存在して
\[
f(B_X(x,\delta))\subset B_Y(f(x),\varepsilon)
\]
となる。従って $B_X(x,\delta)\subset f^{-1}(Y\setminus F)$ である。
各点 $x$ の周りにこのような開球が取れるので $f^{-1}(Y\setminus F)$ は開であり、
\[
f^{-1}(Y\setminus F)=X\setminus f^{-1}(F)
\]
である。従って $f^{-1}(F)$ の補集合は開であり、$f^{-1}(F)$ は閉集合である。

\item[(3)] $x\in\overline A$ を取り、$f(x)\in\overline{f(A)}$ を示す。
任意の $\varepsilon>0$ に対して、$x$ における連続性より、ある $\delta>0$ が存在して
\[
f(B_X(x,\delta))\subset B_Y(f(x),\varepsilon)
\]
となる。$x\in\overline A$ だから $B_X(x,\delta)\cap A\ne\varnothing$ である。
この共通部分から $a$ を一つ取ると
\[
f(a)\in B_Y(f(x),\varepsilon)\cap f(A)
\]
となる。従って $f(x)$ の任意の開球は $f(A)$ と交わるので
$f(x)\in\overline{f(A)}$ である。以上より
\[
f(\overline A)\subset\overline{f(A)}
\]
を得る。
\end{enumerate}
\end{proof}

\chapter{2017年度（平成29年度）}
\label{ch:2017}

\section{専門基礎科目（3題必修）}

\begin{problem}{第1問｜不変部分空間}{2017-b-1}
$f:\R^3\to\R^3$ を
\[
f(x)=\begin{pmatrix}1&0&0\\1&3&-3\\-2&2&-4\end{pmatrix}x
\]
で定める。$f(W)\subset W$ となる1次元および2次元の部分空間 $W$ を全て求めよ。
\end{problem}
\begin{proof}[解答]
固有値は $1,2,-3$ で相異なる。対応する固有ベクトルを
\[
v_1=(-4,11,6)^{\mathsf T},\quad v_2=(0,3,1)^{\mathsf T},\quad
v_3=(0,1,2)^{\mathsf T}
\]
と取れる。従って1次元不変部分空間は $\langle v_1\rangle,\langle v_2\rangle,\langle v_3\rangle$。2次元不変部分空間は $\langle v_1,v_2\rangle,\langle v_1,v_3\rangle,\langle v_2,v_3\rangle$ であり，これが全てである。
\end{proof}

\begin{problem}{第2問｜近似恒等核}{2017-b-2}
$f$ は $[0,\infty)$ 上の有界連続関数とし
$I_n=\int_0^\infty n f(x)e^{-nx}\,\d x$ とおく。
\begin{enumerate}
  \item[(1)] $I_n$ が有限であることを示せ。
  \item[(2)] 任意の $\delta>0$ に対し $\displaystyle\int_\delta^\infty nf(x)e^{-nx}\,\d x\to0$ を示せ。
  \item[(3)] $I_n\to f(0)$ を示せ。
\end{enumerate}
\end{problem}
\begin{proof}[解答]
$|f|\le M$ とすれば (1) は $\int_0^\infty nMe^{-nx}=M$，(2) は絶対値が $Me^{-n\delta}$ 以下。(3) $\int_0^\infty ne^{-nx}=1$ より
\[
I_n-f(0)=\int_0^\infty n(f(x)-f(0))e^{-nx}\,\d x.
\]
連続性で $[0,\delta]$ の差を小さくし，残りは (2) で抑える。
\end{proof}

\begin{problem}{第3問｜有理数の部分空間}{2017-b-3}
$\Q$ に通常の距離を入れ，$A=\{x\in\Q\mid1\le x<2\}$ とする。$A$ は閉，コンパクト，連結か。
\end{problem}
\begin{proof}[解答]
全て否定。
\begin{itemize}
  \item $2-1/n\in A$ は $\Q$ 内で $2\notin A$ に収束するので閉でない。
  \item 距離空間のコンパクト集合は閉だから，コンパクトでない。
  \item 無理数 $c\in(1,2)$ を取ると $A\cap(-\infty,c)$ と $A\cap(c,\infty)$ が非空な開分離を与えるので連結でない。
\end{itemize}
\end{proof}

\section{専門科目（8題から2題選択）}

\begin{problem}{第1問｜群の交換子}{2017-s-1}
$G$ を有限群，$K\subset H$ を部分群とし
$X=\{x\in G\mid[x,h]\in K\ (\forall h\in H)\}$，$[x,y]=x^{-1}y^{-1}xy$ とおく。
\begin{enumerate}
  \item[(1)] $x\in X$ なら $x^{-1}Kx=K$ を示せ。
  \item[(2)] $X$ が $G$ の部分群であることを示せ。
  \item[(3)] $K\triangleleft H$ かつ $H/K$ が可換であることと $H\subset X$ が同値であることを示せ。
\end{enumerate}
\end{problem}
\begin{proof}[解答]
(1) $[x,k]\in K$ から $x^{-1}kx\in K$。有限性により包含は等号になる。(2) $x,y\in X$ に対し
$[xy,h]=y^{-1}[x,h]y[y,h]\in K$。逆元も同様で，単位元も入る。(3) $H\subset X$ は全ての $h_1,h_2\in H$ で $[h_1,h_2]\in K$ という条件。これは $K$ の正規性と商群の可換性に同値。
\end{proof}

\begin{problem}{第2問｜剰余環と局所化}{2017-s-2}
$R$ を可換環，$a\in R$，$S=\{a^n\mid n\ge0\}$ とする。自然な写像 $\pi:R\to R/(a)$，$\tau:R\to S^{-1}R$ を考える。
\begin{enumerate}
  \item[(1)] $x\in\ker\pi\cap\ker\tau$ なら，ある $n$ で $x^n=0$ となることを示せ。
  \item[(2)] $\tau$ が全射なら $\ker\pi+\ker\tau=R$ を示せ。
\end{enumerate}
\end{problem}
\begin{proof}[解答]
(1) $x=ar$ かつ，ある $m$ で $a^mx=0$。従って $x^{m+1}=a^mr^mx=0$。(2) 全射性により $b/1=1/a$ となる $b\in R$ がある。従ってある $n$ で $a^n(ab-1)=0$，すなわち $ab-1\in\ker\tau$。$ab\in(a)=\ker\pi$ だから $1=ab-(ab-1)$ は二つの核の和に入る。
\end{proof}

\begin{problem}{第3問｜被覆空間と Betti 数}{2017-s-3}
\begin{enumerate}
  \item[(1)] $F:\R\to S^1$，$F(\theta)=(\cos\theta,\sin\theta)$ が普遍被覆であることを示し，被覆変換群を求めよ。
  \item[(2)] 四面体の境界のうち，三つの面からなる複体を $K_1$，残り一面を $K_2$ とする。$K_1\cap K_2,K_1,K_2,K_1\cup K_2$ の Betti 数を求めよ。
\end{enumerate}
\end{problem}
\begin{proof}[解答]
(1) 十分短い円弧の逆像は，$2\pi\Z$ だけ平行移動された区間の非交和になる。$\R$ は単連結なので普遍被覆。被覆変換は $T_k(\theta)=\theta+2\pi k$ で，群は $\Z$。
(2) $K_1\simeq D^2$，$K_2\simeq D^2$，$K_1\cap K_2\simeq S^1$，和は四面体境界 $S^2$。従って
\[
\begin{array}{c|ccc}
&b_0&b_1&b_2\\\hline
K_1\cap K_2&1&1&0\\
K_1&1&0&0\\
K_2&1&0&0\\
K_1\cup K_2&1&0&1
\end{array}
\]
で，高次は0。最後の行は Mayer--Vietoris 完全系列からも直ちに出る。
\end{proof}

\begin{problem}{第4問｜管状曲面}{2017-s-4}
$0<a<1$ とし
\[
p(u,v)=(u,u^2/2,0)+\frac{a\cos v}{\sqrt{u^2+1}}(-u,1,0)+a\sin v(0,0,1).
\]
これが曲面であることを示し，Gauss 曲率・平均曲率を求めよ。また $q(s)=p(u_0,s/a)$ が測地線であることを示せ。
\end{problem}
\begin{proof}[解答]
これは平面曲線 $c(u)=(u,u^2/2,0)$ の半径 $a$ の管状曲面である。中心曲線の曲率を $\kappa=(1+u^2)^{-3/2}$ とすると $a\kappa<1$ なので正則。外向き法線を選ぶと主曲率は
\[
k_1=\frac{\kappa\cos v}{1-a\kappa\cos v},\qquad k_2=-\frac1a.
\]
従って
\[
K=-\frac{\kappa\cos v}{a(1-a\kappa\cos v)},\qquad
H=\frac12\left(\frac{\kappa\cos v}{1-a\kappa\cos v}-\frac1a\right).
\]
$q$ は単位速度の半径 $a$ の円で，加速度は $-1/a$ 倍の曲面法線。接成分が0なので測地線である。
\end{proof}

\begin{problem}{第5問｜扇形上の正則関数}{2017-s-5}
$S=\{z\in\mathbb C\mid0<|z|<r,\ \alpha<\arg z<\beta\}$，$\beta<\alpha+2\pi$ とする。
\begin{enumerate}
  \item[(1)] $f(z)$ は $z\to0$ で極限を持つが $f'(z)$ は持たない例を挙げよ。
  \item[(2)] $0<r'<r$，$\alpha<\alpha'<\beta'<\beta$ で切り取った閉じた小扇形の各点を中心に，$S$ に含まれる円板があることを示せ。
  \item[(3)] $f$ が $S$ 上有界正則なら，境界から $\varepsilon$ だけ離した小扇形 $T_\varepsilon$ 上で $|f'(z)|\le M_\varepsilon/|z|$ となることを示せ。
\end{enumerate}
\end{problem}
\begin{proof}[解答]
(1) $S$ 上で対数の枝を取り $f(z)=\exp(\tfrac12\log z)$ とする。$f(z)\to0$ だが $|f'(z)|=1/(2\sqrt{|z|})\to\infty$。(2) $S$ は開集合なので各点に含まれる小円板がある。(3) 角度と外周から一様に離れているため，ある $c_\varepsilon>0$ があり $B(z,c_\varepsilon|z|)\subset S$。$|f|\le B$ と Cauchy の評価を用いれば $|f'(z)|\le B/(c_\varepsilon|z|)$。
\end{proof}

\begin{problem}{第6問｜線形微分方程式}{2017-s-6}
\begin{enumerate}
  \item[(1)] $y''+2y'+2y=0$，$y(0)=1,y'(0)=0$ を解き，零点を全て求めよ。
  \item[(2)] $y''+2y=\cos(\omega x)$，$y(0)=0,y'(0)=1$ の解が $x>0$ で非有界となる $\omega$ と，その解を求めよ。
\end{enumerate}
\end{problem}
\begin{proof}[解答]
(1)
\[
y=e^{-x}(\cos x+\sin x),\qquad x=-\frac\pi4+k\pi\quad(k\in\Z).
\]
(2) 共振する $\omega=\pm\sqrt2$ のときだけ非有界で，
\[
y(x)=\frac{x+2}{2\sqrt2}\sin(\sqrt2x).
\]
\end{proof}

\begin{problem}{第7問｜パラメータ積分}{2017-s-7}
$f$ は $(0,\infty)$ 上可測で $\int_0^\infty x|f(x)|\,\d x<\infty$ を満たす。次を求めよ。
\[
\lim_{a\downarrow0}\frac1a\int_0^\infty(\log(1+ax)-ax)f(x)\,\d x.
\]
\end{problem}
\begin{proof}[解答]
各 $x$ で $(\log(1+ax)-ax)/a\to0$。また $0\le ax-\log(1+ax)\le ax$ なので絶対値は $x$ 以下。$x|f(x)|$ を優関数として優収束定理を使い，答えは $0$。
\end{proof}

\begin{problem}{第8問｜確率収束}{2017-s-8}
実確率変数列 $X_n$ と実確率変数 $X$ について答えよ。
\begin{enumerate}
  \item[(1)] $\lim_{N\to\infty}P(|X|>N)=0$ を示せ。
  \item[(2)] $X_n\to X$ が確率収束し，$f:\R\to\R$ が連続なら $f(X_n)\to f(X)$ も確率収束することを示せ。
\end{enumerate}
\end{problem}
\begin{proof}[解答]
(1) 事象 $\{|X|>N\}$ は単調減少し，共通部分は空集合。確率の上からの連続性を使う。(2) (1) から $P(|X|>M)$ を小さくする。コンパクト区間 $[-M-1,M+1]$ 上で $f$ は一様連続なので，ある $\delta>0$ で $|x-y|<\delta$ なら $|f(x)-f(y)|<\varepsilon$。従って
\[
P(|f(X_n)-f(X)|\ge\varepsilon)
\le P(|X|>M)+P(|X_n-X|\ge\delta),
\]
右辺を順に小さくできる。
\end{proof}

\chapter{2015年度（平成27年度）}
\label{ch:2015}

\section{第1期・専門基礎科目（3題必修）}

\begin{problem}{第1問｜基底と上三角化}{2015-b-1}
$V$ を有限次元実ベクトル空間とする。
\begin{enumerate}
  \item[(1)] 次の定義を述べよ。
  \begin{enumerate}
    \item[(i)] $\{w_1,w_2,\ldots,w_m\}\subset V$ が一次独立であること。
    \item[(ii)] $\{v_1,v_2,\ldots,v_n\}\subset V$ が $V$ の基底であること。
  \end{enumerate}
  \item[(2)] $\{v_1,\ldots,v_n\}$ を $V$ の基底とし，$f:V\to V$ を実ベクトル空間としての同型写像とする。このとき $\{f(v_1),\ldots,f(v_n)\}$ も $V$ の基底であることを示せ。
  \item[(3)] $V=\R^n$ として，$\{v_1,\ldots,v_n\}$ を $V$ の基底とする。各 $i$ $(1\le i\le n)$ に対して $W_i$ を $v_1,\ldots,v_i$ で生成される部分空間とする。$n$ 次正方行列 $A$ が与えられ，各 $i$ に対して $AW_i\subset W_i$ が成り立っているとする。このとき，$n$ 次正則行列 $P$ と $n$ 次上三角行列 $T$ が存在して $P^{-1}AP=T$ とできることを示せ。
\end{enumerate}
\end{problem}
\begin{memo*}
\textbf{解答の見通し。}
一次独立とは、線形結合が0になる方法が係数を全て0にする場合しかないこと、
基底とは、一次独立で空間全体を生成するベクトルの組である。
同型写像は線形な全単射である。線形性で一次結合を $f$ の内側へ入れ、
一次独立性には単射性を、生成することには全射性を使う。

上三角化では、基底ベクトルを列に並べた $P=(v_1\ \cdots\ v_n)$ を作る。
$P^{-1}AP$ は、線形変換 $A$ を基底 $v_1,\ldots,v_n$ の座標で表した行列である。
$Av_i\in\langle v_1,\ldots,v_i\rangle$ を係数表示すると、その係数が
$P^{-1}AP$ の第 $i$ 列になり、第 $i$ 成分より下が全て0になる。
\end{memo*}
\begin{proof}[解答]
\begin{enumerate}
  \item[(1)] (i) $w_1,\ldots,w_m$ が一次独立であるとは
  \[
  \sum_{i=1}^m a_iw_i=0\quad\Longrightarrow\quad a_1=\cdots=a_m=0
  \]
  が任意の $a_i\in\R$ に対して成り立つことをいう。

  (ii) $v_1,\ldots,v_n$ が $V$ の基底であるとは、一次独立であり、かつ
  \[
  V=\langle v_1,\ldots,v_n\rangle
  \]
  となることをいう。

  \item[(2)] $\sum_i a_if(v_i)=0$ なら、線形性と単射性から
  $f(\sum_i a_iv_i)=f(0)$、従って $\sum_i a_iv_i=0$ なので全ての $a_i$ は0。
  また任意の $v\in V$ に対し、全射性から $v=f(w)$ となる $w\in V$ があり、
  $w=\sum_i a_iv_i$ と書けば $v=\sum_i a_if(v_i)$。従って $\{f(v_i)\}$ は基底である。

  \item[(3)] $P=(v_1\ \cdots\ v_n)$ とおく。$\{v_i\}$ は基底なので $P$ は正則である。
  $Av_i\in W_i$ より
  \[
  Av_i=\sum_{j=1}^i t_{ji}v_j.
  \]
  $T=(t_{ji})$ とおけば $T$ は上三角であり、
  \[
  AP=(Av_1\ \cdots\ Av_n)=PT.
  \]
  従って $P^{-1}AP=T$。
\end{enumerate}
\end{proof}

\begin{problem}{第2問｜広義積分}{2015-b-2}
$a$ を実数とする。区間 $(0,\infty)$ 上の関数 $f_a(x)$ を
\[
f_a(x)=x^a\log x
\]
によって定める。
\begin{enumerate}
  \item[(1)] 関数 $f_a(x)$ の不定積分を求めよ。
  \item[(2)] 正数 $R>0$ に対して
  \[
  I(a,R)=\iint_{B(R)}f_a(1+x^2+y^2)\,\d x\d y
  \]
  とおく。ただし $B(R)$ は $xy$ 平面上の原点を中心とする半径 $R$ の円の内部とする。このとき $I(0,1)$ の値を求めよ。
  \item[(3)] (2) で定めた $I(a,R)$ に対して $\displaystyle\lim_{R\to\infty}I(a,R)$ が収束するような $a$ の範囲を求めよ。
\end{enumerate}
\end{problem}
\begin{memo*}
\textbf{解答の見通し。}
不定積分では $\log x$ を微分すると $1/x$ になり、$x^a$ は積分しやすいので部分積分を使う。
$\int x^a\,\d x=x^{a+1}/(a+1)$ には $a\ne-1$ が必要なため、$a=-1$ だけ別に計算する。

二重積分の領域は原点中心の円板で、被積分関数も $x^2+y^2$ だけで決まる。
そこで極座標を使うと面積要素は $\d x\d y=r\,\d r\d\theta$ となり、
角度積分は係数 $2\pi$ になる。$u=1+r^2$ と置けば
不定積分と同じ $\int u^a\log u\,\d u$ へ戻る。
最後は $t=\log u$ と置換する。$\log$ が $t$ に化けて $\int_0^\infty te^{(a+1)t}\,\d t$ になるので、
指数が減衰するかどうか、すなわち $a+1$ の符号だけを見ればよい。
\end{memo*}
\begin{proof}[解答]
\begin{enumerate}
  \item[(1)]
  \[
  \int x^a\log x\,\d x=
  \begin{cases}
  \displaystyle
  x^{a+1}\left(\frac{\log x}{a+1}-\frac1{(a+1)^2}\right)+C &(a\ne-1),\\[2mm]
  \displaystyle\frac12(\log x)^2+C &(a=-1).
  \end{cases}
  \]

  \item[(2)] 極座標と $u=1+r^2$ により
  \[
  I(a,R)=2\pi\int_0^R(1+r^2)^a\log(1+r^2)r\,\d r
  =\pi\int_1^{1+R^2}u^a\log u\,\d u.
  \]
  従って
  \[
  I(0,1)=\pi[u\log u-u]_1^2=\pi(2\log2-1).
  \]

  \item[(3)] $t=\log u$ とおくと
  \[
  \int_1^\infty u^a\log u\,\d u
  =\int_0^\infty te^{(a+1)t}\,\d t.
  \]
  これは $a+1<0$ のとき、かつそのときに限り収束する。従って
  \[
  a<-1.
  \]
  なお、このとき極限値は $\pi/(a+1)^2$ である。
\end{enumerate}
\end{proof}

\begin{problem}{第3問｜二重周期関数}{2015-b-3}
連続写像 $f:\R^2\to\R$ が，任意の $(x,y)\in\R^2$ について
$f(x+1,y)=f(x,y+1)=f(x,y)$ を満たすとする。このとき，以下の問いに答えよ。
\begin{enumerate}
  \item[(1)] $f(\R^2)=f([0,1]\times[0,1])$ を示せ。
  \item[(2)] 写像 $f$ は最大値と最小値を持つことを示せ。
  \item[(3)] 写像 $f$ の最大値を $b$，最小値を $a$ とする。このとき，連続写像 $\sigma:[0,1]\to\R^2$ で，合成 $f\circ\sigma:[0,1]\to[a,b]$ が全射となるようなものが存在することを示せ。
\end{enumerate}
\end{problem}
\begin{memo*}
\textbf{解答の見通し。}
$\lfloor x\rfloor$ は $x$ 以下の最大の整数であり、$0\le x-\lfloor x\rfloor<1$ を満たす。
各座標から整数部分を引けば、周期性により関数値を変えずに任意の点を
基本領域 $[0,1]^2$ へ移せる。
$[0,1]^2$ は閉かつ有界なのでコンパクトであり、連続関数はそこで最大値・最小値を取る。

最小値を取る点 $p$ と最大値を取る点 $q$ を線分 $c(t)=(1-t)p+tq$ で結ぶ。
この線分は正方形の中に残り、一変数連続関数 $f(c(t))$ は端点で最小値と最大値を取る。
従って中間値の定理により、その間の全ての値を通る。
\end{memo*}
\begin{proof}[解答]
\begin{enumerate}
  \item[(1)] $m=\lfloor x\rfloor$, $n=\lfloor y\rfloor$ とすると、周期性より
  \[
  f(x,y)=f(x-m,y-n),\qquad (x-m,y-n)\in[0,1)^2.
  \]
  従って $f(\R^2)\subset f([0,1]^2)$。逆の包含は明らかなので等号が成り立つ。

  \item[(2)] $[0,1]^2$ はコンパクトだから、$f$ はそこで最大値と最小値を取る。
  (1) より、それらは $\R^2$ 全体での最大値・最小値でもある。

  \item[(3)] $f(p)=a$, $f(q)=b$ となる $p,q\in[0,1]^2$ を取り
  \[
  \sigma(t)=(1-t)p+tq\qquad(0\le t\le1)
  \]
  とおく。$f\circ\sigma$ は連続で端点の値が $a,b$ だから、
  中間値の定理より $[a,b]$ への全射である。
\end{enumerate}
\end{proof}

\section{専門科目（8題から2題選択）}

\begin{problem}{第1問｜有限体上の行列群}{2015-s-1}
$\mathbb F$ を素数 $p$ を位数とする有限体とする。$GL_n(\mathbb F)$ を $\mathbb F$ の元を成分とする
$n$ 次正則行列全体とし，
\[
G=\{A=(a_{ij})\in GL_3(\mathbb F)\mid\det A=1,\ a_{12}=0,\ a_{22}=1,\ a_{32}=0\}
\]
とする。
\begin{enumerate}
  \item[(1)] $G$ が行列の積に関して群をなすことを示せ。
  \item[(2)] 写像 $(a_{ij})\mapsto\begin{pmatrix}a_{11}&a_{13}\\a_{31}&a_{33}\end{pmatrix}$ により群準同型 $\varphi:G\to GL_2(\mathbb F)$ が得られることを示せ。
  \item[(3)] (2) の写像 $\varphi$ の核を $K$ とする。$G$ における $K$ の中心化群 $C_G(K)$ を求めよ。
  \item[(4)] $X=\begin{pmatrix}1&0&0\\0&1&0\\1&0&1\end{pmatrix}$ とし，$K$ と $X$ で生成される部分群を $H$ とする。$H$ の位数を求めよ。また，$H$ の中心 $Z(H)$ を求めよ。
\end{enumerate}
\end{problem}
\begin{memo*}
\textbf{解答の見通し。}
基底順を $(e_1,e_3;e_2)$ に並べ替えると、元は
線形部分 $B\in SL_2(\mathbb F_p)$ と平行移動部分 $u\in\mathbb F_p^2$ の組 $(B,u)$ になる。
この積公式で群・準同型・核をまとめて扱う。
中心化群では全ての $u$ を固定する $B$ を求め、
$H$ は $K$ と位数 $p$ の元による半直積として数える。
\end{memo*}
\begin{proof}[解答]
(1) 基底順を $(e_1,e_3;e_2)$ と見れば元は
$\begin{pmatrix}B&0\\u&1\end{pmatrix}$（$B\in SL_2,u\in\mathbb F_p^2$）と書ける。
この表示で積と逆元は
\[
\begin{pmatrix}B&0\\u&1\end{pmatrix}
\begin{pmatrix}C&0\\v&1\end{pmatrix}
=\begin{pmatrix}BC&0\\uC+v&1\end{pmatrix},
\]
\[
\begin{pmatrix}B&0\\u&1\end{pmatrix}^{-1}
=\begin{pmatrix}B^{-1}&0\\-uB^{-1}&1\end{pmatrix}.
\]
$BC,B^{-1}\in SL_2(\mathbb F_p)$ なので $G$ は群である。

(2) $\varphi(B,u)=B$ について $\varphi((B,u)(C,v))=BC=\varphi(B,u)\varphi(C,v)$ だから、$\varphi$ は群準同型である。

(3) $K=\ker\varphi=\{(I,u)\}\cong\mathbb F_p^2$ である。$(B,w)$ が全ての $(I,u)\in K$ と可換する条件を積の公式へ代入すると
\[
(B,w)(I,u)=(B,w+u),\qquad
(I,u)(B,w)=(B,uB+w).
\]
従って全ての行ベクトル $u$ について $uB=u$、すなわち $B=I$ が必要十分である。
よって $C_G(K)=K$。

(4) $B_0=\varphi(X)=\begin{pmatrix}1&0\\1&1\end{pmatrix}$ とおくと
\[
B_0^j=\begin{pmatrix}1&0\\j&1\end{pmatrix}\qquad(j\in\mathbb F_p)
\]
なので $B_0$ の位数は $p$。$K\cap\langle X\rangle=\{1\}$ で、$X$ は共役により
$K$ を保つから、$H=K\rtimes\langle X\rangle$ であり
$|H|=p^2p=p^3$ となる。

中心の元は特に $K$ と可換するので、上の議論から $(I,u)$ の形に限る。
$u=(c,d)$ とすると
\[
(I,u)(B_0,0)=(B_0,uB_0),\qquad
(B_0,0)(I,u)=(B_0,u),
\]
かつ $uB_0=(c+d,d)$ だから、$X$ と可換する条件は $d=0$ である。従って中心は
\[
Z(H)=\left\{\begin{pmatrix}1&0&0\\c&1&0\\0&0&1\end{pmatrix}\middle|c\in\mathbb F_p\right\},
\]
従って位数 $p$。
\end{proof}

\begin{problem}{第2問｜整数係数多項式環}{2015-s-2}
整数係数1変数多項式環 $\Z[X]$ から実数環の直積環 $\R\times\R$ への環準同型写像 $\Phi$ を，
$f(X)\in\Z[X]$ に対し
\[
\Phi(f(X))=\bigl(f(\sqrt7),\,f(\sqrt{13})\bigr)
\]
で定める。
\begin{enumerate}
  \item[(1)] $f(X)\in\Z[X]$ に対し，$f(\sqrt7)=0$ となるなら，$f(X)$ は $X^2-7$ で割り切れることを示せ。
  \item[(2)] $\ker\Phi=((X^2-7)(X^2-13))$ を示せ。
  \item[(3)] $\Phi(X^2-7)$ で生成される $\operatorname{Im}\Phi$ のイデアルを $P_1$，$\Phi(X^2-13)$ で生成される $\operatorname{Im}\Phi$ のイデアルを $P_2$ とする。$P_1$ および $P_2$ は $\operatorname{Im}\Phi$ の素イデアルであるが極大イデアルではないことを示せ。
  \item[(4)] $\operatorname{Im}\Phi$ の素イデアルで，$P_1$ と $P_2$ の両方を含むものの個数を求めよ。
\end{enumerate}
\end{problem}
\begin{memo*}
\textbf{解答の見通し。}
$X^2-7$ で割った余りは一次式なので、$\sqrt7$ の無理性から余りが0と分かる。
この議論を $\sqrt{13}$ にも順に使って核を決める。
$P_1,P_2$ の素性は各成分への射影の商が整域になることから示す。
両方を含む素イデアルは差6も含むため、標数2と3に分けて因数分解を数える。
\end{memo*}
\begin{proof}[解答]
(1) 多項式の除法により
\[
f(X)=(X^2-7)q(X)+rX+s\qquad(r,s\in\Z)
\]
と書ける。$f(\sqrt7)=0$ なら $r\sqrt7+s=0$。もし $r\ne0$ なら
$\sqrt7=-s/r\in\Q$ となって矛盾するので $r=0$、さらに $s=0$ である。
従って $X^2-7$ は $f$ を割る。

(2) $f\in\ker\Phi$ とする。(1) から $f=(X^2-7)q$ と書ける。
$0=f(\sqrt{13})=6q(\sqrt{13})$ なので $q(\sqrt{13})=0$。
同じ議論により $X^2-13$ は $q$ を割るから
\[
f\in((X^2-7)(X^2-13)).
\]
逆にこの積の倍数は $\sqrt7$, $\sqrt{13}$ の両方で0になるので、
\[
\ker\Phi=((X^2-7)(X^2-13)).
\]

(3) $\operatorname{Im}\Phi$ から第1成分を取る写像の像は $\Z[\sqrt7]$ で、核は
$P_1$ である。従って準同型定理から
\[
\operatorname{Im}\Phi/P_1\cong\Z[\sqrt7].
\]
これは整域なので $P_1$ は素イデアルだが、例えば2は逆元を持たないので体ではなく、
$P_1$ は極大でない。第2成分への射影
$\Z[\sqrt7]\times\Z[\sqrt{13}]\to\Z[\sqrt{13}]$
を $\Phi$ と合成した写像の核は $P_2$、像は $\Z[\sqrt{13}]$ である。従って
$\operatorname{Im}\Phi/P_2\cong\Z[\sqrt{13}]$ であり、$P_2$ も素だが極大でない。

(4) $P_1,P_2$ の両方を含む素イデアル $\mathfrak p$ は
\[
\Phi(X^2-7)-\Phi(X^2-13)=\Phi(6)
\]
を含む。従って $6\in\mathfrak p$ であり、素性から $2\in\mathfrak p$ または
$3\in\mathfrak p$ である。法2では二つの関係式はいずれも
\[
X^2-1=(X-1)^2
\]
となるので素イデアルは $(2,X-1)$ に対応する1個。法3では
\[
X^2-1=(X-1)(X+1)
\]
となるので $(3,X-1)$ と $(3,X+1)$ に対応する2個である。
従って求める素イデアルは合計 $1+2=3$ 個である。
\end{proof}

\begin{problem}{第3問｜リフトと基本群}{2015-s-3}
$X,Y$ を弧状連結な位相空間，$x_0$ を $X$ の点とする。連続写像 $f:Y\to X$ は次の性質を満たすとする。
\begin{itemize}
  \item $F=f^{-1}(\{x_0\})$ とおく。$F$ の相対位相は離散位相である。
  \item $A$ を $[0,1]$（または $[0,1]\times[0,1]$）とする。このとき $g:A\to X$ と $f(b)=g(0)$（または $g(0,0)$）となる $b\in Y$ に対して，連続写像 $\tilde g:A\to Y$ で $\tilde g(0)=b$（または $\tilde g(0,0)=b$）かつ $f\circ\tilde g=g$ となるものが唯一つ存在する。この $\tilde g$ を $g$ のリフトと呼ぶ。
\end{itemize}
以下の問いに答えよ。
\begin{enumerate}
  \item[(1)] $f_*:\pi_1(Y,y_0)\to\pi_1(X,x_0)$ は単射であることを示せ。ただし，$y_0\in F$ であり，$\pi_1(X,x_0),\pi_1(Y,y_0)$ はそれぞれの位相空間の基本群（1次ホモトピー群）であり，$f_*:\pi_1(Y,y_0)\to\pi_1(X,x_0)$ は $f_*([\delta])=[f\circ\delta]$ によって定められる基本群の間の準同型写像である。
  \item[(2)] $\varphi:\pi_1(X,x_0)\to F$ を $\varphi([\gamma])=\tilde\gamma(1)$ で定めるとき，以下に答えよ。ただし $\tilde\gamma$ は $\tilde\gamma(0)=y_0$ を満たす $\gamma$ のリフトである。\\
  (i) $\varphi$ の定義はホモトピー類の代表元の取り方によらないことを示せ。\\
  (ii) $\varphi$ は全射であることを示せ。\\
  (iii) $\varphi([\gamma_0])=\varphi([\gamma_1])$ が成り立つことと $[\gamma_0][\gamma_1]^{-1}\in f_*(\pi_1(Y,y_0))$ は同値であることを示せ。
\end{enumerate}
\end{problem}
\begin{memo*}
\textbf{解答の見通し。}
ホモトピーを持ち上げると、その端点はファイバー $F$ の中を連続に動く。
$F$ は離散、パラメータ区間は連結なので端点は動けない、という事実を繰り返し使う。
$f_*$ の単射性、$\varphi$ の代表元によらないこと、二つのリフトの終点条件を
全てこの一つの原理から示す。
\end{memo*}
\begin{proof}[解答]
(1) $y_0$ を基点とする閉道 $\eta:[0,1]\to Y$ に対して
\[
f_*([\eta])=[f\circ\eta]=1
\]
と仮定する。このとき $f\circ\eta$ と定値道 $c_{x_0}$ の間に、端点を固定するホモトピー
\[
H:[0,1]\times[0,1]\to X
\]
が存在する。$H(0,t)=f(\eta(t))$ とし、このホモトピーを初期リフト $\eta$ から持ち上げて
$\widetilde H:[0,1]^2\to Y$ を得る。

各 $s\in[0,1]$ について $H(s,0)=H(s,1)=x_0$ だから
\[
\widetilde H(s,0),\ \widetilde H(s,1)\in F=f^{-1}(x_0).
\]
二つの写像 $s\mapsto\widetilde H(s,0)$ と $s\mapsto\widetilde H(s,1)$ は連続であり、
$[0,1]$ は連結、$F$ は離散なので、それぞれ一定である。$s=0$ では両端点が $y_0$ だから
\[
\widetilde H(s,0)=\widetilde H(s,1)=y_0
\qquad(0\le s\le1).
\]
さらに $H(1,t)=x_0$ であり、その $y_0$ から始まるリフトは道のリフトの一意性により
定値道 $c_{y_0}$ である。従って $\eta=\widetilde H(0,\cdot)$ は
$c_{y_0}=\widetilde H(1,\cdot)$ と端点を固定してホモトピックであり、$[\eta]=1$ である。
よって $\ker f_*=\{1\}$、すなわち $f_*$ は単射である。

(2)(i) $x_0$ を基点とする閉道 $\gamma$ の $y_0$ から始まる一意なリフトを
$\widetilde\gamma$ とし
\[
\varphi([\gamma])=\widetilde\gamma(1)\in F
\]
と定める。同じ類を表す閉道 $\gamma_0,\gamma_1$ の間の端点固定ホモトピーを持ち上げると、
リフトの終点は $F$ の中を連続に動く。$F$ は離散なので終点は一定であり
\[
\widetilde\gamma_0(1)=\widetilde\gamma_1(1).
\]
従って $\varphi$ は代表元の取り方によらず well-defined である。

(ii) 任意の $y\in F$ に対して、$Y$ の弧状連結性より
$y_0$ から $y$ への道 $\alpha$ を取れる。このとき
\[
\gamma=f\circ\alpha
\]
は $x_0$ を基点とする閉道であり、$\alpha$ 自身が $\gamma$ の $y_0$ から始まるリフトである。
従って $\varphi([\gamma])=\alpha(1)=y$ となり、$\varphi$ は全射である。

(iii) $\gamma_0,\gamma_1$ の $y_0$ から始まるリフトを
$\widetilde\gamma_0,\widetilde\gamma_1$ とする。もし
\[
\varphi([\gamma_0])=\varphi([\gamma_1])=:y
\]
なら、道
\[
\widetilde\gamma_0*\widetilde\gamma_1^{-1}
\]
は $y_0$ を基点とする $Y$ の閉道であり、その $f$ による像は
$\gamma_0*\gamma_1^{-1}$ である。従って
\[
[\gamma_0][\gamma_1]^{-1}
=f_*([\widetilde\gamma_0*\widetilde\gamma_1^{-1}])
\in f_*(\pi_1(Y,y_0)).
\]

逆に $[\gamma_0][\gamma_1]^{-1}\in f_*(\pi_1(Y,y_0))$ とする。
このとき $\gamma_0*\gamma_1^{-1}$ は、ある $Y$ の閉道の像と端点を固定してホモトピックである。
ホモトピーを持ち上げ、上と同じく終点が離散集合 $F$ の中で一定であることを使うと、
$\gamma_0*\gamma_1^{-1}$ の $y_0$ から始まるリフトも $y_0$ で終わる。
このリフトの後半を逆向きに見ると、それは $\gamma_1$ の $y_0$ から始まるリフトであり、
その終点は $\widetilde\gamma_0(1)$ である。道のリフトの一意性より、このリフトは
$\widetilde\gamma_1$ と一致する。従って
\[
\widetilde\gamma_0(1)=\widetilde\gamma_1(1),
\]
すなわち $\varphi([\gamma_0])=\varphi([\gamma_1])$ である。
\end{proof}

\begin{problem}{第4問｜楕円放物面}{2015-s-4}
$\R^2$ 上の関数 $f$ を $f(x,y)=2x^2+y^2$ で定める。$f$ のグラフを $S$ で表し，
$S$ 内の曲線 $C$ を $C=\{(t,t,f(t,t))\mid t\in\R\}$ で定義する。
\begin{enumerate}
  \item[(1)] $S$ の各点でのガウス曲率および平均曲率を求めよ。
  \item[(2)] $C$ の接線に関する $S$ の法曲率を求めよ。
  \item[(3)] (2) で求めた法曲率が $S$ の平均曲率に等しいような $C$ の点を全て求めよ。
\end{enumerate}
\end{problem}
\begin{memo*}
\textbf{解答の見通し。}
グラフ表示から第1・第2基本形式の係数を順に計算する。
Gauss 曲率と平均曲率はこれらの標準公式へ代入する。
曲線 $C(t)=X(t,t)$ の座標速度は $(1,1)$ なので、
法曲率は第2基本形式と第1基本形式の比で求まる。
最後は $x=y=t$ を代入し、分母が正であることを確認して方程式を整理する。
\end{memo*}
\begin{proof}[解答]
(1) グラフを $X(x,y)=(x,y,2x^2+y^2)$ とパラメータ表示する。
\[
X_x=(1,0,4x),\qquad X_y=(0,1,2y)
\]
なので、第1基本形式の係数と上向き単位法線は
\[
E=1+16x^2,\quad F=8xy,\quad G=1+4y^2,
\]
\[
n=\frac{(-4x,-2y,1)}W,\qquad W=\sqrt{1+16x^2+4y^2}.
\]
また $X_{xx}=(0,0,4)$, $X_{xy}=0$, $X_{yy}=(0,0,2)$ より
第2基本形式の係数は
\[
e=\frac4W,\qquad f=0,\qquad g=\frac2W.
\]
従って
\begin{align*}
K&=\frac{eg-f^2}{EG-F^2}
=\frac{8/W^2}{W^2}=\frac8{W^4},\\
H&=\frac{eG-2fF+gE}{2(EG-F^2)}\\
&=\frac{4(1+4y^2)+2(1+16x^2)}{2W^3}
=\frac{3+16x^2+8y^2}{W^3}.
\end{align*}
すなわち
\[
K=\frac8{W^4},\qquad H=\frac{3+16x^2+8y^2}{W^3}.
\]
(2) $C(t)=X(t,t)$ の座標方向は $(\d x,\d y)=(1,1)$ なので、法曲率は
第2基本形式と第1基本形式の比から
\begin{align*}
k_n(t)
&=\frac{e+2f+g}{E+2F+G}\bigg|_{(x,y)=(t,t)}\\
&=\frac{6/\sqrt{1+20t^2}}{2+36t^2}\\
&=\frac6{\sqrt{1+20t^2}(2+36t^2)}.
\end{align*}
従って
\[
k_n(t)=\frac6{\sqrt{1+20t^2}(2+36t^2)}.
\]
(3) 平均曲率に $x=y=t$ を代入すると
\[
H(t,t)=\frac{3+24t^2}{(1+20t^2)^{3/2}},\qquad
k_n(t)=\frac6{\sqrt{1+20t^2}(2+36t^2)}.
\]
両辺に正の $(1+20t^2)^{3/2}(2+36t^2)$ を掛けて整理すると
\begin{align*}
6(1+20t^2)&=(3+24t^2)(2+36t^2)\\
&=6+156t^2+864t^4,
\end{align*}
すなわち $36t^2(1+24t^2)=0$。実数では $t=0$ だけなので、点は $(0,0,0)$ のみ。
\end{proof}

\begin{problem}{第5問｜3次方程式}{2015-s-5}
複素変数の3次関数 $f(z)=z^3+az+b$ について，次の問いに答えよ。
\begin{enumerate}
  \item[(1)] $R:=\max\{|a|,|b|\}+1$ とおく。このとき，$|z|\ge R+1$ を満たす $z$ に対して，$f(z)\ne0$ を示せ。
  \item[(2)] 
  \[
  P_n:=P_n(a,b)=\frac1{2\pi i}\int_{|z|=R+1}\frac{z^n(3z^2+a)}{z^3+az+b}\,\d z
  \qquad(n\ge0)
  \]
  とおく。このとき，$\Delta:=P_0P_3^2-\dfrac12P_2^3$ を $a,b$ を用いて表せ。
  \item[(3)] $\Delta=0$ かつ $(a,b)\ne(0,0)$ のとき，3次方程式 $z^3+az+b=0$ の解を全て求めよ。
\end{enumerate}
\end{problem}
\begin{memo*}
\textbf{解答の見通し。}
大円上では最高次項 $z^3$ が $az+b$ より大きいことを示し、全ての根を円内に入れる。
$f'/f$ の積分は偏角の原理により根の冪和 $P_n=\sum\alpha_j^n$ になる。
Vieta の公式で $P_2,P_3$ を $a,b$ に直し、
$\Delta=0$ では重根候補 $r=-3b/(2a)$ を作って多項式を因数分解する。
\end{memo*}
\begin{proof}[解答]
(1) $M=\max\{|a|,|b|\}$ とおくと $R=M+1$ である。$t=|z|\ge R+1=M+2$ なら
\[
|az+b|\le |a|t+|b|\le M(t+1)\le(t-2)(t+1)<t^3=|z^3|.
\]
従って逆三角不等式から
$|f(z)|\ge |z|^3-|az+b|>0$ となり、$f(z)\ne0$ である。

(2) $f'(z)=3z^2+a$ なので、偏角原理（対数微分の留数）から
\[
P_n=\frac1{2\pi i}\int_Cz^n\frac{f'(z)}{f(z)}\,\d z
=\sum_{j=1}^3\alpha_j^n,
\]
ここで $\alpha_1,\alpha_2,\alpha_3$ は重複度を込めた3根である。
$P_0=3$ であり、$\alpha_1+\alpha_2+\alpha_3=0$、
$\alpha_1\alpha_2+\alpha_2\alpha_3+\alpha_3\alpha_1=a$、
$\alpha_1\alpha_2\alpha_3=-b$ を使うと
\begin{align*}
P_2&=(\alpha_1+\alpha_2+\alpha_3)^2
-2(\alpha_1\alpha_2+\alpha_2\alpha_3+\alpha_3\alpha_1)=-2a,\\
P_3&=(\alpha_1+\alpha_2+\alpha_3)^3
-3(\alpha_1+\alpha_2+\alpha_3)
(\alpha_1\alpha_2+\alpha_2\alpha_3+\alpha_3\alpha_1)
+3\alpha_1\alpha_2\alpha_3=-3b.
\end{align*}
従って
\[
\Delta=3(-3b)^2-\frac12(-2a)^3=4a^3+27b^2.
\]

(3) $\Delta=0$ かつ $(a,b)\ne(0,0)$ なら $a\ne0$ である。
$r=-3b/(2a)$ とおく。$4a^3+27b^2=0$ を用いると
\[
-3r^2=a,\qquad 2r^3=b.
\]
従って係数を比較して
\[
(z-r)^2(z+2r)=z^3-3r^2z+2r^3=z^3+az+b=f(z).
\]
よって根は $r=-3b/(2a)$（重複度2）と $-2r=3b/a$（重複度1）である。
\end{proof}

\begin{problem}{第6問｜常微分方程式}{2015-s-6}
\begin{enumerate}
  \item[(1)] 微分方程式の初期値問題
  \[
  \begin{cases}
  \dfrac{\d y}{\d x}=\sin y,\\[1mm]
  y(0)=\dfrac\pi2
  \end{cases}
  \]
  の解を求めよ。
  \item[(2)] $f(x,y)$ を $\R^2$ 上の $C^1$ 級関数とする。微分方程式
  \[
  \frac{\d y}{\d x}=\sin f(x,y)
  \]
  の任意の解の定義域は $\R$ 全体となることを示せ。
\end{enumerate}
\end{problem}
\begin{memo*}
\textbf{解答の見通し。}
初めの方程式は変数分離し、$1/\sin y$ の原始関数を
$\log\tan(y/2)$ として求める。
一般の方程式では右辺の絶対値が1以下なので、解は有限時間に高々その時間幅しか動かない。
従って有限端点で発散できず、局所存在・一意性定理によりさらに延長できる。
\end{memo*}
\begin{proof}[解答]
(1) $0<y<\pi$ の範囲で変数分離すると
\[
\frac{\d y}{\sin y}=\d x,\qquad
\int\frac{\d y}{\sin y}=\log\tan\frac y2.
\]
ここで原始関数は
\[
\frac{\d}{\d y}\log\tan\frac y2
=\frac{1}{\tan(y/2)}\cdot\frac12\sec^2\frac y2
=\frac1{2\sin(y/2)\cos(y/2)}
=\frac1{\sin y}
\]
から確かめられる。
従って $\log\tan(y/2)=x+C$。初期条件 $y(0)=\pi/2$ から $C=0$ なので
\[
\tan\frac y2=e^x,\qquad y(x)=2\arctan(e^x).
\]
(2) 以下、解とは延長できないもの、すなわち極大解を指す。
$y'=\sin f(x,y)$ では $|y'|\le1$ である。解が $x_0$ から有限時刻 $b$ まで存在するとき
\[
|y(x)-y(x_0)|\le|x-x_0|
\]
だから $y$ は有限区間で発散しない。右辺は $C^1$ なので局所存在・一意性定理により
有限な端点の近くからさらに解を延長できる。従って極大解の定義域の端点は有限にならず、定義域は $\R$ 全体である。
\end{proof}

\begin{problem}{第7問｜収束定理}{2015-s-7}
$p>0$ を定数とする。測度空間 $(X,\mathfrak B,\mu)$ 上の $p$ 乗可積分関数 $f$ は
$0<f<\infty$ かつ $\int_Xf^p\,\d\mu=1$ を満たしているとする。このとき次の値を求めよ。
\[
\lim_{n\to\infty}\int_X n\left|\sin\left(\frac{f(x)}n\right)\right|^p\,\d\mu(x).
\]
（収束定理を使う場合はその根拠を述べよ。）
\end{problem}
\begin{memo*}
\textbf{解答の見通し。}
被積分関数の $n$ の効き方を先に取り出す。$u_n=f(x)/n$ とおくと
$|\sin u_n|^p=u_n^p(|\sin u_n|/u_n)^p$ であり、$u_n^p=f(x)^p/n^p$ だから
被積分関数は $n^{1-p}f(x)^p$ に、1へ収束する因子を掛けた形になる。
従って $n^{1-p}$ の挙動、すなわち $p$ と1の大小で結論が三つに分かれる。

$|\sin u|\le|u|$ による上からの評価 $n|\sin u_n|^p\le n^{1-p}f^p$ は、
$p>1$ ならそのまま積分して0を与える。$p=1$ ではこの評価が可積分な優関数 $f$ を与えるので
優収束定理が使える。$p<1$ では被積分関数が各点で $\infty$ へ発散するので、
下から評価する Fatou の補題を使う。
\end{memo*}
\begin{proof}[解答]
各 $x\in X$ について $u_n=f(x)/n$ とおく。$0<f(x)<\infty$ だから $u_n>0$ かつ $u_n\to0$ であり、
$\sin u/u\to1$ $(u\to0)$ より
\[
\frac{|\sin u_n|}{u_n}\longrightarrow1.
\]
$|\sin u_n|^p=u_n^p\left(|\sin u_n|/u_n\right)^p$ と $u_n^p=f(x)^p/n^p$ を使うと
\[
n\left|\sin\left(\frac{f(x)}n\right)\right|^p
=n^{1-p}f(x)^p
\left(\frac{|\sin u_n|}{u_n}\right)^p
\tag{$*$}
\]
であり、右端の因子は各点で1へ収束する。
また $u\ge0$ で $|\sin u|\le u$ だから
\[
0\le n\left|\sin\left(\frac{f(x)}n\right)\right|^p
\le n\,u_n^p=n^{1-p}f(x)^p
\tag{$**$}
\]
が全ての $n$ と $x$ で成り立つ。以下、$p$ と1の大小で場合を分ける。

$p>1$ のとき。式 ($**$) を積分し、仮定 $\int_Xf^p\,\d\mu=1$ を使うと
\[
0\le\int_X n\left|\sin\left(\frac{f(x)}n\right)\right|^p\,\d\mu
\le n^{1-p}\int_Xf^p\,\d\mu=n^{1-p}.
\]
$1-p<0$ だから $n^{1-p}\to0$ であり、はさみうちにより極限は $0$ である。

$p=1$ のとき。式 ($**$) は $n|\sin(f(x)/n)|\le f(x)$ となり、
$\int_Xf\,\d\mu=1<\infty$ だから $f$ は可積分な優関数である。
一方 式 ($*$) で $n^{1-p}=1$ なので、被積分関数は各点で $f(x)$ へ収束する。
よって優収束定理を適用でき
\[
\lim_{n\to\infty}\int_X n\left|\sin\left(\frac{f(x)}n\right)\right|\,\d\mu
=\int_Xf\,\d\mu=1.
\]

$0<p<1$ のとき。$1-p>0$ だから $n^{1-p}\to\infty$ であり、$f(x)^p>0$ と
式 ($*$) から、被積分関数は各点で $\infty$ へ発散する。
被積分関数は非負なので Fatou の補題が使えて
\[
\liminf_{n\to\infty}\int_X n\left|\sin\left(\frac{f(x)}n\right)\right|^p\,\d\mu
\ge\int_X\liminf_{n\to\infty}
n\left|\sin\left(\frac{f(x)}n\right)\right|^p\,\d\mu
=\int_X\infty\,\d\mu.
\]
ここで $\int_Xf^p\,\d\mu=1\ne0$ より $\mu(X)>0$ だから右辺は $\infty$ である。
従って極限は $+\infty$ である。

以上をまとめると、求める極限は
\[
\lim_{n\to\infty}\int_X n\left|\sin\left(\frac{f(x)}n\right)\right|^p\,\d\mu
=\begin{cases}
+\infty,&0<p<1,\\
1,&p=1,\\
0,&p>1
\end{cases}
\]
である。
\end{proof}

\begin{problem}{第8問｜独立確率変数の級数}{2015-s-8}
実数値確率変数 $X$ は連続型で，確率密度関数を $f$ とする。さらに，$X$ は2乗可積分であるとし，
$X$ の平均 $E[X]$ を $\mu$，分散 $V[X]$ を $\sigma^2$ とおく。
実数値確率変数列 $\{X_n\}_{n=1}^\infty$ は独立で，各 $X_n$ の分布関数は
\[
\int_{-\infty}^xnf(ny)\,\d y\qquad(x\in\R)
\]
によって与えられるものとする。整数 $n\ge1$ に対して
\[
S_n=\sum_{k=1}^nX_k
\]
とおくとき，以下の問いに答えよ。
\begin{enumerate}
  \item[(1)] 整数 $n\ge1$ に対して，$X_n$ の平均 $E[X_n]$ および分散 $V[X_n]$ を求めよ。
  \item[(2)] 確率変数列 $\{S_n\}_{n=1}^\infty$ が概収束するための必要十分条件は $\mu=0$ であることを示せ。
  \item[(3)] $X$ が標準正規分布にしたがうとき，(2) より $\{S_n\}_{n=1}^\infty$ はある確率変数 $S$ に概収束することがわかる。$S$ のしたがう確率分布を求めよ。必要があれば，
  \[
  \sum_{n=1}^\infty\frac1{n^2}=\frac{\pi^2}6
  \]
  を用いてもよい。
\end{enumerate}
\end{problem}
\begin{memo*}
\textbf{解答の見通し。}
密度の変数変換から $X_n$ は $X/n$ と同分布だと読み取る。
平均を引いた $Y_n=X_n-\mu/n$ は分散の総和が有限なので、その級数は概収束する。
残る平均部分は調和級数であり、$\mu=0$ かどうかを判定する。
標準正規の場合は独立な正規分布の和なので、分散の極限から極限分布が分かる。
\end{memo*}
\begin{proof}[解答]
(1) $X_n$ の分布関数を変数変換すると、任意の有界可測関数 $\psi$ に対し
\[
E[\psi(X_n)]=\int_{\R}\psi(x)nf(nx)\,\d x
=\int_{\R}\psi(u/n)f(u)\,\d u.
\]
従って $X_n$ は $X/n$ と同分布であり
\[
E[X_n]=\frac\mu n,\qquad V[X_n]=\frac{\sigma^2}{n^2}.
\]

(2) $Y_n=X_n-\mu/n$ とおくと、$Y_n$ は独立、平均0で
\[
\sum_{n=1}^\infty V[Y_n]
=\sigma^2\sum_{n=1}^\infty\frac1{n^2}<\infty.
\]
Kolmogorov の収束定理により $\sum_nY_n$ はほとんど確実に収束する。
従って
\[
S_n=\sum_{k=1}^nY_k+\mu\sum_{k=1}^n\frac1k.
\]
$\mu=0$ なら右辺第1項だけなので $S_n$ は概収束する。
$\mu\ne0$ なら第1項は有限値へ収束する一方、調和級数を含む第2項は
$\operatorname{sgn}(\mu)\infty$ へ発散するため、$S_n$ は有限な確率変数へ概収束しない。

(3) $X$ が標準正規分布なら $X_n\sim N(0,1/n^2)$ であり、独立性から
\[
S_n\sim N\left(0,\sum_{k=1}^n\frac1{k^2}\right).
\]
分散は $\sum_{k=1}^n1/k^2\to\pi^2/6$ だから、概収束極限の分布は
$N(0,\pi^2/6)$ である。
\end{proof}

\section{第2期・専門基礎科目（3題必修）}

\begin{problem}{第1問｜像・核と直和}{2015-2-1}
\begin{enumerate}
  \item[(1)] 
  \[
  A=\begin{pmatrix}1&2&3&0\\-1&0&-1&2\\2&3&5&-1\\1&1&2&-1\end{pmatrix},\quad
  B=\begin{pmatrix}1&2&3&0\\1&2&2&1\\2&4&3&3\\3&6&7&2\end{pmatrix}
  \]
  とし，列ベクトル空間 $\R^4$ のベクトル $v$ に対して $f_A(v)=Av$, $f_B(v)=Bv$ とおく。\\
  (i) $\operatorname{Im}f_A$ および $\operatorname{Im}f_B$ の基底を1組ずつ求めよ。また，$\operatorname{Im}f_A\cap\operatorname{Im}f_B$ を求めよ。\\
  (ii) $\dim(\ker f_A+\ker f_B)$ を求めよ。
  \item[(2)] $n$ を自然数とする。$f,g$ を $\R^n$ から $\R^n$ への線形写像とし，$\R^n=\operatorname{Im}f\oplus\operatorname{Im}g$ とする。\\
  (i) $\ker(f+g)=\ker f\cap\ker g$ であることを示せ。\\
  (ii) $\R^n=\ker f+\ker g$ であるとき，$f+g$ は全単射であることを示せ。
\end{enumerate}
\end{problem}
\begin{memo*}
\textbf{解答の見通し。}
行列の像は列ベクトルの全ての線形結合、核は $Ax=0$ を満たすベクトル全体である。
まず各列を少数の独立な列の線形結合として式で表し、その独立な列を像の基底にする。
共通部分のベクトルは $A$ 側と $B$ 側の二通りに基底表示できるため、二表示を等置して係数を解く。
核については、列の線形関係を $Ax$ の式へ代入して連立方程式を作る。

$U\oplus W$ は $U+W$ に加えて $U\cap W=\{0\}$ を満たす直和を意味する。
従って $f(x)=-g(x)$ が両方の像に属すれば0でなければならない。
最後は部分空間の和の次元公式を書き下し、$\ker(f+g)=\{0\}$、すなわち単射を示す。
同じ有限次元の空間同士では単射と全射が同値なので、全単射が従う。
\end{memo*}
\begin{proof}[解答]
\begin{enumerate}
  \item[(1)] (i) $A,B$ の列を調べると
  \[
  \operatorname{Im}f_A
  =\left\langle
  \begin{pmatrix}1\\-1\\2\\1\end{pmatrix},
  \begin{pmatrix}2\\0\\3\\1\end{pmatrix}
  \right\rangle,\qquad
  \operatorname{Im}f_B
  =\left\langle
  \begin{pmatrix}1\\1\\2\\3\end{pmatrix},
  \begin{pmatrix}3\\2\\3\\7\end{pmatrix}
  \right\rangle.
  \]
  これら4ベクトルの係数を成分比較すると、両方の像に属するベクトルの係数は
  全て0となる。従って
  \[
  \operatorname{Im}f_A\cap\operatorname{Im}f_B=\{0\}.
  \]

  (ii) 同じ列の関係から
  \[
  \ker f_A=\langle(-1,-1,1,0),(2,-1,0,1)\rangle,
  \]
  \[
  \ker f_B=\langle(-2,1,0,0),(-3,0,1,1)\rangle.
  \]
  $p=(-1,-1,1,0)$, $q=(2,-1,0,1)$, $r=(-2,1,0,0)$ とすると、
  第4のベクトルは $p+q+2r$ であり、$p,q,r$ は一次独立である。従って
  \[
  \dim(\ker f_A+\ker f_B)=3.
  \]

  \item[(2)] (i) $x\in\ker(f+g)$ なら
  $f(x)=-g(x)\in\operatorname{Im}f\cap\operatorname{Im}g=\{0\}$。
  従って $x\in\ker f\cap\ker g$。逆の包含は明らかだから
  \[
  \ker(f+g)=\ker f\cap\ker g.
  \]

  (ii) 階数・退化次数の定理と像の直和から
  \[
  n=\dim\operatorname{Im}f+\dim\operatorname{Im}g
  =2n-\dim\ker f-\dim\ker g.
  \]
  一方、$\R^n=\ker f+\ker g$ だから
  \[
  n=\dim\ker f+\dim\ker g-\dim(\ker f\cap\ker g).
  \]
  よって $\ker f\cap\ker g=\{0\}$。(i) より $f+g$ は単射であり、
  同じ有限次元空間の間の線形写像なので全単射である。
\end{enumerate}
\end{proof}

\begin{problem}{第2問｜滑らかな関数の分解}{2015-2-2}
$f(x,y)$ を $\R^2$ における $C^\infty$ 級関数で $f(0,0)=0$ を満たすものとする。このとき，次の問いに答えよ。
\begin{enumerate}
  \item[(1)] $f(x,y)=\displaystyle\int_0^1\frac{\d}{\d t}f(tx,ty)\,\d t$ を用いて
  \[
  f(x,y)=xg(x,y)+yh(x,y)
  \]
  となる $\R^2$ における $C^\infty$ 級関数 $g(x,y),h(x,y)$ が存在することを示せ。
  \item[(2)] $\R^2$ における $C^\infty$ 級関数 $G(x,y)$ と $\R$ における $C^\infty$ 級関数 $H(y)$ が存在して
  \[
  f(x,y)=xG(x,y)+yH(y)
  \]
  が成り立つことを示せ。
\end{enumerate}
\end{problem}
\begin{memo*}
\textbf{解答の見通し。}
目標の $f=xg+yh$ では係数 $x,y$ を作る必要がある。
そこで原点から $(x,y)$ へ向かう線分 $(tx,ty)$ を考え、
$\varphi(t)=f(tx,ty)$ を連鎖律で微分する。すると
$\varphi'(t)=x\frac{\partial f}{\partial x}(tx,ty)+y\frac{\partial f}{\partial y}(tx,ty)$ となり、必要な $x,y$ が現れる。
$t=0$ から1まで積分すれば、微積分学の基本定理により $f(x,y)-f(0,0)$ が戻る。

二つ目では $H$ を $y$ だけの関数にする必要がある。そのため
$f(x,y)-f(0,y)$ を $x$ 方向だけに積分し、残る $f(0,y)-f(0,0)$ を
$y$ 方向だけに積分する。
\end{memo*}
\begin{proof}[解答]
\begin{enumerate}
  \item[(1)] 連鎖律と $f(0,0)=0$ より
  \begin{align*}
  f(x,y)
  &=\int_0^1\frac{\d}{\d t}f(tx,ty)\,\d t\\
  &=x\int_0^1f_x(tx,ty)\,\d t
   +y\int_0^1f_y(tx,ty)\,\d t.
  \end{align*}
  従って
  \[
  g(x,y)=\int_0^1f_x(tx,ty)\,\d t,\qquad
  h(x,y)=\int_0^1f_y(tx,ty)\,\d t
  \]
  とすればよい。積分記号下で任意回微分できるので $g,h\in C^\infty$。

  \item[(2)]
  \[
  f(x,y)-f(0,y)=x\int_0^1f_x(tx,y)\,\d t,\qquad
  f(0,y)=y\int_0^1f_y(0,ty)\,\d t.
  \]
  よって
  \[
  G(x,y)=\int_0^1f_x(tx,y)\,\d t,\qquad
  H(y)=\int_0^1f_y(0,ty)\,\d t
  \]
  とおけば $f=xG+yH$ であり、$G,H$ はともに $C^\infty$ 級である。
\end{enumerate}
\end{proof}

\begin{problem}{第3問｜コンパクト性と連結性}{2015-2-3}
\begin{enumerate}
  \item[(1)] $S$ をコンパクトな位相空間とし，$T$ を Hausdorff 空間とする。$f:S\to T$ を $S$ から $T$ への連続写像とする。$f(S)$ の1点 $p$ に対し，$f^{-1}(\{p\})$ はコンパクト集合であることを示せ。
  \item[(2)] $\R$ を全ての実数からなる集合とし，$\Q$ を全ての有理数からなる集合とする。$\R^2$ の部分集合 $X,Y$ を $X=\Q^2$, $Y=\R^2\setminus X$ で定める。$X,Y$ に $\R^2$ からの相対位相を導入するとき，$X,Y$ が連結かどうかを判定せよ。
\end{enumerate}
\end{problem}
\begin{memo*}
\textbf{解答の見通し。}
ファイバー $f^{-1}(\{p\})$ は、$p$ へ写る点を全て集めた集合である。
Hausdorff 空間では一点集合が閉、連続写像では閉集合の逆像が閉、
コンパクト空間の閉部分集合はコンパクト、という三段階を順に使う。

$\Q^2$ の非連結性は、$x$ 座標がある無理数 $c$ より小さい側と大きい側に分けて示す。
$c$ 自身は有理点に現れないため二集合で全体を覆える。
$Y=\R^2\setminus\Q^2$ では、各点は少なくとも一方の座標が無理数である。
その無理数座標を固定した縦線または横線は $Y$ の外へ出ないので、折れ線で共通点へ結べる。
\end{memo*}
\begin{proof}[解答]
\begin{enumerate}
  \item[(1)] Hausdorff 空間では $\{p\}$ は閉である。従って連続性より
  $f^{-1}(\{p\})$ はコンパクト空間 $S$ の閉部分集合であり、コンパクトである。

  \item[(2)] $X=\Q^2$ は非連結である。例えば
  \[
  U=X\cap\{x<\sqrt2\},\qquad V=X\cap\{x>\sqrt2\}
  \]
  は非自明な開分離を与える。

  一方 $Y=\R^2\setminus\Q^2$ は道連結である。無理数 $\alpha,\beta$ を固定する。
  $(x,y)\in Y$ について、$x$ が無理数なら
  \[
  (x,y)\longrightarrow(x,\beta)\longrightarrow(\alpha,\beta),
  \]
  $x$ が有理数なら $y$ は無理数なので
  \[
  (x,y)\longrightarrow(\alpha,y)\longrightarrow(\alpha,\beta)
  \]
  という折れ線が全て $Y$ 内にある。従って $Y$ は連結である。
\end{enumerate}
\end{proof}

\section{応用数理コース・専門基礎科目（4題必修）}

\begin{problem}{第1問｜連立方程式と行列}{2015-a-1}
次の問いに答えよ。
\begin{enumerate}
  \item[(1)] 次の連立1次方程式が解をもつときの定数 $a$ の値を求め，そのときのすべての解を求めよ。
  \[
  \begin{cases}
  x+2y+3z=1\\
  4x+5y+6z=-2\\
  7x+8y+9z=a
  \end{cases}
  \]
  \item[(2)] 3つの行列
  \[
  A=\begin{pmatrix}1&2&3\\2&3&1\\3&1&2\end{pmatrix},\quad
  B=\begin{pmatrix}1&2&3\\8&5&5\\3&1&2\end{pmatrix},\quad
  C=\begin{pmatrix}2&1&3\\3&2&1\\1&3&2\end{pmatrix}
  \]
  に対して，$XA=B$, $AY=C$ をみたす行列 $X,Y$ を求めよ。
  \item[(3)] $A$ を $n$ 次正方行列，$I$ を $n$ 次単位行列とするとき，次を証明せよ。\\
  (i) $A^2+A-2I=O$ ならば，$A$ は正則である。ただし $O$ は零行列とする。\\
  (ii) $A$ がべき零行列ならば，$A$ は正則でなく，$I+A$ と $I-A$ は正則である。
\end{enumerate}
\end{problem}
\begin{memo*}
\textbf{解答の見通し。}
連立方程式では、第2式から第1式の4倍、第3式から第1式の7倍を引いて $x$ を消去する。
残った二式の左辺は定数倍の関係になるため、右辺も同じ倍率でなければ解が存在しない。

行列の積は一般に順序を交換できない。$XA=B$ から $X$ を孤立させるには右から
$A^{-1}$、$AY=C$ から $Y$ を孤立させるには左から $A^{-1}$ を掛ける。
その前に $|A|\ne0$ を確認し、余因子行列から逆行列を計算する。
最後の正則性では、逆行列の候補を多項式恒等式から直接作り、左右から掛けて単位行列になることを確認する。
\end{memo*}
\begin{proof}[解答]
\begin{enumerate}
  \item[(1)] 第2式から第1式の4倍、第3式から第1式の7倍を引くと
  \[
  -3y-6z=-6,\qquad -6y-12z=a-7.
  \]
  従って解をもつのは $a=-5$ のときであり、全解は
  \[
  (x,y,z)=(t-3,\,2-2t,\,t)\qquad(t\in\R).
  \]

  \item[(2)]
  \[
  A^{-1}=\frac1{18}
  \begin{pmatrix}-5&1&7\\1&7&-5\\7&-5&1\end{pmatrix}
  \]
  だから
  \[
  X=BA^{-1}=\begin{pmatrix}1&0&0\\0&1&2\\0&0&1\end{pmatrix},
  \qquad
  Y=A^{-1}C=\begin{pmatrix}0&1&0\\1&0&0\\0&0&1\end{pmatrix}.
  \]

  \item[(3)] (i) $A^2+A=2I$ より
  \[
  A\frac{A+I}{2}=\frac{A+I}{2}A=I.
  \]
  従って $A^{-1}=(A+I)/2$。

  (ii) $A^m=O$ とする。$A$ が正則なら $A^m$ も正則となるので矛盾。
  また
  \[
  (I-A)(I+A+\cdots+A^{m-1})=I-A^m=I,
  \]
  \[
  (I+A)(I-A+\cdots+(-1)^{m-1}A^{m-1})=I-(-A)^m=I.
  \]
  従って $I-A$, $I+A$ は正則である。
\end{enumerate}
\end{proof}

\begin{problem}{第2問｜対角化・漸化式・直交行列}{2015-a-2}
\begin{enumerate}
  \item[(1)] $A=\begin{pmatrix}-2&1&2\\1&0&0\\0&1&0\end{pmatrix}$ とするとき，次の問いに答えよ。\\
  (i) $A$ の固有値と各固有値に対する固有空間を求めよ。\\
  (ii) $A$ が対角化可能かどうか理由を含めて答えよ。また，対角化可能ならば対角化し，対角化するための正則行列とその逆行列も求めよ。\\
  (iii) $A^n$ の1行目の成分をすべて求めよ。\\
  (iv) 漸化式
  \[
  x_n=-2x_{n-1}+x_{n-2}+2x_{n-3}\qquad(n\ge4)
  \]
  について，初期条件が $x_1=1$, $x_2=0$, $x_3=0$ の場合と，初期条件が $x_1=0$, $x_2=-1$, $x_3=1$ の場合の一般項 $x_n$ をそれぞれ求めよ。
  \item[(2)] $A$ を $n$ 次正方行列とし，
  \begin{enumerate}
    \item[(ア)] $A$ は直交行列である
    \item[(イ)] $\R^n$ に属する任意のベクトル $x$ に対して，$\norm{Ax}=\norm{x}$ である
    \item[(ウ)] $\R^n$ に属する任意のベクトル $x,y$ に対して，$\inner{Ax}{Ay}=\inner{x}{y}$ である
  \end{enumerate}
  とする。ただし $\R^n$ に属するベクトル $a,b$ に対して，$\inner{a}{b}$ は $a$ と $b$ の内積を表し，$\norm{a}=\sqrt{\inner{a}{a}}$ とする。このとき，次を証明せよ。\\
  (i) (ア) が成立するならば，(イ) が成立する。\\
  (ii) (イ) が成立するならば，(ウ) が成立する。\\
  (iii) (ウ) が成立するならば，(ア) が成立する。
\end{enumerate}
\end{problem}
\begin{memo*}
\textbf{解答の見通し。}
まず $|\lambda I-A|=0$ を因数分解して固有値を求め、
各 $\lambda$ を $(A-\lambda I)x=0$ へ代入して固有ベクトルを得る。
固有ベクトルを列に並べた $P$ では $AP=PD$ が成り立つため、左から $P^{-1}$ を掛けて
$P^{-1}AP=D$、さらに $A^n=PD^nP^{-1}$ とできる。

漸化式へ試行解 $x_n=r^n$ を代入すると、行列と同じ三根 $1,-1,-2$ を持つ特性方程式が現れる。
根が相異なるので一般解を三つの指数列の線形結合として置き、初期値を代入して係数を解く。
直交行列の同値条件では、$\norm{x}^2=\inner{x}{x}$ と
$\norm{x+y}^2$ の展開を使い、「行列条件・長さ・内積」を順につなぐ。
\end{memo*}
\begin{proof}[解答]
\begin{enumerate}
  \item[(1)] (i)
  \[
  \det(\lambda I-A)=(\lambda-1)(\lambda+1)(\lambda+2),
  \]
  \[
  E_1=\langle(1,1,1)^{\mathsf T}\rangle,\quad
  E_{-1}=\langle(1,-1,1)^{\mathsf T}\rangle,\quad
  E_{-2}=\langle(4,-2,1)^{\mathsf T}\rangle.
  \]

  (ii) 固有値が相異なるので対角化可能である。
  \[
  P=\begin{pmatrix}1&1&4\\1&-1&-2\\1&1&1\end{pmatrix},\qquad
  P^{-1}=\frac16\begin{pmatrix}1&3&2\\-3&-3&6\\2&0&-2\end{pmatrix},
  \]
  \[
  P^{-1}AP=\diag(1,-1,-2).
  \]

  (iii) $A^n=P\diag(1,(-1)^n,(-2)^n)P^{-1}$ より、第1行は
  \[
  \left(
  \frac16-\frac{(-1)^n}{2}+\frac43(-2)^n,\quad
  \frac{1-(-1)^n}{2},\quad
  \frac13+(-1)^n-\frac43(-2)^n
  \right).
  \]

  (iv) 特性方程式は $(r-1)(r+1)(r+2)=0$ だから
  \[
  x_n=\alpha+\beta(-1)^n+\gamma(-2)^n.
  \]
  各初期条件を代入すると
  \[
  (x_1,x_2,x_3)=(1,0,0):
  \quad x_n=\frac13-(-1)^n+\frac16(-2)^n,
  \]
  \[
  (x_1,x_2,x_3)=(0,-1,1):
  \quad x_n=-\frac13-\frac16(-2)^n.
  \]

  \item[(2)] (i) $A^{\mathsf T}A=I$ なら
  \[
  \|Ax\|^2=x^{\mathsf T}A^{\mathsf T}Ax=\|x\|^2.
  \]

  (ii) 長さを保つなら、実内積の極化恒等式より
  \[
  \langle Ax,Ay\rangle
  =\frac{\|A(x+y)\|^2-\|Ax\|^2-\|Ay\|^2}{2}
  =\langle x,y\rangle.
  \]

  (iii) 内積を保つなら、標準基底 $e_i,e_j$ に対して
  \[
  (A^{\mathsf T}A)_{ij}
  =\langle Ae_i,Ae_j\rangle
  =\langle e_i,e_j\rangle=\delta_{ij}.
  \]
  従って $A^{\mathsf T}A=I$、すなわち $A$ は直交行列である。
\end{enumerate}
\end{proof}

\begin{problem}{第3問｜初等微積分}{2015-a-3}
次の問いに答えよ。
\begin{enumerate}
  \item[(1)] 次の関数の導関数を求めよ。\\
  (i) $\log|\tan x|$\\
  (ii) $x^{\sin x}$ $(x>0)$
  \item[(2)] (i) $x+\sqrt{x^2+1}=t$ のとき，$x$ を $t$ だけを用いて表わせ。\\
  (ii) 次の定積分の値を求めよ。
  \[
  \text{(a)}\ \int_0^1\frac{\d x}{\sqrt{x^2+1}}
  \qquad
  \text{(b)}\ \int_0^1\sqrt{x^2+1}\,\d x
  \]
\end{enumerate}
\end{problem}
\begin{memo*}
\textbf{解答の見通し。}
$\log|\tan x|$ は $\tan x\ne0$ の区間ごとに、合成関数の微分公式
$(\log|u|)'=u'/u$ を使う。$x^{\sin x}$ は実数値として $x>0$ で考え、
$x^{\sin x}=\exp(\sin x\log x)$ と書いて指数関数と積の微分へ直す。

置換 $t=x+\sqrt{x^2+1}$ では、共役な式を掛けると
\[
(\sqrt{x^2+1}+x)(\sqrt{x^2+1}-x)=1
\]
となるため、$t^{-1}=\sqrt{x^2+1}-x$ が得られる。
$t$ と $t^{-1}$ を加減して $x$ と平方根を表し、$\d x$ と積分端点も全て $t$ へ変換する。
\end{memo*}
\begin{proof}[解答]
\begin{enumerate}
  \item[(1)] (i)
  \[
  \frac{\d}{\d x}\log|\tan x|
  =\frac{\sec^2x}{\tan x}
  =\frac1{\sin x\cos x}.
  \]
  (ii) $x^{\sin x}=e^{\sin x\log x}$ より
  \[
  \frac{\d}{\d x}x^{\sin x}
  =x^{\sin x}\left(\cos x\log x+\frac{\sin x}{x}\right).
  \]

  \item[(2)] (i)
  \[
  t^{-1}=\sqrt{x^2+1}-x,\qquad
  x=\frac{t-t^{-1}}2.
  \]

  (ii) さらに
  \[
  \sqrt{x^2+1}=\frac{t+t^{-1}}2,\qquad
  \d x=\frac{1+t^{-2}}2\,\d t,
  \]
  で、$x=0,1$ は $t=1,1+\sqrt2$ に対応する。従って
  \[
  \text{(a)}\quad
  \int_0^1\frac{\d x}{\sqrt{x^2+1}}
  =\int_1^{1+\sqrt2}\frac{\d t}{t}
  =\log(1+\sqrt2),
  \]
  \[
  \text{(b)}\quad
  \int_0^1\sqrt{x^2+1}\,\d x
  =\frac14\int_1^{1+\sqrt2}(t+2t^{-1}+t^{-3})\,\d t
  =\frac{\sqrt2+\log(1+\sqrt2)}2.
  \]
\end{enumerate}
\end{proof}

\begin{problem}{第4問｜最大最小と三重積分}{2015-a-4}
次の問いに答えよ。
\begin{enumerate}
  \item[(1)] $\sqrt x+\sqrt y=1$ の条件の下で，$x^2+xy+y^2$ の最大値と最小値を求めよ。
  \item[(2)] 次の積分の値を求めよ。
  \[
  \iiint_V(x+y+z)\,\d x\d y\d z
  \]
  ただし $V=\{(x,y,z)\mid x^2+y^2+z^2\le a^2\}$ $(a>0)$ である。
\end{enumerate}
\end{problem}
\begin{memo*}
\textbf{解答の見通し。}
平方根が現れるので $u=\sqrt x$, $v=\sqrt y$ と置く。このとき $x=u^2$, $y=v^2$、
$u,v\ge0$ で、制約は直線 $u+v=1$ になる。
目的関数を展開し、$u+v=1$ を使って積 $q=uv$ だけの式へ直す。
相加相乗平均から $0\le q\le1/4$ なので、一変数二次関数の増減として最大・最小を決める。

三重積分では、球内の各点 $(x,y,z)$ と反対側の点 $(-x,-y,-z)$ が対になり、
領域の体積要素は変わらない一方、$x+y+z$ の符号だけが反転する。
変数変換を式にすると積分値 $J$ が $J=-J$ を満たし、$J=0$ となる。
\end{memo*}
\begin{proof}[解答]
\begin{enumerate}
  \item[(1)] $u=\sqrt x$, $v=\sqrt y$, $q=uv$ とおくと
  $u,v\ge0$, $u+v=1$, $0\le q\le1/4$ であり、
  \[
  x^2+xy+y^2=u^4+u^2v^2+v^4=1-4q+3q^2.
  \]
  これは $0\le q\le1/4$ で単調減少する。従って
  \[
  \max=1\quad ((x,y)=(1,0),(0,1)),
  \]
  \[
  \min=\frac3{16}\quad ((x,y)=(1/4,1/4)).
  \]

  \item[(2)] $V$ は原点対称で、$x+y+z$ は奇関数だから
  \[
  \iiint_V(x+y+z)\,\d x\d y\d z=0.
  \]
\end{enumerate}
\end{proof}

\chapter{年度不明の問題}
\label{ch:unknown}

\begin{memo*}
元資料に年度表記がないため，推測で年度を付けずに掲載する。写真の問題番号から，少なくとも二つの試験冊子が混在している。
\end{memo*}

\section{専門基礎科目と思われる問題}

\begin{problem}{行列｜直交対角化}{unknown-1}
行列
\[
P=\begin{pmatrix}
1/\sqrt3&1/\sqrt2&p\\
q&r&s\\
-1/\sqrt3&1/\sqrt2&-1/\sqrt6
\end{pmatrix}
\qquad(p,q,r,s\in\R)
\]
は直交行列であるとする。
\begin{enumerate}
  \item[(1)] $p,q,r,s$ の取りうる組み合わせを全て求めよ。
  \item[(2)] 行列 ${}^{\mathsf T}(P^{10})\,P^{19}\,({}^{\mathsf T}P)^{10}$ を求めよ。
  \item[(3)] 行列 $A$ は $P$ で対角化されるとし，$A$ の固有多項式 $g_A(x)$ は $x^3+\sqrt[10]{2}\,x^2$ であるとする。また，ベクトル $(1,0,1)^{\mathsf T}$ は固有値 $-\sqrt[10]{2}$ に属する $A$ の固有ベクトルであるとする。このとき $A^{10}$ を求めよ。
  \item[(4)] 直交行列で対角化される行列は対称行列であることを証明せよ。
\end{enumerate}
\end{problem}
\begin{memo*}
\textbf{解答の見通し。}
直交行列は $P^{\mathsf T}P=I$ を満たす行列である。これは、各列の長さが1で、
異なる二列の内積が0であることと同値なので、その条件を式にして $p,q,r,s$ を順に決める。
$P^{\mathsf T}=P^{-1}$ だから、中央の行列積は同じ底 $P$ の整数乗へ直して指数を加減できる。

$A$ の固有値は固有多項式から0,0,$-\sqrt[10]{2}$ と分かる。
非零固有値の単位固有ベクトルを $u$ とし、残り二本の単位固有ベクトルを加えて直交行列 $Q$ を作ると、
$A=Q\diag(-\sqrt[10]{2},0,0)Q^{\mathsf T}$ である。この積を列ごとに展開すると
$A=-\sqrt[10]{2}\,uu^{\mathsf T}$ となり、10乗を計算できる。
\end{memo*}
\begin{proof}[解答]
\begin{enumerate}
  \item[(1)] 列ベクトルの正規直交条件から
  \[
  r=0,\qquad q=\frac{\varepsilon}{\sqrt3},\qquad
  p=\frac1{\sqrt6},\qquad
  s=-\varepsilon\sqrt{\frac23}
  \qquad(\varepsilon=\pm1).
  \]
  これらはいずれも第3列の長さも1にするので、全ての組である。

  \item[(2)] $P^{\mathsf T}=P^{-1}$ より
  \[
  (P^{10})^{\mathsf T}P^{19}(P^{\mathsf T})^{10}
  =P^{-10}P^{19}P^{-10}=P^{-1}=P^{\mathsf T}.
  \]

  \item[(3)] 固有値は $0,0,-\sqrt[10]{2}$ である。
  \[
  u=\frac1{\sqrt2}(1,0,1)^{\mathsf T}
  \]
  とすると、直交対角化より $A=-\sqrt[10]{2}\,uu^{\mathsf T}$。
  $u^{\mathsf T}u=1$ だから $(uu^{\mathsf T})^k=uu^{\mathsf T}$ であり、
  \[
  A^{10}=2uu^{\mathsf T}
  =\begin{pmatrix}1&0&1\\0&0&0\\1&0&1\end{pmatrix}.
  \]

  \item[(4)] 直交行列 $Q$ と実対角行列 $D$ により
  $Q^{-1}AQ=D$ と書けるなら $A=QDQ^{\mathsf T}$。従って
  \[
  A^{\mathsf T}=QD^{\mathsf T}Q^{\mathsf T}=QDQ^{\mathsf T}=A.
  \]
\end{enumerate}
\end{proof}

\begin{problem}{微分積分｜積分記号下の微分}{unknown-2}
$a,b,c,d\in\R$ に対し $a<b$, $c<d$ を仮定する。$D$ は $\R^2$ の開集合で $[a,b]\times[c,d]$ を含むとする。$f(x,y)$ を $D$ 上の $C^1$ 級関数とする。
\begin{enumerate}
  \item[(1)] $[a,b]$ 上で
  \[
  \frac{\d}{\d x}\int_c^d f(x,y)\,\d y=\int_c^d\frac{\partial f}{\partial x}(x,y)\,\d y
  \]
  が成り立つことを示せ。ただし $G(x)=\int_c^d\frac{\partial f}{\partial x}(x,y)\,\d y$ が $[a,b]$ 上で連続であることを仮定して良い。
  \item[(2)] $\phi(x),\psi(x)$ は $[a,b]$ を含む開区間上の $C^1$ 級関数で，任意の $x\in[a,b]$ に対し $\phi(x),\psi(x)\in[c,d]$ が成り立つとする。このとき
  \[
  \frac{\d}{\d x}\int_{\phi(x)}^{\psi(x)}f(x,y)\,\d y
  =f(x,\psi(x))\frac{\d\psi}{\d x}(x)-f(x,\phi(x))\frac{\d\phi}{\d x}(x)
  +\int_{\phi(x)}^{\psi(x)}\frac{\partial f}{\partial x}(x,y)\,\d y
  \]
  が成り立つことを示せ。
  \item[(3)] $x>0$ に対し
  \[
  p(x)=\int_{x^2}^{x^3}\frac{\sin(xy)}y\,\d y
  \]
  とおくとき，$\dfrac{\d p}{\d x}(x)$ を求めよ。
\end{enumerate}
\end{problem}
\begin{memo*}
\textbf{解答の見通し。}
固定端点とは積分区間 $[c,d]$ が $x$ によらない場合である。
差商 $\{F(x+h)-F(x)\}/h$ を書き、$f$ の差だけを積分した形へ直す。
一変数の微積分学の基本定理を $x$ 方向へ使うと、その差商は $\partial f/\partial x$ の平均になる。
$\partial f/\partial x$ の連続性により平均が一様に $\dfrac{\partial f}{\partial x}(x,y)$ へ近づくため、極限と積分を交換できる。

可変端点では $Q(x,v)=\int_c^v f(x,y)\,\d y$ と置き、
$Q(x,\psi(x))-Q(x,\phi(x))$ を二変数の連鎖律で微分する。
これにより、上端の移動、下端の移動、被積分関数自身の変化という三つの項が現れる。
\end{memo*}
\begin{proof}[解答]
\begin{enumerate}
  \item[(1)] $F(x)=\int_c^d f(x,y)\,\d y$ とおく。差商は
  \[
  \frac{F(x+h)-F(x)}h
  =\int_c^d\int_0^1 f_x(x+th,y)\,\d t\d y.
  \]
  $f_x$ の連続性により $h\to0$ で被積分関数は一様に $f_x(x,y)$ へ収束するから
  \[
  F'(x)=\int_c^d f_x(x,y)\,\d y.
  \]

  \item[(2)] $Q(x,v)=\int_c^v f(x,y)\,\d y$ とおくと
  \[
  Q_x(x,v)=\int_c^v f_x(x,y)\,\d y,\qquad Q_v(x,v)=f(x,v).
  \]
  $\int_{\phi(x)}^{\psi(x)}f(x,y)\,\d y
  =Q(x,\psi(x))-Q(x,\phi(x))$ に連鎖律を用いれば
  \[
  \frac{\d}{\d x}\int_{\phi(x)}^{\psi(x)}f(x,y)\,\d y
  =f(x,\psi(x))\psi'(x)-f(x,\phi(x))\phi'(x)
  +\int_{\phi(x)}^{\psi(x)}f_x(x,y)\,\d y.
  \]

  \item[(3)] $f(x,y)=\sin(xy)/y$, $\phi(x)=x^2$, $\psi(x)=x^3$ として (2) を使うと
  \begin{align*}
  p'(x)
  &=\frac{3\sin(x^4)-2\sin(x^3)}x
    +\int_{x^2}^{x^3}\cos(xy)\,\d y\\
  &=\frac{4\sin(x^4)-3\sin(x^3)}x.
  \end{align*}
\end{enumerate}
\end{proof}

\begin{problem}{距離空間｜コンパクト集合の近傍}{unknown-3}
$(X,d)$ を距離空間とし，$K$ を $X$ のコンパクトな部分集合とする。関数 $\rho:X\to\R$ を
\[
\rho(x)=\inf\{d(x,z)\mid z\in K\}\qquad(x\in X)
\]
で定める。
\begin{enumerate}
  \item[(1)] すべての $x,y\in X$ に対して $|\rho(x)-\rho(y)|\le d(x,y)$ が成り立つことを示せ。
  \item[(2)] 任意の $x\in X$ に対して，ある $a\in K$ が存在して $\rho(x)=d(x,a)$ を満たすことを示せ。
  \item[(3)] $L$ は $K$ を含む $X$ の開集合とする。このとき $K_r=\{x\in X\mid\rho(x)<r\}\subset L$ となる実数 $r>0$ が存在することを示せ。
\end{enumerate}
\end{problem}
\begin{memo*}
\textbf{解答の見通し。}
$\rho(x)$ は、$x$ と $K$ の各点との距離の下限である。
三角不等式を全ての $z\in K$ について書き、下限を取ると
$|\rho(x)-\rho(y)|\le d(x,y)$ が得られる。
最短点の存在は、連続関数 $z\mapsto d(x,z)$ がコンパクト集合 $K$ 上で最小値を取ることから示す。

$K$ を含む開集合 $L$ の外側を $F=X\setminus L$ とする。
$F$ が空なら $L=X$ なので、どの $r>0$ を選んでも $\{\rho<r\}\subset L$ である。
$F$ が空でなければ、各 $z\in K$ と閉集合 $F$ の距離は正であり、
それを $K$ 上で最小化した値 $\delta$ も正になる。$0<r<\delta$ を選べば、
$K$ から距離 $r$ 未満の点が $F$ に入ることはできない。
\end{memo*}
\begin{proof}[解答]
\begin{enumerate}
  \item[(1)] 任意の $z\in K$ に対し
  \[
  \rho(x)\le d(x,z)\le d(x,y)+d(y,z).
  \]
  $z$ について下限を取り、$x,y$ を交換した式と合わせれば
  \[
  |\rho(x)-\rho(y)|\le d(x,y).
  \]

  \item[(2)] $z\mapsto d(x,z)$ は $K$ 上連続である。$K$ はコンパクトだから
  最小値を取る $a\in K$ があり、
  \[
  \rho(x)=\min_{z\in K}d(x,z)=d(x,a).
  \]

  \item[(3)] 各 $z\in K$ に対し、$L$ が開なので
  $B(z,2\varepsilon_z)\subset L$ となる $\varepsilon_z>0$ を取る。
  コンパクト性より
  \[
  K\subset B(z_1,\varepsilon_{z_1})\cup\cdots\cup
  B(z_m,\varepsilon_{z_m}).
  \]
  $r=\min_i\varepsilon_{z_i}>0$ とする。$\rho(x)<r$ なら (2) より
  $d(x,a)<r$ となる $a\in K$ がある。ある $i$ で
  $d(a,z_i)<\varepsilon_{z_i}$ だから
  \[
  d(x,z_i)<r+\varepsilon_{z_i}\le2\varepsilon_{z_i}.
  \]
  従って $x\in L$、すなわち $K_r\subset L$。
\end{enumerate}
\end{proof}

\section{専門科目と思われる問題}

\begin{problem}{回転面｜トーラス}{unknown-4}
$R>r>0$ とする。$\R^3$ において $xz$ 平面上の円
\[
\{(x,0,z)\in\R^3\mid(x-R)^2+z^2=r^2\}
\]
を $z$ 軸のまわりに回転して得られる曲面を $S$ とする。
\begin{enumerate}
  \item[(1)] 曲面 $S$ のガウス曲率 $K$ と平均曲率 $H$ を求めよ。
  \item[(2)] 曲面 $S$ 上の測地線の例を一つあげ，それが測地線となることを説明せよ。
\end{enumerate}
\end{problem}
\begin{memo*}
\textbf{解答の見通し。}
回転角 $u$ と母円の角 $v$ で標準的にパラメータ表示する。
二つの座標ベクトルは直交するため、第1・第2基本形式から主曲率を直接読める。
測地線の例には外側の赤道 $v=0$ を選び、
その加速度が曲面の法線方向だけであることを確認する。
\end{memo*}
\begin{proof}[解答]
\begin{enumerate}
\item[(1)]
\[
X(u,v)=((R+r\cos v)\cos u,(R+r\cos v)\sin u,r\sin v)
\]
とおく。まず偏微分は
\[
X_u=(-(R+r\cos v)\sin u,(R+r\cos v)\cos u,0),
\]
\[
X_v=(-r\sin v\cos u,-r\sin v\sin u,r\cos v).
\]
従って第1基本形式の係数は
\[
E=(R+r\cos v)^2,\qquad F=0,\qquad G=r^2.
\]
$R>r$ より $R+r\cos v>0$ なので $EG-F^2>0$、従ってこれは正則な曲面である。
外向き単位法線を
\[
N=(\cos v\cos u,\cos v\sin u,\sin v)
\]
と取る。二階偏微分と内積を取ると
\[
L=X_{uu}\cdot N=-(R+r\cos v)\cos v,\quad
M=X_{uv}\cdot N=0,\quad
N_2=X_{vv}\cdot N=-r.
\]
$F=M=0$ なので座標方向が主方向であり、主曲率は
\[
k_1=\frac LE=-\frac{\cos v}{R+r\cos v},\qquad
k_2=\frac{N_2}G=-\frac1r.
\]
従って
\[
K=\frac{\cos v}{r(R+r\cos v)},\qquad
H=-\frac{R+2r\cos v}{2r(R+r\cos v)}.
\]
\item[(2)] $v=0$ と固定した外側の赤道を $\gamma(u)=X(u,0)$ とおくと
\[
\gamma''(u)=-(R+r)(\cos u,\sin u,0)=-(R+r)N(u,0).
\]
加速度の接平面成分が0、すなわち測地的曲率が0なので、この円は測地線である。
\end{enumerate}
\end{proof}

\begin{problem}{Euler 型微分方程式}{unknown-5}
$p,q\in\R$ とする。微分方程式に対する初期値問題
\[
x^2\frac{\d^2y}{\d x^2}+px\frac{\d y}{\d x}+qy=0\quad(x\ge1),
\qquad
y(1)=0,\quad \frac{\d y}{\d x}(1)=1
\tag{1}
\]
を考える。
\begin{enumerate}
  \item[(1)] $p\ne1$ かつ $q=0$ とするとき，(1) の解 $y(x)$ を $p$ を用いて表せ。
  \item[(2)] 変数変換 $x=e^t$ によって，(1) の第一式を変形して得られる方程式を求めよ。
  \item[(3)] $p=2$ とするとき，(1) の解 $y(x)$ を $q$ を用いて表せ。
  \item[(4)] (3) で得られた解 $y(x)$ について，$y(x)=0$ となる $x>1$ が存在するような $q$ の範囲を求め，そのような $x>1$ の最小値を $q$ を用いて表せ。
\end{enumerate}
\end{problem}
\begin{memo*}
\textbf{解答の見通し。}
Euler 型方程式は対数変数 $t=\log x$ へ移すと定数係数方程式になる。
$p=2$ では判別式 $d=1-4q$ の符号で、双曲線関数・重根・三角関数の三場合に分ける。
$x>1$ は $t>0$ に対応し、零点が繰り返し現れるのは三角関数になる
$d<0$、すなわち $q>1/4$ の場合だけである。
\end{memo*}
\begin{proof}[解答]
\begin{enumerate}
\item[(1)] $q=0$ では $x^2y''+pxy'=0$ である。$z=y'$ とおくと
\[
xz'+pz=0,\qquad \frac{z'}z=-\frac px.
\]
従って $z=Cx^{-p}$。初期条件 $y'(1)=1$ から $C=1$ で、$y(1)=0$ を使って
\[
y(x)=\int_1^x s^{-p}\,\d s=\frac{x^{1-p}-1}{1-p}\qquad(p\ne1).
\]

\item[(2)] $x=e^t$、$Y(t)=y(e^t)$ とおく。連鎖律から
\[
Y'(t)=e^ty'(e^t)=xy'(x),
\]
\[
Y''(t)=xy'(x)+x^2y''(x),
\]
従って $x^2y''=Y''-Y'$ である。元の方程式は
\[
Y''+(p-1)Y'+qY=0,\qquad Y(0)=0, Y'(0)=1.
\]
\item[(3)] $p=2$ では特性方程式が $\lambda^2+\lambda+q=0$ となる。
$d=1-4q$ とおく。

$d>0$ のとき、特性根は $(-1\pm\sqrt d)/2$ なので
\[
Y(t)=e^{-t/2}\left(C_1e^{\sqrt d\,t/2}+C_2e^{-\sqrt d\,t/2}\right).
\]
$Y(0)=0$ から $C_2=-C_1$。従って
$Y(t)=2C_1e^{-t/2}\sinh(\sqrt d\,t/2)$ であり、
\[
Y'(0)=2C_1\cdot\frac{\sqrt d}{2}=C_1\sqrt d=1
\]
から $C_1=1/\sqrt d$ を得る。

$d=0$ のときは重根 $-1/2$ を持つので
$Y(t)=(C_1+C_2t)e^{-t/2}$。初期条件から $C_1=0$, $C_2=1$ である。

$d<0$ のとき、$b=\sqrt{-d}$ とおけば特性根は $(-1\pm ib)/2$ であり
\[
Y(t)=e^{-t/2}\left(C_1\cos\frac{bt}{2}+C_2\sin\frac{bt}{2}\right).
\]
$Y(0)=0$ から $C_1=0$、さらに
$Y'(0)=C_2b/2=1$ から $C_2=2/b$ となる。

最後に $t=\log x$、$e^{-t/2}=x^{-1/2}$ を戻すと
\[
y=x^{-1/2}\begin{cases}
\dfrac{2}{\sqrt d}\sinh(\frac{\sqrt d}{2}\log x),&d>0,\\
\log x,&d=0,\\
\dfrac{2}{\sqrt{-d}}\sin(\frac{\sqrt{-d}}2\log x),&d<0.
\end{cases}
\]
\item[(4)] $d\ge0$ の場合、$x>1$ では $\log x>0$ であり、$\sinh u$ と $u$ は正なので零点を持たない。
$d<0$、すなわち $q>1/4$ の場合は
\[
\sin\left(\frac{\sqrt{4q-1}}2\log x\right)=0
\]
より、正の零点は
\[
\frac{\sqrt{4q-1}}2\log x=k\pi\qquad(k=1,2,\ldots).
\]
従って $x>1$ の零点があるのは $q>1/4$ のときに限り、最小の零点は
$\exp(2\pi/\sqrt{4q-1})$ である。
\end{enumerate}
\end{proof}

\begin{problem}{解析｜弱い特異核}{unknown-6}
$\alpha$ を $0<\alpha<1$ を満たす実数とする。$[0,\infty)$ 上で有界なルベーグ可測関数 $\phi(t)$ に対し，$[0,\infty)$ 上の関数 $f(x)$ を
\[
f(x)=\int_0^x\phi(t)\,\d t
\]
で定める。
\begin{enumerate}
  \item[(1)] $f(x)$ は $[0,\infty)$ 上で連続であることを示せ。
  \item[(2)] $(0,\infty)$ 上の関数 $g(x)$ を
  \[
  g(x)=\int_0^x\frac{f(s)}{(x-s)^\alpha}\,\d s
  \]
  で定める。任意の $x\in(0,\infty)$ に対して
  \[
  g(x)=\frac1{1-\alpha}\int_0^x\phi(t)(x-t)^{1-\alpha}\,\d t
  \]
  が成り立つことを示せ。
  \item[(3)] $g(x)$ を (2) で定めた $(0,\infty)$ 上の関数とする。$\{h_n\}$ は正の実数列で，$n\to\infty$ のとき $h_n\to0$ とする。任意の $x\in(0,\infty)$ に対して
  \[
  \lim_{n\to\infty}\frac{g(x+h_n)-g(x)}{h_n}=\int_0^x\frac{\phi(t)}{(x-t)^\alpha}\,\d t
  \]
  が成り立つことを示せ。
\end{enumerate}
\end{problem}
\begin{memo*}
\textbf{解答の見通し。}
$\phi$ の有界性から原始関数 $f$ は Lipschitz 連続になる。
$f(s)=\int_0^s\phi(t)\,\d t$ を代入すると三角形領域の二重積分になるので、
$0<\alpha<1$ による絶対可積分性を確認して Fubini の定理を使う。
導関数候補 $G$ を先に定め、$\int_0^xG(u)\,\d u=g(x)$ を示して
微積分学の基本定理へ帰着する。
\end{memo*}
\begin{proof}[解答]
\begin{enumerate}
\item[(1)] $|f(y)-f(x)|\le\norm\phi_\infty|y-x|$ なので $f$ は連続である。
\item[(2)]$f(s)=\int_0^s\phi(t)dt$ を代入し，Fubini で積分順序を交換すると
\begin{align*}
g(x)
&=\int_0^x\int_0^s\phi(t)(x-s)^{-\alpha}\,\d t\d s\\
&=\int_0^x\phi(t)\int_t^x(x-s)^{-\alpha}\,\d s\d t.
\end{align*}
ここで $|\phi(t)|(x-s)^{-\alpha}$ は三角形 $0\le t\le s\le x$ 上可積分である。実際、
\[
\int_0^x\int_0^s\norm\phi_\infty(x-s)^{-\alpha}\,\d t\d s
=\norm\phi_\infty\int_0^x s(x-s)^{-\alpha}\,\d s<\infty
\]
なので Fubini の定理を適用できる。内側を $u=x-s$ で積分すると
\[
\int_t^x(x-s)^{-\alpha}\,\d s
=\int_0^{x-t}u^{-\alpha}\,\d u
=\frac{(x-t)^{1-\alpha}}{1-\alpha},
\]
従って指定された表示式を得る。

\item[(3)] 導関数の候補を
\[
G(x)=\int_0^x\phi(t)(x-t)^{-\alpha}\,\d t
\]
とおく。$0<\alpha<1$ なのでこの積分は絶対収束する。さらに Fubini の定理により
\begin{align*}
\int_0^xG(u)\,\d u
&=\int_0^x\int_0^u\phi(t)(u-t)^{-\alpha}\,\d t\d u\\
&=\int_0^x\phi(t)\int_t^x(u-t)^{-\alpha}\,\d u\d t\\
&=\frac1{1-\alpha}\int_0^x\phi(t)(x-t)^{1-\alpha}\,\d t\\
&=g(x).
\end{align*}
$G$ は局所連続である。実際、核 $u^{-\alpha}\mathbf1_{(0,L)}(u)$ は
$L^1(\R)$ に属し、その平行移動は $L^1$ で連続、$\phi$ は有界だからである。
従って微積分学の基本定理を適用でき
\[
g'(x)=G(x)=\int_0^x\phi(t)(x-t)^{-\alpha}\,\d t
\]
を得る。特に問題の極限は $h_n>0$, $h_n\to0$ に対する差商の極限なので、この値に等しい。
\end{enumerate}
\end{proof}

\chapter{専門基礎科目の出題傾向と予想問題集}
\label{ch:prediction}

\begin{memo*}
この章は、収録済みの2015・2017・2021・2022・2026年度の問題から、
繰り返し現れる分野と解法を整理して作成した非公式の予想問題集である。
実際の出題を保証するものではない。まず問題だけを時間を測って解き、
その後に「解答の見通し」と解答を確認する使い方を想定している。
三題を必修と想定し、一題40分を目安にすると、120分の試験として練習できる。
\end{memo*}

\section{過去問から見える傾向}

\subsection{専門基礎科目}

年度によって募集期や問題数に違いはあるが、数学コースの専門基礎では
次の三分野が繰り返し組み合わされている。
\begin{enumerate}
  \item \textbf{線形代数}：
  固有値・固有空間、対角化、Jordan 標準形、不変部分空間、直交行列・ユニタリ行列。
  単に答えを出すだけでなく、基底を具体的に構成する問題が多い。
  \item \textbf{微分積分}：
  多変数関数、広義積分、重積分、Gamma・Beta 関数、積分記号下の微分。
  積分領域の分割や変数変換を自分で選ぶ力が問われる。
  \item \textbf{位相・距離空間}：
  コンパクト性、連結性、連続写像、部分空間、集合までの距離。
  定義を量化記号の順序まで正確に使う証明が中心である。
\end{enumerate}
特に2021・2022年度の複数の募集期と2026年度第2期では、
この三分野がほぼ一題ずつ並んでいる。従って、予想問題も三題必修の形にした。

\section{予想問題（専門基礎3題必修を想定）}

\begin{problem}{予想第1問｜実対称行列の直交対角化}{prediction-b-1}
\[
A=\begin{pmatrix}
2&1&1\\
1&2&1\\
1&1&2
\end{pmatrix}
\]
とする。
\begin{enumerate}
  \item[(1)] $|A|$ を求めよ。
  \item[(2)] $A$ の固有値と各固有空間を求めよ。
  \item[(3)] $A$ を直交行列によって対角化せよ。
  \item[(4)] $x^2+y^2+z^2=1$ のもとで
  \[
  Q(x,y,z)=
  \begin{pmatrix}x&y&z\end{pmatrix}
  A
  \begin{pmatrix}x\\y\\z\end{pmatrix}
  \]
  の最大値と最小値を求め、等号成立条件も述べよ。
\end{enumerate}
\end{problem}
\begin{memo*}
\textbf{出題根拠。}
固有値、固有空間、直交対角化は、収録した専門基礎の線形代数で最も反復の多い流れである。

\textbf{解答の見通し。}
固有値 $\lambda$ は、連立方程式
$(A-\lambda I_3)\boldsymbol{x}=\boldsymbol{0}$ が
零ベクトル以外の解をもつような数である。その条件は
$|A-\lambda I_3|=0$ である。
そこで、まず $A-\lambda I_3$ を書き、
第2行と第3行から第1行を引く。この行基本変形によって
特性多項式を $(1-\lambda)^2(4-\lambda)$ まで因数分解する。
次に $\lambda=1,4$ を一つずつ
$(A-\lambda I_3)\boldsymbol{x}=\boldsymbol{0}$ へ代入して固有空間を求める。
最後に、各固有空間から互いに直交する長さ1のベクトルを選び、
それらを列に並べて直交行列を作る。
\end{memo*}
\begin{proof}[解答]
\begin{enumerate}
  \item[(1)] 行列式では、ある行に別の行の定数倍を加えても値は変わらない。
  第2行、第3行からそれぞれ第1行を引くと
  \[
  |A|
  =\begin{vmatrix}
  2&1&1\\1&2&1\\1&1&2
  \end{vmatrix}
  =\begin{vmatrix}
  2&1&1\\-1&1&0\\-1&0&1
  \end{vmatrix}.
  \]
  右辺を第1行で展開すると
  \[
  |A|
  =2\begin{vmatrix}1&0\\0&1\end{vmatrix}
  -1\begin{vmatrix}-1&0\\-1&1\end{vmatrix}
  +1\begin{vmatrix}-1&1\\-1&0\end{vmatrix}
  =2+1+1=4.
  \]

  \item[(2)] $I_3$ は対角成分が1、その他の成分が0である単位行列なので、
  $A-\lambda I_3$ では $A$ の対角成分だけから $\lambda$ を引く。従って
  \[
  A-\lambda I_3
  =\begin{pmatrix}
  2-\lambda&1&1\\
  1&2-\lambda&1\\
  1&1&2-\lambda
  \end{pmatrix}.
  \]
  この行列式で $R_2\leftarrow R_2-R_1$、
  $R_3\leftarrow R_3-R_1$ と変形すると
  \begin{align*}
  |A-\lambda I_3|
  &=\begin{vmatrix}
  2-\lambda&1&1\\
  \lambda-1&1-\lambda&0\\
  \lambda-1&0&1-\lambda
  \end{vmatrix}\\
  &=(1-\lambda)^2
  \begin{vmatrix}
  2-\lambda&1&1\\
  -1&1&0\\
  -1&0&1
  \end{vmatrix}\\
  &=(1-\lambda)^2\{(2-\lambda)+1+1\}\\
  &=(1-\lambda)^2(4-\lambda).
  \end{align*}
  2行目と3行目では
  $\lambda-1=-(1-\lambda)$ であるため、
  各行から $1-\lambda$ を一つずつくくり出した。
  よって固有方程式は
  \[
  (1-\lambda)^2(4-\lambda)=0
  \]
  であり、固有値は $\lambda=1,4$ である。

  $\lambda=1$ のとき
  \[
  (A-I_3)\begin{pmatrix}x\\y\\z\end{pmatrix}
  =\begin{pmatrix}
  1&1&1\\1&1&1\\1&1&1
  \end{pmatrix}
  \begin{pmatrix}x\\y\\z\end{pmatrix}
  =\begin{pmatrix}x+y+z\\x+y+z\\x+y+z\end{pmatrix}.
  \]
  従って $(A-I_3)\boldsymbol{x}=\boldsymbol{0}$ は
  $x+y+z=0$ と同値であり、
  \[
  W_1=\ker(A-I_3)
  =\left\{
  \begin{pmatrix}x\\y\\z\end{pmatrix}\in\R^3
  \mathrel{}\middle|\mathrel{} x+y+z=0
  \right\}.
  \]
  $\lambda=4$ のとき、連立方程式は
  \[
  \begin{cases}
  -2x+y+z=0,\\
  x-2y+z=0,\\
  x+y-2z=0.
  \end{cases}
  \]
  第1式から第2式を引くと $-3x+3y=0$ なので $x=y$、
  第2式から第3式を引くと $-3y+3z=0$ なので $y=z$ である。従って
  \[
  W_4=\ker(A-4I_3)
  =\left\langle
  \begin{pmatrix}1\\1\\1\end{pmatrix}
  \right\rangle.
  \]

  \item[(3)] 長さが1で、異なる二本の内積が0であるベクトルの組を
  「正規直交している」という。$W_1$ からそのような二本を選び、
  $W_4$ から長さ1の一本を選ぶと
  \[
  \boldsymbol{u}_1=\frac1{\sqrt2}
  \begin{pmatrix}1\\-1\\0\end{pmatrix},
  \qquad
  \boldsymbol{u}_2=\frac1{\sqrt6}
  \begin{pmatrix}1\\1\\-2\end{pmatrix},
  \qquad
  \boldsymbol{u}_3=\frac1{\sqrt3}
  \begin{pmatrix}1\\1\\1\end{pmatrix}
  \]
  を取る。互いの内積は
  \[
  \boldsymbol{u}_1^{\mathsf T}\boldsymbol{u}_2
  =\frac{1-1+0}{\sqrt{12}}=0,\qquad
  \boldsymbol{u}_1^{\mathsf T}\boldsymbol{u}_3
  =\frac{1-1+0}{\sqrt6}=0,
  \]
  \[
  \boldsymbol{u}_2^{\mathsf T}\boldsymbol{u}_3
  =\frac{1+1-2}{\sqrt{18}}=0
  \]
  であり、長さについても
  \[
  \|\boldsymbol{u}_1\|^2=\frac{1+1}{2}=1,\quad
  \|\boldsymbol{u}_2\|^2=\frac{1+1+4}{6}=1,\quad
  \|\boldsymbol{u}_3\|^2=\frac{1+1+1}{3}=1
  \]
  と確認できる。
  従って
  \[
  P=(\boldsymbol{u}_1\ \boldsymbol{u}_2\ \boldsymbol{u}_3)
  \]
  は直交行列である。また
  \[
  A\boldsymbol{u}_1=\boldsymbol{u}_1,\quad
  A\boldsymbol{u}_2=\boldsymbol{u}_2,\quad
  A\boldsymbol{u}_3=4\boldsymbol{u}_3
  \]
  だから
  \[
  P^{\mathsf T}AP=\diag(1,1,4).
  \]

  \item[(4)] 単位ベクトル
  $\boldsymbol{x}=(x,y,z)^{\mathsf T}$ を
  \[
  \boldsymbol{x}=c_1\boldsymbol{u}_1+
  c_2\boldsymbol{u}_2+c_3\boldsymbol{u}_3
  \]
  と表す。$\boldsymbol{u}_1,\boldsymbol{u}_2,\boldsymbol{u}_3$ は
  長さ1で互いに直交する基底なので
  $c_1^2+c_2^2+c_3^2=1$ であり、(3) より
  \begin{align*}
  Q(x,y,z)
  &=c_1^2+c_2^2+4c_3^2\\
  &=1+3c_3^2.
  \end{align*}
  従って最小値は $1$、最大値は $4$ である。
  ここで
  \[
  c_3=\boldsymbol{u}_3^{\mathsf T}\boldsymbol{x}
  =\frac{x+y+z}{\sqrt3}
  \]
  だから、最小値は $c_3=0$、すなわち $x+y+z=0$ のときに成立する。
  最大値は $c_3^2=1$ のときに成立する。このとき
  $c_1^2+c_2^2=0$ なので $c_1=c_2=0$ であり、
  \[
  (x,y,z)=\pm\frac1{\sqrt3}(1,1,1)
  \]
  のときに成立する。
\end{enumerate}
\end{proof}

\begin{problem}{予想第2問｜重積分と変数変換}{prediction-b-2}
\[
D=\{(x,y)\in\R^2\mid x\ge0,\ y\ge0,\ 1\le x+y\le2\}
\]
とし
\[
u=x+y,\qquad v=x-y
\]
とおく。
\begin{enumerate}
  \item[(1)] $x,y$ を $u,v$ で表し、
  $\left|\dfrac{\partial(x,y)}{\partial(u,v)}\right|$ を求めよ。
  \item[(2)] $D$ に対応する $uv$ 平面上の領域を求めよ。
  \item[(3)] 次の重積分を求めよ。
  \[
  \iint_D\frac{(x-y)^2}{x+y}\,\d x\d y
  \]
\end{enumerate}
\end{problem}
\begin{memo*}
\textbf{出題根拠。}
積分領域を読み取り、変数変換の Jacobian と積分区間を決める問題は繰り返し出題されている。

\textbf{解答の見通し。}
$u=x+y$, $v=x-y$ の二式を足し引きすると
\[
x=\frac{u+v}{2},\qquad y=\frac{u-v}{2}
\]
となる。この逆変換を使って
$x\ge0$, $y\ge0$, $1\le x+y\le2$ を一つずつ書き換えると
\[
1\le u\le2,\qquad -u\le v\le u
\]
を得る。また面積要素は
\[
\frac{\partial(x,y)}{\partial(u,v)}
=\begin{vmatrix}1/2&1/2\\1/2&-1/2\end{vmatrix}
=-\frac12,
\qquad
\left|\frac{\partial(x,y)}{\partial(u,v)}\right|=\frac12
\]
だから、$\d x\d y=(1/2)\,\d u\d v$ と変換する。
最後に、被積分関数の $x+y$ を $u$、$x-y$ を $v$ に置き換え、
$u$ を外側、$v$ を内側の変数として積分する。
\end{memo*}
\begin{proof}[解答]
\begin{enumerate}
  \item[(1)] 二式を辺々加えると $u+v=2x$、
  $u$ から $v$ を引くと $u-v=2y$ なので
  \[
  x=\frac{u+v}{2},\qquad y=\frac{u-v}{2}.
  \]
  従って
  \[
  \frac{\partial(x,y)}{\partial(u,v)}
  =
  \begin{vmatrix}
  1/2&1/2\\
  1/2&-1/2
  \end{vmatrix}
  =-\frac12
  \]
  であり、その絶対値は
  \[
  \left|\frac{\partial(x,y)}{\partial(u,v)}\right|=\frac12
  \]
  である。

  \item[(2)] $1\le x+y\le2$ は
  \[
  1\le u\le2
  \]
  となる。また
  \[
  x=\frac{u+v}{2}\ge0\quad\Longleftrightarrow\quad v\ge-u,
  \]
  \[
  y=\frac{u-v}{2}\ge0\quad\Longleftrightarrow\quad v\le u.
  \]
  従って対応する領域は
  \[
  \Omega=\{(u,v)\in\R^2\mid1\le u\le2,\ -u\le v\le u\}
  \]
  である。

  \item[(3)] 被積分関数は
  \[
  \frac{(x-y)^2}{x+y}=\frac{v^2}{u}
  \]
  となる。(1), (2) より
  \begin{align*}
  \iint_D\frac{(x-y)^2}{x+y}\,\d x\d y
  &=\int_1^2\int_{-u}^{u}
  \frac{v^2}{u}\cdot\frac12\,\d v\d u\\
  &=\int_1^2\frac1{2u}
  \left[\frac{v^3}{3}\right]_{-u}^{u}\d u\\
  &=\int_1^2\frac{u^2}{3}\,\d u\\
  &=\frac13\left[\frac{u^3}{3}\right]_1^2
  =\frac79.
  \end{align*}
\end{enumerate}
\end{proof}

\begin{problem}{予想第3問｜コンパクト集合と距離}{prediction-b-3}
$(X,d)$ を距離空間、$K\subset X$ を空でないコンパクト集合とし
\[
d(x,K)=\inf\{d(x,k)\mid k\in K\}
\]
とおく。
\begin{enumerate}
  \item[(1)] 任意の $x,y\in X$ に対し
  \[
  |d(x,K)-d(y,K)|\le d(x,y)
  \]
  を示せ。
  \item[(2)] 任意の $x\in X$ に対し、
  $d(x,K)=d(x,k_x)$ となる $k_x\in K$ が存在することを示せ。
  \item[(3)] $F\subset X$ が空でない閉集合で $K\cap F=\varnothing$ を満たすとき
  \[
  d(K,F):=\inf\{d(k,f)\mid k\in K,\ f\in F\}>0
  \]
  を示せ。
  \item[(4)] $U$ が $K$ を含む開集合なら、ある $\varepsilon>0$ が存在して
  \[
  \{x\in X\mid d(x,K)<\varepsilon\}\subset U
  \]
  となることを示せ。
\end{enumerate}
\end{problem}
\begin{memo*}
\textbf{出題根拠。}
連続性、コンパクト性、閉集合、距離関数を定義からつなぐ証明は専門基礎で頻出している。

\textbf{解答の見通し。}
集合 $S$ までの距離は
\[
d(x,S)=\inf_{s\in S}d(x,s)
\]
である。ここで下限とは、全ての $d(x,s)$ 以下にある数のうち最大のものである。
(1) では三角不等式
$d(x,s)\le d(x,y)+d(y,s)$ の両辺で $s\in S$ に関する下限を取る。
(2) では $d(x,S)$ に近づく $s_n\in S$ を選び、コンパクト性から収束部分列を取る。
(3) では連続関数 $x\mapsto d(x,F)$ を $K$ 上で最小化して
$d(K,F)=\min_{x\in K}d(x,F)>0$ を示す。
(4) は $F=X\setminus U$ と置き、(3) の正の距離より小さい $\varepsilon$ を選ぶ。
\end{memo*}
\begin{proof}[解答]
\begin{enumerate}
  \item[(1)] 任意の $k\in K$ に対して三角不等式より
  \[
  d(x,k)\le d(x,y)+d(y,k).
  \]
  $d(x,K)\le d(x,k)$ でもあるから、全ての $k\in K$ について
  \[
  d(x,K)\le d(x,y)+d(y,k)
  \]
  が成り立つ。右辺について $k\in K$ の下限を取ると
  \[
  d(x,K)\le d(x,y)+d(y,K).
  \]
  従って
  \[
  d(x,K)-d(y,K)\le d(x,y).
  \]
  $x,y$ を交換して
  \[
  d(y,K)-d(x,K)\le d(x,y)
  \]
  も得られるので、二式を合わせれば結論を得る。

  \item[(2)] 下限の定義により、各 $n\in\N$ に対して $k_n\in K$ を
  \[
  d(x,K)\le d(x,k_n)<d(x,K)+\frac1n
  \]
  となるように取れる。$K$ はコンパクトだから、ある部分列
  $(k_{n_j})$ と $k_x\in K$ が存在して $k_{n_j}\to k_x$ となる。
  (1) と同じ三角不等式から写像 $k\mapsto d(x,k)$ は連続なので
  \[
  d(x,k_x)
  =\lim_{j\to\infty}d(x,k_{n_j})
  =d(x,K).
  \]

  \item[(3)] (1) の $K$ を $F$ に置き換えると
  $h(k)=d(k,F)$ は連続である。連続関数 $h$ はコンパクト集合 $K$ 上で
  最小値を取るから、ある $k_0\in K$ が存在して
  \[
  \inf_{k\in K}d(k,F)=\min_{k\in K}d(k,F)=d(k_0,F)
  \]
  となる。
  一方、二重の下限を順に取れば
  \[
  d(K,F)
  =\inf_{k\in K}\inf_{f\in F}d(k,f)
  =\inf_{k\in K}d(k,F)
  \]
  である。従って $d(K,F)=d(k_0,F)$ である。

  もし $d(k_0,F)=0$ なら、各 $n$ に対して $f_n\in F$ を
  $d(k_0,f_n)<1/n$ となるように取れる。すると $f_n\to k_0$ である。
  $F$ は閉集合だから $k_0\in F$ となり、$k_0\in K$ と合わせて
  $K\cap F\ne\varnothing$ となる。これは仮定に反する。
  従って $d(K,F)>0$ である。

  \item[(4)] $F=X\setminus U$ とおく。$F=\varnothing$ なら任意の
  $\varepsilon>0$ で結論が成り立つ。$F\ne\varnothing$ とする。
  $U$ は開集合だから $F$ は閉集合であり、$K\subset U$ より
  $K\cap F=\varnothing$ である。(3) により
  \[
  \delta=d(K,F)>0.
  \]
  $\varepsilon=\delta/2$ と取る。もし $d(x,K)<\varepsilon$ かつ
  $x\notin U$ なら $x\in F$ なので
  \[
  \delta=d(K,F)\le d(x,K)<\frac{\delta}{2},
  \]
  となり矛盾する。従って $d(x,K)<\varepsilon$ なら $x\in U$ である。
\end{enumerate}
\end{proof}


\end{document}
